% Documentul master: pagina de titlu + harta a celor 8 PDF-uri.
% Compilare: pdflatex main.tex (de doua ori, pentru cuprins).

\documentclass[11pt,a4paper]{article}
% Preambul comun pentru toate handout-urile training-ului.
% Fiecare fisier .tex face % Preambul comun pentru toate handout-urile training-ului.
% Fiecare fisier .tex face % Preambul comun pentru toate handout-urile training-ului.
% Fiecare fisier .tex face \input{preamble} dupa \documentclass.

\usepackage[utf8]{inputenc}
\usepackage[T1]{fontenc}
% Body font: Lato -- modern, lizibil, profesional. Fallback la TeX Gyre Heros
% (Helvetica-like, mereu disponibil) daca Lato nu e instalat in MiKTeX.
\IfFileExists{lato.sty}{%
  \usepackage[default]{lato}%
}{%
  \usepackage{tgheros}%
  \renewcommand{\familydefault}{\sfdefault}%
}
\usepackage[romanian]{babel}
\usepackage[a4paper, margin=2.5cm, headheight=16pt, top=2.4cm, bottom=2.4cm]{geometry}
\usepackage[final, protrusion=true, expansion=false]{microtype}
\usepackage{setspace}
\usepackage{xcolor}
\usepackage{amssymb}
\usepackage{listings}
\usepackage{hyperref}
\usepackage{enumitem}
\usepackage{titlesec}
\usepackage{fancyhdr}
\usepackage{parskip}
\usepackage[most]{tcolorbox}
\usepackage{tocloft}
\usepackage{graphicx}
\usepackage{tikz}
\usepackage{fontawesome5}
\usepackage{lettrine}
\usetikzlibrary{shadows.blur, shapes.geometric, positioning, calc,
  decorations.pathmorphing, patterns, fadings, backgrounds}

% =========================================================
%   PALETA DE CULORI
% =========================================================

% Paleta finala (3 BEST + 3 derivate in acelasi ton)
\definecolor{themePrereq}{HTML}{5C6B7A}    % slate (pre-rechizite, neutral)
\definecolor{themeHtml}{HTML}{FAA51E}      % BEST orange
\definecolor{themeCss}{HTML}{0673BA}       % BEST blue
\definecolor{themeJs}{HTML}{ECC100}        % gold (analog cu orange, distinct)
\definecolor{themeDeploy}{HTML}{75C044}    % BEST green
\definecolor{themeSeo}{HTML}{B344C9}       % violet (analog cu blue)
\definecolor{themeLighthouse}{HTML}{DD4444}% coral red (complementar cu green)
\definecolor{themeResurse}{HTML}{5C6B7A}   % slate (resurse, neutral)
\definecolor{themeBonus}{HTML}{7E22CE}     % deep purple (bonus / avansat, self-study only)
% Aliasuri pentru compatibilitate (folosit in main.tex existent)
\colorlet{themeIntro}{themeHtml}
\colorlet{themeExercitiu}{themeLighthouse}

% Background pagina si elemente -- calduros
\definecolor{pageBg}{HTML}{FFFCF6}         % crema foarte subtila
\definecolor{accentSoft}{HTML}{F4ECDD}     % accent cald
\definecolor{textDim}{HTML}{64748B}        % gri pentru text secundar

% Culoarea curenta a partii (suprascrisa de fiecare fisier prin \setpart)
\colorlet{partcolor}{themeIntro}

% Iconita partii curente (suprascrisa prin \setpart)
\newcommand{\particon}{\faStar}
\newcommand{\partlabel}{Training Web}

% \setpart{culoare}{iconita}
\newcommand{\setpart}[2]{%
  \colorlet{partcolor}{#1}%
  \renewcommand{\particon}{#2}%
}
% \setpart{culoare}{iconita}{eticheta} -- versiunea cu eticheta custom in header
\newcommand{\setpartlabel}[3]{%
  \colorlet{partcolor}{#1}%
  \renewcommand{\particon}{#2}%
  \renewcommand{\partlabel}{#3}%
}

% Culori pentru callout boxes
\definecolor{noteBg}{HTML}{FEF9C3}        \definecolor{noteFg}{HTML}{A16207}
\definecolor{tipBg}{HTML}{DBEAFE}         \definecolor{tipFg}{HTML}{1D4ED8}
\definecolor{warnBg}{HTML}{FEE2E2}        \definecolor{warnFg}{HTML}{B91C1C}
\definecolor{gameBg}{HTML}{FFEDD5}        \definecolor{gameFg}{HTML}{C2410C}
\definecolor{recapBg}{HTML}{DCFCE7}       \definecolor{recapFg}{HTML}{15803D}
\definecolor{exerBg}{HTML}{F3E8FF}        \definecolor{exerFg}{HTML}{7C3AED}
\definecolor{triviaBg}{HTML}{FCE7F3}      \definecolor{triviaFg}{HTML}{BE185D}

% Culori pentru cod -- paleta "GitHub Light" (calm, profesional)
\definecolor{codeBg}{HTML}{F6F8FA}
\definecolor{codeBorder}{HTML}{D0D7DE}
\definecolor{codeComment}{HTML}{6E7781}
\definecolor{codeKeyword}{HTML}{CF222E}
\definecolor{codeString}{HTML}{0A3069}
\definecolor{codeTag}{HTML}{116329}

% =========================================================
%   PAGINA & LIZIBILITATE (warmth)
% =========================================================
\pagecolor{pageBg}
\setstretch{1.18}
\setlength{\parskip}{0.5em}

% =========================================================
%   CUPRINS (TOC)
% =========================================================
\setcounter{tocdepth}{0}         % numai entry-uri \part (capitole), fara sectiuni
\renewcommand{\contentsname}{Cuprins}

% Stil cuprins -- numai capitolele, asezate elegant
\renewcommand{\cftpartfont}{\sffamily\bfseries\color{themeText}}
\renewcommand{\cftpartpagefont}{\sffamily\bfseries\color{themeText}}
\renewcommand{\cftpartleader}{\hfill}
\renewcommand{\cftpartafterpnum}{\par\addvspace{2pt}}
\setlength{\cftbeforepartskip}{4pt}

% =========================================================
%   HYPERREF
% =========================================================
\hypersetup{
  colorlinks=true,
  linkcolor=partcolor,
  urlcolor=partcolor,
  citecolor=partcolor,
  pdfborderstyle={/S/U/W 1}
}

% =========================================================
%   HEADER & FOOTER -- minimaliste
% =========================================================
\pagestyle{fancy}
\fancyhf{}
\fancyhead[L]{\footnotesize\sffamily\color{black!55}\partlabel}
\fancyhead[R]{\footnotesize\sffamily\color{black!55}BEST Ia\textcommabelow{s}i}
\fancyfoot[C]{\small\sffamily\color{partcolor}\textbf{\thepage}}
\renewcommand{\headrulewidth}{0pt}
\renewcommand{\footrulewidth}{0pt}

% =========================================================
%   STIL TITLURI -- sobru, stiintific
% =========================================================
\titleformat{\section}
  {\Large\sffamily\bfseries\color{themeText}}
  {\color{partcolor}\sffamily\bfseries\thesection.\hspace{0.5em}}
  {0pt}{}
\titlespacing*{\section}{0pt}{1.4em}{0.4em}

\titleformat{\subsection}
  {\normalsize\sffamily\bfseries\color{black!75}}
  {\color{partcolor}\thesubsection\hspace{0.5em}}{0pt}{}
\titlespacing*{\subsection}{0pt}{0.8em}{0.2em}

\titleformat{\subsubsection}
  {\normalsize\sffamily\bfseries\color{black!70}}
  {}{0pt}{}

\definecolor{themeText}{HTML}{1A1A1A}

% =========================================================
%   LISTINGS -- COD
% =========================================================
\lstdefinestyle{base}{
  backgroundcolor=\color{codeBg},
  basicstyle=\ttfamily\small,
  commentstyle=\color{codeComment}\itshape,
  keywordstyle=\color{codeKeyword}\bfseries,
  stringstyle=\color{codeString},
  numberstyle=\tiny\color{codeComment},
  numbers=left,
  numbersep=8pt,
  frame=leftline,
  framerule=2.5pt,
  rulecolor=\color{partcolor},
  framexleftmargin=4pt,
  breaklines=true,
  breakatwhitespace=true,
  showstringspaces=false,
  tabsize=2,
  xleftmargin=14pt,
  xrightmargin=4pt,
  aboveskip=10pt,
  belowskip=10pt,
  literate=
    {ă}{{\u a}}1 {Ă}{{\u A}}1
    {â}{{\^a}}1 {Â}{{\^A}}1
    {î}{{\^\i}}1 {Î}{{\^I}}1
    {ș}{{\textcommabelow s}}1 {Ș}{{\textcommabelow S}}1
    {ț}{{\textcommabelow t}}1 {Ț}{{\textcommabelow T}}1
}

\lstdefinelanguage{HTML5}{
  morekeywords={html,head,body,title,meta,link,script,style,
    h1,h2,h3,h4,h5,h6,p,a,img,div,span,
    header,main,section,article,aside,nav,footer,
    ul,ol,li,button,input,form,label,
    strong,em,br,hr,DOCTYPE},
  sensitive=true,
  morecomment=[s]{<!--}{-->},
  morestring=[b]",
  morestring=[b]'
}
\lstdefinestyle{html}{style=base, language=HTML5,
  keywordstyle=\color{codeTag}\bfseries}

\lstdefinelanguage{CSSLang}{
  morekeywords={display,flex,grid,block,inline,none,
    color,background,background-color,
    padding,padding-top,padding-right,padding-bottom,padding-left,
    margin,margin-top,margin-right,margin-bottom,margin-left,
    border,border-radius,border-color,outline,
    width,max-width,min-width,height,max-height,min-height,
    font-family,font-size,font-weight,line-height,text-align,
    text-decoration,letter-spacing,
    flex-direction,justify-content,align-items,gap,flex-wrap,flex-grow,
    grid-template-columns,grid-template-rows,grid-column,grid-row,
    position,top,right,bottom,left,z-index,
    box-sizing,opacity,transition,transform,
    cursor,overflow,visibility,
    media,keyframes,from,to,important,root,hover,focus,active,
    auto,solid,dashed,dotted,bold,normal,center,stretch,
    space-between,space-around,flex-start,flex-end,column,row,wrap},
  sensitive=true,
  morecomment=[s]{/*}{*/},
  morestring=[b]",
  morestring=[b]'
}
\lstdefinestyle{css}{style=base, language=CSSLang,
  keywordstyle=\color{codeKeyword}\bfseries}

\lstdefinelanguage{JSLang}{
  morekeywords={let,const,var,function,return,if,else,for,while,
    true,false,null,undefined,new,this,class,extends,
    document,console,window,
    addEventListener,querySelector,querySelectorAll,
    getElementById,getElementsByClassName,
    classList,toggle,add,remove,contains,
    style,textContent,innerHTML,value,
    log,warn,error},
  sensitive=true,
  morecomment=[l]{//},
  morecomment=[s]{/*}{*/},
  morestring=[b]",
  morestring=[b]',
  morestring=[b]`
}
\lstdefinestyle{js}{style=base, language=JSLang,
  keywordstyle=\color{codeKeyword}\bfseries}

\lstdefinelanguage{BashLang}{
  morekeywords={git,cd,mkdir,ls,code,echo,cat,curl,wget,sudo,
    init,add,commit,push,pull,clone,remote,branch,checkout,merge,status,log,
    config,fetch,rebase,reset,stash,diff,tag,
    install,update,run,build,start,test},
  sensitive=true,
  morecomment=[l]{\#},
  morestring=[b]",
  morestring=[b]'
}
\lstdefinestyle{bash}{style=base, language=BashLang,
  keywordstyle=\color{codeKeyword}\bfseries}

% Inline code -- mai discret
\newcommand{\code}[1]{\textcolor{partcolor!85!black}{\ttfamily\small #1}}

% =========================================================
%   CALLOUT BOXES -- minimaliste (bara stanga subtila + label discret)
% =========================================================
\tcbset{
  calloutbase/.style={
    enhanced,
    breakable,
    sharp corners,
    boxrule=0pt,
    leftrule=2.5pt,
    arc=0pt,
    colback=pageBg,
    fonttitle=\sffamily\small\bfseries,
    coltitle=tcbcolframe,
    attach title to upper={\par\vspace{2pt}},
    title after break={},
    before skip=10pt,
    after skip=10pt,
    left=12pt, right=4pt, top=4pt, bottom=4pt,
    fontupper=\small,
  }
}

\newtcolorbox{nota}[1][NOT\u A]{
  calloutbase,
  colframe=noteFg,
  title={\faInfoCircle\ \uppercase\expandafter{#1}},
}
\newtcolorbox{ponturi}[1][PONT]{
  calloutbase,
  colframe=tipFg,
  title={\faLightbulb[regular]\ \uppercase\expandafter{#1}},
}
\newtcolorbox{atentie}[1][ATEN\textcommabelow{T}IE]{
  calloutbase,
  colframe=warnFg,
  title={\faExclamationTriangle\ \uppercase\expandafter{#1}},
}
\newtcolorbox{joc}[1][APLICA\textcommabelow{T}IE PRACTIC\u A]{
  calloutbase,
  colframe=partcolor,
  title={\faGamepad\ \uppercase\expandafter{#1}},
}
\newtcolorbox{recap}[1][RECAP]{
  calloutbase,
  colframe=recapFg,
  title={\faCheckCircle\ \uppercase\expandafter{#1}},
}
\newtcolorbox{curiozitate}[1][\textcommabelow{S}TIAI C\u A?]{
  calloutbase,
  colframe=triviaFg,
  title={\faMagic\ \uppercase\expandafter{#1}},
}
\newtcolorbox{exercitiu}[1][EXERCI\textcommabelow{T}IU]{
  calloutbase,
  colframe=exerFg,
  title={\faPen\ \uppercase\expandafter{#1}},
}

% --- Outcome / Why / Self-check (structura tutorial) ----
\definecolor{outBg}{HTML}{E0F2FE}        \definecolor{outFg}{HTML}{0369A1}
\definecolor{whyBg}{HTML}{FEF3C7}        \definecolor{whyFg}{HTML}{B45309}
\definecolor{checkBg}{HTML}{F3E8FF}      \definecolor{checkFg}{HTML}{7C3AED}

\newtcolorbox{outcomes}[1][CE INVE\textcommabelow{T}I]{
  calloutbase,
  colframe=outFg,
  title={\faBullseye\ \uppercase\expandafter{#1}},
}
\newtcolorbox{whymatters}[1][DE CE]{
  calloutbase,
  colframe=whyFg,
  title={\faQuestionCircle\ \uppercase\expandafter{#1}},
}
\newtcolorbox{selfcheck}[1][PROVOCARE]{
  calloutbase,
  colframe=checkFg,
  title={\faTasks\ \uppercase\expandafter{#1}},
}

% --- Bonus / Avansat (continut auto-studiu, distinct vizual) ----
% Diferit fata de celelalte callout-uri: fundal usor colorat (nu pageBg),
% bara stanga mai groasa, icon racheta. Beginners vad clar si pot sari peste.
\definecolor{bonusBg}{HTML}{F6EEFB}        \definecolor{bonusFg}{HTML}{6B21A8}

\newtcolorbox{bonusbox}[1][BONUS \,\textbullet\, AVANSAT]{
  enhanced,
  breakable,
  sharp corners,
  boxrule=0pt,
  leftrule=4pt,
  arc=0pt,
  colback=bonusBg,
  colframe=themeBonus,
  fonttitle=\sffamily\small\bfseries,
  coltitle=bonusFg,
  attach title to upper={\par\vspace{2pt}},
  title after break={\faRocket\ \textbf{\color{bonusFg}continuare bonus}},
  before skip=14pt,
  after skip=14pt,
  left=14pt, right=6pt, top=6pt, bottom=6pt,
  fontupper=\small,
  title={\faRocket\ \uppercase\expandafter{#1}},
}

% Bridge la finalul fiecarui capitol
\newcommand{\bridge}[1]{%
  \par\vspace{4pt}%
  \begin{flushleft}%
    {\sffamily\small\color{partcolor}$\rightarrow$\ \textbf{Urmeaza:}}\quad{\small\color{black!70}#1}%
  \end{flushleft}%
}

% --- Asa da / Asa nu (compare bune vs rele practici) ----
\definecolor{daBg}{HTML}{DCFCE7}        \definecolor{daFg}{HTML}{15803D}
\definecolor{nuBg}{HTML}{FEE2E2}        \definecolor{nuFg}{HTML}{B91C1C}

\newtcolorbox{dado}[1][A\textcommabelow{S}A DA]{
  calloutbase,
  colframe=daFg,
  title={\faCheckCircle\ \uppercase\expandafter{#1}},
}
\newtcolorbox{dano}[1][A\textcommabelow{S}A NU]{
  calloutbase,
  colframe=nuFg,
  title={\faTimesCircle\ \uppercase\expandafter{#1}},
}

% =========================================================
%   BANNER MINIMALIST (\handout{titlu}{subtitlu})
%   Eticheta sus, titlu mare, subtitlu, linie subtila colorata.
% =========================================================
\newcommand{\handout}[2]{%
  \noindent\begin{minipage}{0.04\linewidth}\color{partcolor}\rule{4pt}{2.6em}\end{minipage}%
  \begin{minipage}{0.94\linewidth}
    {\sffamily\fontsize{30}{34}\selectfont\textbf{\color{themeText}#1}}\\[4pt]
    {\sffamily\large\color{black!55}#2}
  \end{minipage}
  \par\vspace{14pt}
  {\color{partcolor!40}\hrule height 0.5pt}
  \par\vspace{14pt}
}

% =========================================================
%   BULLETS COLORATI (itemize)
% =========================================================
\setlist[itemize,1]{label={\color{partcolor}\textbf{\textbullet}}}
\setlist[itemize,2]{label={\color{partcolor!70}\textbf{\textendash}}}
\setlist[enumerate,1]{label={\color{partcolor}\textbf{\arabic*.}}}

% =========================================================
%   PULL QUOTE / ACCENT
% =========================================================
\newcommand{\pullquote}[1]{%
  \begin{tcolorbox}[
    enhanced,
    colback=partcolor!8,
    colframe=partcolor,
    boxrule=0pt,
    leftrule=4pt,
    sharp corners,
    arc=0pt,
    fontupper=\itshape\large,
    left=14pt, right=14pt, top=8pt, bottom=8pt,
    overlay={%
      \node[partcolor!40, font=\fontsize{40}{40}\selectfont, anchor=north west]
        at ([xshift=2pt,yshift=2pt]frame.north west) {\textbf{``}};
    },
  ]
    #1
  \end{tcolorbox}
}

% =========================================================
%   DIVIDER DECORATIV (foloseste \partdivider in document)
%   Atentie: NU folosi numele "\sectionbreak" -- e rezervat
%   intern de titlesec si va aparea automat intre sectiuni.
% =========================================================
\newcommand{\partdivider}{%
  \par\vspace{0.6em}
  \begin{center}
    \tikz[baseline=-0.5ex]{
      \fill[partcolor] (0,0) circle (2pt);
      \fill[partcolor!60] (10pt,0) circle (2pt);
      \fill[partcolor!30] (20pt,0) circle (2pt);
    }
  \end{center}
  \vspace{0.3em}
}

% =========================================================
%   BOX MODEL DIAGRAM (TikZ)
% =========================================================
\newcommand{\boxmodelfigure}{%
\begin{center}
\begin{tikzpicture}[scale=1, transform shape, every node/.style={font=\sffamily\small}]
  % Margin (outermost)
  \fill[orange!25, rounded corners=4pt] (0,0) rectangle (8,5);
  \node[anchor=north west] at (0.15, 4.85) {\textbf{\color{orange!70!black}\faExpand\ margin}};

  % Border
  \fill[gray!40, rounded corners=3pt] (1,0.7) rectangle (7, 4.3);
  \draw[ultra thick, gray!70!black, rounded corners=3pt] (1,0.7) rectangle (7,4.3);
  \node[anchor=north west] at (1.1, 4.2) {\textbf{\color{gray!50!black}\faSquare\ border}};

  % Padding
  \fill[green!25, rounded corners=2pt] (1.3, 1) rectangle (6.7, 4);
  \node[anchor=north west] at (1.4, 3.9) {\textbf{\color{green!50!black}\faCompress\ padding}};

  % Content (innermost)
  \fill[blue!25, rounded corners=2pt] (2.2, 1.7) rectangle (5.8, 3.3);
  \node[anchor=north west] at (2.3, 3.2) {\textbf{\color{blue!60!black}\faFile*\ content}};
  \node at (4, 2.5) {\textit{textul, imaginea, ce e inauntru}};
\end{tikzpicture}
\end{center}
}

% =========================================================
%   CARD "ICON + TEXT" (utilitate pentru handout-uri)
% =========================================================
\newcommand{\iconcard}[3]{%
  \begin{tcolorbox}[
    enhanced,
    colback=partcolor!8,
    colframe=partcolor!30,
    boxrule=0.5pt,
    arc=6pt,
    left=12pt, right=12pt, top=8pt, bottom=8pt,
  ]
  \begin{minipage}[t]{0.08\linewidth}
    \color{partcolor}\fontsize{20}{20}\selectfont #1
  \end{minipage}%
  \begin{minipage}[t]{0.9\linewidth}
    \textbf{\color{partcolor!85!black}#2}\\
    #3
  \end{minipage}
  \end{tcolorbox}%
}
 dupa \documentclass.

\usepackage[utf8]{inputenc}
\usepackage[T1]{fontenc}
% Body font: Lato -- modern, lizibil, profesional. Fallback la TeX Gyre Heros
% (Helvetica-like, mereu disponibil) daca Lato nu e instalat in MiKTeX.
\IfFileExists{lato.sty}{%
  \usepackage[default]{lato}%
}{%
  \usepackage{tgheros}%
  \renewcommand{\familydefault}{\sfdefault}%
}
\usepackage[romanian]{babel}
\usepackage[a4paper, margin=2.5cm, headheight=16pt, top=2.4cm, bottom=2.4cm]{geometry}
\usepackage[final, protrusion=true, expansion=false]{microtype}
\usepackage{setspace}
\usepackage{xcolor}
\usepackage{amssymb}
\usepackage{listings}
\usepackage{hyperref}
\usepackage{enumitem}
\usepackage{titlesec}
\usepackage{fancyhdr}
\usepackage{parskip}
\usepackage[most]{tcolorbox}
\usepackage{tocloft}
\usepackage{graphicx}
\usepackage{tikz}
\usepackage{fontawesome5}
\usepackage{lettrine}
\usetikzlibrary{shadows.blur, shapes.geometric, positioning, calc,
  decorations.pathmorphing, patterns, fadings, backgrounds}

% =========================================================
%   PALETA DE CULORI
% =========================================================

% Paleta finala (3 BEST + 3 derivate in acelasi ton)
\definecolor{themePrereq}{HTML}{5C6B7A}    % slate (pre-rechizite, neutral)
\definecolor{themeHtml}{HTML}{FAA51E}      % BEST orange
\definecolor{themeCss}{HTML}{0673BA}       % BEST blue
\definecolor{themeJs}{HTML}{ECC100}        % gold (analog cu orange, distinct)
\definecolor{themeDeploy}{HTML}{75C044}    % BEST green
\definecolor{themeSeo}{HTML}{B344C9}       % violet (analog cu blue)
\definecolor{themeLighthouse}{HTML}{DD4444}% coral red (complementar cu green)
\definecolor{themeResurse}{HTML}{5C6B7A}   % slate (resurse, neutral)
\definecolor{themeBonus}{HTML}{7E22CE}     % deep purple (bonus / avansat, self-study only)
% Aliasuri pentru compatibilitate (folosit in main.tex existent)
\colorlet{themeIntro}{themeHtml}
\colorlet{themeExercitiu}{themeLighthouse}

% Background pagina si elemente -- calduros
\definecolor{pageBg}{HTML}{FFFCF6}         % crema foarte subtila
\definecolor{accentSoft}{HTML}{F4ECDD}     % accent cald
\definecolor{textDim}{HTML}{64748B}        % gri pentru text secundar

% Culoarea curenta a partii (suprascrisa de fiecare fisier prin \setpart)
\colorlet{partcolor}{themeIntro}

% Iconita partii curente (suprascrisa prin \setpart)
\newcommand{\particon}{\faStar}
\newcommand{\partlabel}{Training Web}

% \setpart{culoare}{iconita}
\newcommand{\setpart}[2]{%
  \colorlet{partcolor}{#1}%
  \renewcommand{\particon}{#2}%
}
% \setpart{culoare}{iconita}{eticheta} -- versiunea cu eticheta custom in header
\newcommand{\setpartlabel}[3]{%
  \colorlet{partcolor}{#1}%
  \renewcommand{\particon}{#2}%
  \renewcommand{\partlabel}{#3}%
}

% Culori pentru callout boxes
\definecolor{noteBg}{HTML}{FEF9C3}        \definecolor{noteFg}{HTML}{A16207}
\definecolor{tipBg}{HTML}{DBEAFE}         \definecolor{tipFg}{HTML}{1D4ED8}
\definecolor{warnBg}{HTML}{FEE2E2}        \definecolor{warnFg}{HTML}{B91C1C}
\definecolor{gameBg}{HTML}{FFEDD5}        \definecolor{gameFg}{HTML}{C2410C}
\definecolor{recapBg}{HTML}{DCFCE7}       \definecolor{recapFg}{HTML}{15803D}
\definecolor{exerBg}{HTML}{F3E8FF}        \definecolor{exerFg}{HTML}{7C3AED}
\definecolor{triviaBg}{HTML}{FCE7F3}      \definecolor{triviaFg}{HTML}{BE185D}

% Culori pentru cod -- paleta "GitHub Light" (calm, profesional)
\definecolor{codeBg}{HTML}{F6F8FA}
\definecolor{codeBorder}{HTML}{D0D7DE}
\definecolor{codeComment}{HTML}{6E7781}
\definecolor{codeKeyword}{HTML}{CF222E}
\definecolor{codeString}{HTML}{0A3069}
\definecolor{codeTag}{HTML}{116329}

% =========================================================
%   PAGINA & LIZIBILITATE (warmth)
% =========================================================
\pagecolor{pageBg}
\setstretch{1.18}
\setlength{\parskip}{0.5em}

% =========================================================
%   CUPRINS (TOC)
% =========================================================
\setcounter{tocdepth}{0}         % numai entry-uri \part (capitole), fara sectiuni
\renewcommand{\contentsname}{Cuprins}

% Stil cuprins -- numai capitolele, asezate elegant
\renewcommand{\cftpartfont}{\sffamily\bfseries\color{themeText}}
\renewcommand{\cftpartpagefont}{\sffamily\bfseries\color{themeText}}
\renewcommand{\cftpartleader}{\hfill}
\renewcommand{\cftpartafterpnum}{\par\addvspace{2pt}}
\setlength{\cftbeforepartskip}{4pt}

% =========================================================
%   HYPERREF
% =========================================================
\hypersetup{
  colorlinks=true,
  linkcolor=partcolor,
  urlcolor=partcolor,
  citecolor=partcolor,
  pdfborderstyle={/S/U/W 1}
}

% =========================================================
%   HEADER & FOOTER -- minimaliste
% =========================================================
\pagestyle{fancy}
\fancyhf{}
\fancyhead[L]{\footnotesize\sffamily\color{black!55}\partlabel}
\fancyhead[R]{\footnotesize\sffamily\color{black!55}BEST Ia\textcommabelow{s}i}
\fancyfoot[C]{\small\sffamily\color{partcolor}\textbf{\thepage}}
\renewcommand{\headrulewidth}{0pt}
\renewcommand{\footrulewidth}{0pt}

% =========================================================
%   STIL TITLURI -- sobru, stiintific
% =========================================================
\titleformat{\section}
  {\Large\sffamily\bfseries\color{themeText}}
  {\color{partcolor}\sffamily\bfseries\thesection.\hspace{0.5em}}
  {0pt}{}
\titlespacing*{\section}{0pt}{1.4em}{0.4em}

\titleformat{\subsection}
  {\normalsize\sffamily\bfseries\color{black!75}}
  {\color{partcolor}\thesubsection\hspace{0.5em}}{0pt}{}
\titlespacing*{\subsection}{0pt}{0.8em}{0.2em}

\titleformat{\subsubsection}
  {\normalsize\sffamily\bfseries\color{black!70}}
  {}{0pt}{}

\definecolor{themeText}{HTML}{1A1A1A}

% =========================================================
%   LISTINGS -- COD
% =========================================================
\lstdefinestyle{base}{
  backgroundcolor=\color{codeBg},
  basicstyle=\ttfamily\small,
  commentstyle=\color{codeComment}\itshape,
  keywordstyle=\color{codeKeyword}\bfseries,
  stringstyle=\color{codeString},
  numberstyle=\tiny\color{codeComment},
  numbers=left,
  numbersep=8pt,
  frame=leftline,
  framerule=2.5pt,
  rulecolor=\color{partcolor},
  framexleftmargin=4pt,
  breaklines=true,
  breakatwhitespace=true,
  showstringspaces=false,
  tabsize=2,
  xleftmargin=14pt,
  xrightmargin=4pt,
  aboveskip=10pt,
  belowskip=10pt,
  literate=
    {ă}{{\u a}}1 {Ă}{{\u A}}1
    {â}{{\^a}}1 {Â}{{\^A}}1
    {î}{{\^\i}}1 {Î}{{\^I}}1
    {ș}{{\textcommabelow s}}1 {Ș}{{\textcommabelow S}}1
    {ț}{{\textcommabelow t}}1 {Ț}{{\textcommabelow T}}1
}

\lstdefinelanguage{HTML5}{
  morekeywords={html,head,body,title,meta,link,script,style,
    h1,h2,h3,h4,h5,h6,p,a,img,div,span,
    header,main,section,article,aside,nav,footer,
    ul,ol,li,button,input,form,label,
    strong,em,br,hr,DOCTYPE},
  sensitive=true,
  morecomment=[s]{<!--}{-->},
  morestring=[b]",
  morestring=[b]'
}
\lstdefinestyle{html}{style=base, language=HTML5,
  keywordstyle=\color{codeTag}\bfseries}

\lstdefinelanguage{CSSLang}{
  morekeywords={display,flex,grid,block,inline,none,
    color,background,background-color,
    padding,padding-top,padding-right,padding-bottom,padding-left,
    margin,margin-top,margin-right,margin-bottom,margin-left,
    border,border-radius,border-color,outline,
    width,max-width,min-width,height,max-height,min-height,
    font-family,font-size,font-weight,line-height,text-align,
    text-decoration,letter-spacing,
    flex-direction,justify-content,align-items,gap,flex-wrap,flex-grow,
    grid-template-columns,grid-template-rows,grid-column,grid-row,
    position,top,right,bottom,left,z-index,
    box-sizing,opacity,transition,transform,
    cursor,overflow,visibility,
    media,keyframes,from,to,important,root,hover,focus,active,
    auto,solid,dashed,dotted,bold,normal,center,stretch,
    space-between,space-around,flex-start,flex-end,column,row,wrap},
  sensitive=true,
  morecomment=[s]{/*}{*/},
  morestring=[b]",
  morestring=[b]'
}
\lstdefinestyle{css}{style=base, language=CSSLang,
  keywordstyle=\color{codeKeyword}\bfseries}

\lstdefinelanguage{JSLang}{
  morekeywords={let,const,var,function,return,if,else,for,while,
    true,false,null,undefined,new,this,class,extends,
    document,console,window,
    addEventListener,querySelector,querySelectorAll,
    getElementById,getElementsByClassName,
    classList,toggle,add,remove,contains,
    style,textContent,innerHTML,value,
    log,warn,error},
  sensitive=true,
  morecomment=[l]{//},
  morecomment=[s]{/*}{*/},
  morestring=[b]",
  morestring=[b]',
  morestring=[b]`
}
\lstdefinestyle{js}{style=base, language=JSLang,
  keywordstyle=\color{codeKeyword}\bfseries}

\lstdefinelanguage{BashLang}{
  morekeywords={git,cd,mkdir,ls,code,echo,cat,curl,wget,sudo,
    init,add,commit,push,pull,clone,remote,branch,checkout,merge,status,log,
    config,fetch,rebase,reset,stash,diff,tag,
    install,update,run,build,start,test},
  sensitive=true,
  morecomment=[l]{\#},
  morestring=[b]",
  morestring=[b]'
}
\lstdefinestyle{bash}{style=base, language=BashLang,
  keywordstyle=\color{codeKeyword}\bfseries}

% Inline code -- mai discret
\newcommand{\code}[1]{\textcolor{partcolor!85!black}{\ttfamily\small #1}}

% =========================================================
%   CALLOUT BOXES -- minimaliste (bara stanga subtila + label discret)
% =========================================================
\tcbset{
  calloutbase/.style={
    enhanced,
    breakable,
    sharp corners,
    boxrule=0pt,
    leftrule=2.5pt,
    arc=0pt,
    colback=pageBg,
    fonttitle=\sffamily\small\bfseries,
    coltitle=tcbcolframe,
    attach title to upper={\par\vspace{2pt}},
    title after break={},
    before skip=10pt,
    after skip=10pt,
    left=12pt, right=4pt, top=4pt, bottom=4pt,
    fontupper=\small,
  }
}

\newtcolorbox{nota}[1][NOT\u A]{
  calloutbase,
  colframe=noteFg,
  title={\faInfoCircle\ \uppercase\expandafter{#1}},
}
\newtcolorbox{ponturi}[1][PONT]{
  calloutbase,
  colframe=tipFg,
  title={\faLightbulb[regular]\ \uppercase\expandafter{#1}},
}
\newtcolorbox{atentie}[1][ATEN\textcommabelow{T}IE]{
  calloutbase,
  colframe=warnFg,
  title={\faExclamationTriangle\ \uppercase\expandafter{#1}},
}
\newtcolorbox{joc}[1][APLICA\textcommabelow{T}IE PRACTIC\u A]{
  calloutbase,
  colframe=partcolor,
  title={\faGamepad\ \uppercase\expandafter{#1}},
}
\newtcolorbox{recap}[1][RECAP]{
  calloutbase,
  colframe=recapFg,
  title={\faCheckCircle\ \uppercase\expandafter{#1}},
}
\newtcolorbox{curiozitate}[1][\textcommabelow{S}TIAI C\u A?]{
  calloutbase,
  colframe=triviaFg,
  title={\faMagic\ \uppercase\expandafter{#1}},
}
\newtcolorbox{exercitiu}[1][EXERCI\textcommabelow{T}IU]{
  calloutbase,
  colframe=exerFg,
  title={\faPen\ \uppercase\expandafter{#1}},
}

% --- Outcome / Why / Self-check (structura tutorial) ----
\definecolor{outBg}{HTML}{E0F2FE}        \definecolor{outFg}{HTML}{0369A1}
\definecolor{whyBg}{HTML}{FEF3C7}        \definecolor{whyFg}{HTML}{B45309}
\definecolor{checkBg}{HTML}{F3E8FF}      \definecolor{checkFg}{HTML}{7C3AED}

\newtcolorbox{outcomes}[1][CE INVE\textcommabelow{T}I]{
  calloutbase,
  colframe=outFg,
  title={\faBullseye\ \uppercase\expandafter{#1}},
}
\newtcolorbox{whymatters}[1][DE CE]{
  calloutbase,
  colframe=whyFg,
  title={\faQuestionCircle\ \uppercase\expandafter{#1}},
}
\newtcolorbox{selfcheck}[1][PROVOCARE]{
  calloutbase,
  colframe=checkFg,
  title={\faTasks\ \uppercase\expandafter{#1}},
}

% --- Bonus / Avansat (continut auto-studiu, distinct vizual) ----
% Diferit fata de celelalte callout-uri: fundal usor colorat (nu pageBg),
% bara stanga mai groasa, icon racheta. Beginners vad clar si pot sari peste.
\definecolor{bonusBg}{HTML}{F6EEFB}        \definecolor{bonusFg}{HTML}{6B21A8}

\newtcolorbox{bonusbox}[1][BONUS \,\textbullet\, AVANSAT]{
  enhanced,
  breakable,
  sharp corners,
  boxrule=0pt,
  leftrule=4pt,
  arc=0pt,
  colback=bonusBg,
  colframe=themeBonus,
  fonttitle=\sffamily\small\bfseries,
  coltitle=bonusFg,
  attach title to upper={\par\vspace{2pt}},
  title after break={\faRocket\ \textbf{\color{bonusFg}continuare bonus}},
  before skip=14pt,
  after skip=14pt,
  left=14pt, right=6pt, top=6pt, bottom=6pt,
  fontupper=\small,
  title={\faRocket\ \uppercase\expandafter{#1}},
}

% Bridge la finalul fiecarui capitol
\newcommand{\bridge}[1]{%
  \par\vspace{4pt}%
  \begin{flushleft}%
    {\sffamily\small\color{partcolor}$\rightarrow$\ \textbf{Urmeaza:}}\quad{\small\color{black!70}#1}%
  \end{flushleft}%
}

% --- Asa da / Asa nu (compare bune vs rele practici) ----
\definecolor{daBg}{HTML}{DCFCE7}        \definecolor{daFg}{HTML}{15803D}
\definecolor{nuBg}{HTML}{FEE2E2}        \definecolor{nuFg}{HTML}{B91C1C}

\newtcolorbox{dado}[1][A\textcommabelow{S}A DA]{
  calloutbase,
  colframe=daFg,
  title={\faCheckCircle\ \uppercase\expandafter{#1}},
}
\newtcolorbox{dano}[1][A\textcommabelow{S}A NU]{
  calloutbase,
  colframe=nuFg,
  title={\faTimesCircle\ \uppercase\expandafter{#1}},
}

% =========================================================
%   BANNER MINIMALIST (\handout{titlu}{subtitlu})
%   Eticheta sus, titlu mare, subtitlu, linie subtila colorata.
% =========================================================
\newcommand{\handout}[2]{%
  \noindent\begin{minipage}{0.04\linewidth}\color{partcolor}\rule{4pt}{2.6em}\end{minipage}%
  \begin{minipage}{0.94\linewidth}
    {\sffamily\fontsize{30}{34}\selectfont\textbf{\color{themeText}#1}}\\[4pt]
    {\sffamily\large\color{black!55}#2}
  \end{minipage}
  \par\vspace{14pt}
  {\color{partcolor!40}\hrule height 0.5pt}
  \par\vspace{14pt}
}

% =========================================================
%   BULLETS COLORATI (itemize)
% =========================================================
\setlist[itemize,1]{label={\color{partcolor}\textbf{\textbullet}}}
\setlist[itemize,2]{label={\color{partcolor!70}\textbf{\textendash}}}
\setlist[enumerate,1]{label={\color{partcolor}\textbf{\arabic*.}}}

% =========================================================
%   PULL QUOTE / ACCENT
% =========================================================
\newcommand{\pullquote}[1]{%
  \begin{tcolorbox}[
    enhanced,
    colback=partcolor!8,
    colframe=partcolor,
    boxrule=0pt,
    leftrule=4pt,
    sharp corners,
    arc=0pt,
    fontupper=\itshape\large,
    left=14pt, right=14pt, top=8pt, bottom=8pt,
    overlay={%
      \node[partcolor!40, font=\fontsize{40}{40}\selectfont, anchor=north west]
        at ([xshift=2pt,yshift=2pt]frame.north west) {\textbf{``}};
    },
  ]
    #1
  \end{tcolorbox}
}

% =========================================================
%   DIVIDER DECORATIV (foloseste \partdivider in document)
%   Atentie: NU folosi numele "\sectionbreak" -- e rezervat
%   intern de titlesec si va aparea automat intre sectiuni.
% =========================================================
\newcommand{\partdivider}{%
  \par\vspace{0.6em}
  \begin{center}
    \tikz[baseline=-0.5ex]{
      \fill[partcolor] (0,0) circle (2pt);
      \fill[partcolor!60] (10pt,0) circle (2pt);
      \fill[partcolor!30] (20pt,0) circle (2pt);
    }
  \end{center}
  \vspace{0.3em}
}

% =========================================================
%   BOX MODEL DIAGRAM (TikZ)
% =========================================================
\newcommand{\boxmodelfigure}{%
\begin{center}
\begin{tikzpicture}[scale=1, transform shape, every node/.style={font=\sffamily\small}]
  % Margin (outermost)
  \fill[orange!25, rounded corners=4pt] (0,0) rectangle (8,5);
  \node[anchor=north west] at (0.15, 4.85) {\textbf{\color{orange!70!black}\faExpand\ margin}};

  % Border
  \fill[gray!40, rounded corners=3pt] (1,0.7) rectangle (7, 4.3);
  \draw[ultra thick, gray!70!black, rounded corners=3pt] (1,0.7) rectangle (7,4.3);
  \node[anchor=north west] at (1.1, 4.2) {\textbf{\color{gray!50!black}\faSquare\ border}};

  % Padding
  \fill[green!25, rounded corners=2pt] (1.3, 1) rectangle (6.7, 4);
  \node[anchor=north west] at (1.4, 3.9) {\textbf{\color{green!50!black}\faCompress\ padding}};

  % Content (innermost)
  \fill[blue!25, rounded corners=2pt] (2.2, 1.7) rectangle (5.8, 3.3);
  \node[anchor=north west] at (2.3, 3.2) {\textbf{\color{blue!60!black}\faFile*\ content}};
  \node at (4, 2.5) {\textit{textul, imaginea, ce e inauntru}};
\end{tikzpicture}
\end{center}
}

% =========================================================
%   CARD "ICON + TEXT" (utilitate pentru handout-uri)
% =========================================================
\newcommand{\iconcard}[3]{%
  \begin{tcolorbox}[
    enhanced,
    colback=partcolor!8,
    colframe=partcolor!30,
    boxrule=0.5pt,
    arc=6pt,
    left=12pt, right=12pt, top=8pt, bottom=8pt,
  ]
  \begin{minipage}[t]{0.08\linewidth}
    \color{partcolor}\fontsize{20}{20}\selectfont #1
  \end{minipage}%
  \begin{minipage}[t]{0.9\linewidth}
    \textbf{\color{partcolor!85!black}#2}\\
    #3
  \end{minipage}
  \end{tcolorbox}%
}
 dupa \documentclass.

\usepackage[utf8]{inputenc}
\usepackage[T1]{fontenc}
% Body font: Lato -- modern, lizibil, profesional. Fallback la TeX Gyre Heros
% (Helvetica-like, mereu disponibil) daca Lato nu e instalat in MiKTeX.
\IfFileExists{lato.sty}{%
  \usepackage[default]{lato}%
}{%
  \usepackage{tgheros}%
  \renewcommand{\familydefault}{\sfdefault}%
}
\usepackage[romanian]{babel}
\usepackage[a4paper, margin=2.5cm, headheight=16pt, top=2.4cm, bottom=2.4cm]{geometry}
\usepackage[final, protrusion=true, expansion=false]{microtype}
\usepackage{setspace}
\usepackage{xcolor}
\usepackage{amssymb}
\usepackage{listings}
\usepackage{hyperref}
\usepackage{enumitem}
\usepackage{titlesec}
\usepackage{fancyhdr}
\usepackage{parskip}
\usepackage[most]{tcolorbox}
\usepackage{tocloft}
\usepackage{graphicx}
\usepackage{tikz}
\usepackage{fontawesome5}
\usepackage{lettrine}
\usetikzlibrary{shadows.blur, shapes.geometric, positioning, calc,
  decorations.pathmorphing, patterns, fadings, backgrounds}

% =========================================================
%   PALETA DE CULORI
% =========================================================

% Paleta finala (3 BEST + 3 derivate in acelasi ton)
\definecolor{themePrereq}{HTML}{5C6B7A}    % slate (pre-rechizite, neutral)
\definecolor{themeHtml}{HTML}{FAA51E}      % BEST orange
\definecolor{themeCss}{HTML}{0673BA}       % BEST blue
\definecolor{themeJs}{HTML}{ECC100}        % gold (analog cu orange, distinct)
\definecolor{themeDeploy}{HTML}{75C044}    % BEST green
\definecolor{themeSeo}{HTML}{B344C9}       % violet (analog cu blue)
\definecolor{themeLighthouse}{HTML}{DD4444}% coral red (complementar cu green)
\definecolor{themeResurse}{HTML}{5C6B7A}   % slate (resurse, neutral)
\definecolor{themeBonus}{HTML}{7E22CE}     % deep purple (bonus / avansat, self-study only)
% Aliasuri pentru compatibilitate (folosit in main.tex existent)
\colorlet{themeIntro}{themeHtml}
\colorlet{themeExercitiu}{themeLighthouse}

% Background pagina si elemente -- calduros
\definecolor{pageBg}{HTML}{FFFCF6}         % crema foarte subtila
\definecolor{accentSoft}{HTML}{F4ECDD}     % accent cald
\definecolor{textDim}{HTML}{64748B}        % gri pentru text secundar

% Culoarea curenta a partii (suprascrisa de fiecare fisier prin \setpart)
\colorlet{partcolor}{themeIntro}

% Iconita partii curente (suprascrisa prin \setpart)
\newcommand{\particon}{\faStar}
\newcommand{\partlabel}{Training Web}

% \setpart{culoare}{iconita}
\newcommand{\setpart}[2]{%
  \colorlet{partcolor}{#1}%
  \renewcommand{\particon}{#2}%
}
% \setpart{culoare}{iconita}{eticheta} -- versiunea cu eticheta custom in header
\newcommand{\setpartlabel}[3]{%
  \colorlet{partcolor}{#1}%
  \renewcommand{\particon}{#2}%
  \renewcommand{\partlabel}{#3}%
}

% Culori pentru callout boxes
\definecolor{noteBg}{HTML}{FEF9C3}        \definecolor{noteFg}{HTML}{A16207}
\definecolor{tipBg}{HTML}{DBEAFE}         \definecolor{tipFg}{HTML}{1D4ED8}
\definecolor{warnBg}{HTML}{FEE2E2}        \definecolor{warnFg}{HTML}{B91C1C}
\definecolor{gameBg}{HTML}{FFEDD5}        \definecolor{gameFg}{HTML}{C2410C}
\definecolor{recapBg}{HTML}{DCFCE7}       \definecolor{recapFg}{HTML}{15803D}
\definecolor{exerBg}{HTML}{F3E8FF}        \definecolor{exerFg}{HTML}{7C3AED}
\definecolor{triviaBg}{HTML}{FCE7F3}      \definecolor{triviaFg}{HTML}{BE185D}

% Culori pentru cod -- paleta "GitHub Light" (calm, profesional)
\definecolor{codeBg}{HTML}{F6F8FA}
\definecolor{codeBorder}{HTML}{D0D7DE}
\definecolor{codeComment}{HTML}{6E7781}
\definecolor{codeKeyword}{HTML}{CF222E}
\definecolor{codeString}{HTML}{0A3069}
\definecolor{codeTag}{HTML}{116329}

% =========================================================
%   PAGINA & LIZIBILITATE (warmth)
% =========================================================
\pagecolor{pageBg}
\setstretch{1.18}
\setlength{\parskip}{0.5em}

% =========================================================
%   CUPRINS (TOC)
% =========================================================
\setcounter{tocdepth}{0}         % numai entry-uri \part (capitole), fara sectiuni
\renewcommand{\contentsname}{Cuprins}

% Stil cuprins -- numai capitolele, asezate elegant
\renewcommand{\cftpartfont}{\sffamily\bfseries\color{themeText}}
\renewcommand{\cftpartpagefont}{\sffamily\bfseries\color{themeText}}
\renewcommand{\cftpartleader}{\hfill}
\renewcommand{\cftpartafterpnum}{\par\addvspace{2pt}}
\setlength{\cftbeforepartskip}{4pt}

% =========================================================
%   HYPERREF
% =========================================================
\hypersetup{
  colorlinks=true,
  linkcolor=partcolor,
  urlcolor=partcolor,
  citecolor=partcolor,
  pdfborderstyle={/S/U/W 1}
}

% =========================================================
%   HEADER & FOOTER -- minimaliste
% =========================================================
\pagestyle{fancy}
\fancyhf{}
\fancyhead[L]{\footnotesize\sffamily\color{black!55}\partlabel}
\fancyhead[R]{\footnotesize\sffamily\color{black!55}BEST Ia\textcommabelow{s}i}
\fancyfoot[C]{\small\sffamily\color{partcolor}\textbf{\thepage}}
\renewcommand{\headrulewidth}{0pt}
\renewcommand{\footrulewidth}{0pt}

% =========================================================
%   STIL TITLURI -- sobru, stiintific
% =========================================================
\titleformat{\section}
  {\Large\sffamily\bfseries\color{themeText}}
  {\color{partcolor}\sffamily\bfseries\thesection.\hspace{0.5em}}
  {0pt}{}
\titlespacing*{\section}{0pt}{1.4em}{0.4em}

\titleformat{\subsection}
  {\normalsize\sffamily\bfseries\color{black!75}}
  {\color{partcolor}\thesubsection\hspace{0.5em}}{0pt}{}
\titlespacing*{\subsection}{0pt}{0.8em}{0.2em}

\titleformat{\subsubsection}
  {\normalsize\sffamily\bfseries\color{black!70}}
  {}{0pt}{}

\definecolor{themeText}{HTML}{1A1A1A}

% =========================================================
%   LISTINGS -- COD
% =========================================================
\lstdefinestyle{base}{
  backgroundcolor=\color{codeBg},
  basicstyle=\ttfamily\small,
  commentstyle=\color{codeComment}\itshape,
  keywordstyle=\color{codeKeyword}\bfseries,
  stringstyle=\color{codeString},
  numberstyle=\tiny\color{codeComment},
  numbers=left,
  numbersep=8pt,
  frame=leftline,
  framerule=2.5pt,
  rulecolor=\color{partcolor},
  framexleftmargin=4pt,
  breaklines=true,
  breakatwhitespace=true,
  showstringspaces=false,
  tabsize=2,
  xleftmargin=14pt,
  xrightmargin=4pt,
  aboveskip=10pt,
  belowskip=10pt,
  literate=
    {ă}{{\u a}}1 {Ă}{{\u A}}1
    {â}{{\^a}}1 {Â}{{\^A}}1
    {î}{{\^\i}}1 {Î}{{\^I}}1
    {ș}{{\textcommabelow s}}1 {Ș}{{\textcommabelow S}}1
    {ț}{{\textcommabelow t}}1 {Ț}{{\textcommabelow T}}1
}

\lstdefinelanguage{HTML5}{
  morekeywords={html,head,body,title,meta,link,script,style,
    h1,h2,h3,h4,h5,h6,p,a,img,div,span,
    header,main,section,article,aside,nav,footer,
    ul,ol,li,button,input,form,label,
    strong,em,br,hr,DOCTYPE},
  sensitive=true,
  morecomment=[s]{<!--}{-->},
  morestring=[b]",
  morestring=[b]'
}
\lstdefinestyle{html}{style=base, language=HTML5,
  keywordstyle=\color{codeTag}\bfseries}

\lstdefinelanguage{CSSLang}{
  morekeywords={display,flex,grid,block,inline,none,
    color,background,background-color,
    padding,padding-top,padding-right,padding-bottom,padding-left,
    margin,margin-top,margin-right,margin-bottom,margin-left,
    border,border-radius,border-color,outline,
    width,max-width,min-width,height,max-height,min-height,
    font-family,font-size,font-weight,line-height,text-align,
    text-decoration,letter-spacing,
    flex-direction,justify-content,align-items,gap,flex-wrap,flex-grow,
    grid-template-columns,grid-template-rows,grid-column,grid-row,
    position,top,right,bottom,left,z-index,
    box-sizing,opacity,transition,transform,
    cursor,overflow,visibility,
    media,keyframes,from,to,important,root,hover,focus,active,
    auto,solid,dashed,dotted,bold,normal,center,stretch,
    space-between,space-around,flex-start,flex-end,column,row,wrap},
  sensitive=true,
  morecomment=[s]{/*}{*/},
  morestring=[b]",
  morestring=[b]'
}
\lstdefinestyle{css}{style=base, language=CSSLang,
  keywordstyle=\color{codeKeyword}\bfseries}

\lstdefinelanguage{JSLang}{
  morekeywords={let,const,var,function,return,if,else,for,while,
    true,false,null,undefined,new,this,class,extends,
    document,console,window,
    addEventListener,querySelector,querySelectorAll,
    getElementById,getElementsByClassName,
    classList,toggle,add,remove,contains,
    style,textContent,innerHTML,value,
    log,warn,error},
  sensitive=true,
  morecomment=[l]{//},
  morecomment=[s]{/*}{*/},
  morestring=[b]",
  morestring=[b]',
  morestring=[b]`
}
\lstdefinestyle{js}{style=base, language=JSLang,
  keywordstyle=\color{codeKeyword}\bfseries}

\lstdefinelanguage{BashLang}{
  morekeywords={git,cd,mkdir,ls,code,echo,cat,curl,wget,sudo,
    init,add,commit,push,pull,clone,remote,branch,checkout,merge,status,log,
    config,fetch,rebase,reset,stash,diff,tag,
    install,update,run,build,start,test},
  sensitive=true,
  morecomment=[l]{\#},
  morestring=[b]",
  morestring=[b]'
}
\lstdefinestyle{bash}{style=base, language=BashLang,
  keywordstyle=\color{codeKeyword}\bfseries}

% Inline code -- mai discret
\newcommand{\code}[1]{\textcolor{partcolor!85!black}{\ttfamily\small #1}}

% =========================================================
%   CALLOUT BOXES -- minimaliste (bara stanga subtila + label discret)
% =========================================================
\tcbset{
  calloutbase/.style={
    enhanced,
    breakable,
    sharp corners,
    boxrule=0pt,
    leftrule=2.5pt,
    arc=0pt,
    colback=pageBg,
    fonttitle=\sffamily\small\bfseries,
    coltitle=tcbcolframe,
    attach title to upper={\par\vspace{2pt}},
    title after break={},
    before skip=10pt,
    after skip=10pt,
    left=12pt, right=4pt, top=4pt, bottom=4pt,
    fontupper=\small,
  }
}

\newtcolorbox{nota}[1][NOT\u A]{
  calloutbase,
  colframe=noteFg,
  title={\faInfoCircle\ \uppercase\expandafter{#1}},
}
\newtcolorbox{ponturi}[1][PONT]{
  calloutbase,
  colframe=tipFg,
  title={\faLightbulb[regular]\ \uppercase\expandafter{#1}},
}
\newtcolorbox{atentie}[1][ATEN\textcommabelow{T}IE]{
  calloutbase,
  colframe=warnFg,
  title={\faExclamationTriangle\ \uppercase\expandafter{#1}},
}
\newtcolorbox{joc}[1][APLICA\textcommabelow{T}IE PRACTIC\u A]{
  calloutbase,
  colframe=partcolor,
  title={\faGamepad\ \uppercase\expandafter{#1}},
}
\newtcolorbox{recap}[1][RECAP]{
  calloutbase,
  colframe=recapFg,
  title={\faCheckCircle\ \uppercase\expandafter{#1}},
}
\newtcolorbox{curiozitate}[1][\textcommabelow{S}TIAI C\u A?]{
  calloutbase,
  colframe=triviaFg,
  title={\faMagic\ \uppercase\expandafter{#1}},
}
\newtcolorbox{exercitiu}[1][EXERCI\textcommabelow{T}IU]{
  calloutbase,
  colframe=exerFg,
  title={\faPen\ \uppercase\expandafter{#1}},
}

% --- Outcome / Why / Self-check (structura tutorial) ----
\definecolor{outBg}{HTML}{E0F2FE}        \definecolor{outFg}{HTML}{0369A1}
\definecolor{whyBg}{HTML}{FEF3C7}        \definecolor{whyFg}{HTML}{B45309}
\definecolor{checkBg}{HTML}{F3E8FF}      \definecolor{checkFg}{HTML}{7C3AED}

\newtcolorbox{outcomes}[1][CE INVE\textcommabelow{T}I]{
  calloutbase,
  colframe=outFg,
  title={\faBullseye\ \uppercase\expandafter{#1}},
}
\newtcolorbox{whymatters}[1][DE CE]{
  calloutbase,
  colframe=whyFg,
  title={\faQuestionCircle\ \uppercase\expandafter{#1}},
}
\newtcolorbox{selfcheck}[1][PROVOCARE]{
  calloutbase,
  colframe=checkFg,
  title={\faTasks\ \uppercase\expandafter{#1}},
}

% --- Bonus / Avansat (continut auto-studiu, distinct vizual) ----
% Diferit fata de celelalte callout-uri: fundal usor colorat (nu pageBg),
% bara stanga mai groasa, icon racheta. Beginners vad clar si pot sari peste.
\definecolor{bonusBg}{HTML}{F6EEFB}        \definecolor{bonusFg}{HTML}{6B21A8}

\newtcolorbox{bonusbox}[1][BONUS \,\textbullet\, AVANSAT]{
  enhanced,
  breakable,
  sharp corners,
  boxrule=0pt,
  leftrule=4pt,
  arc=0pt,
  colback=bonusBg,
  colframe=themeBonus,
  fonttitle=\sffamily\small\bfseries,
  coltitle=bonusFg,
  attach title to upper={\par\vspace{2pt}},
  title after break={\faRocket\ \textbf{\color{bonusFg}continuare bonus}},
  before skip=14pt,
  after skip=14pt,
  left=14pt, right=6pt, top=6pt, bottom=6pt,
  fontupper=\small,
  title={\faRocket\ \uppercase\expandafter{#1}},
}

% Bridge la finalul fiecarui capitol
\newcommand{\bridge}[1]{%
  \par\vspace{4pt}%
  \begin{flushleft}%
    {\sffamily\small\color{partcolor}$\rightarrow$\ \textbf{Urmeaza:}}\quad{\small\color{black!70}#1}%
  \end{flushleft}%
}

% --- Asa da / Asa nu (compare bune vs rele practici) ----
\definecolor{daBg}{HTML}{DCFCE7}        \definecolor{daFg}{HTML}{15803D}
\definecolor{nuBg}{HTML}{FEE2E2}        \definecolor{nuFg}{HTML}{B91C1C}

\newtcolorbox{dado}[1][A\textcommabelow{S}A DA]{
  calloutbase,
  colframe=daFg,
  title={\faCheckCircle\ \uppercase\expandafter{#1}},
}
\newtcolorbox{dano}[1][A\textcommabelow{S}A NU]{
  calloutbase,
  colframe=nuFg,
  title={\faTimesCircle\ \uppercase\expandafter{#1}},
}

% =========================================================
%   BANNER MINIMALIST (\handout{titlu}{subtitlu})
%   Eticheta sus, titlu mare, subtitlu, linie subtila colorata.
% =========================================================
\newcommand{\handout}[2]{%
  \noindent\begin{minipage}{0.04\linewidth}\color{partcolor}\rule{4pt}{2.6em}\end{minipage}%
  \begin{minipage}{0.94\linewidth}
    {\sffamily\fontsize{30}{34}\selectfont\textbf{\color{themeText}#1}}\\[4pt]
    {\sffamily\large\color{black!55}#2}
  \end{minipage}
  \par\vspace{14pt}
  {\color{partcolor!40}\hrule height 0.5pt}
  \par\vspace{14pt}
}

% =========================================================
%   BULLETS COLORATI (itemize)
% =========================================================
\setlist[itemize,1]{label={\color{partcolor}\textbf{\textbullet}}}
\setlist[itemize,2]{label={\color{partcolor!70}\textbf{\textendash}}}
\setlist[enumerate,1]{label={\color{partcolor}\textbf{\arabic*.}}}

% =========================================================
%   PULL QUOTE / ACCENT
% =========================================================
\newcommand{\pullquote}[1]{%
  \begin{tcolorbox}[
    enhanced,
    colback=partcolor!8,
    colframe=partcolor,
    boxrule=0pt,
    leftrule=4pt,
    sharp corners,
    arc=0pt,
    fontupper=\itshape\large,
    left=14pt, right=14pt, top=8pt, bottom=8pt,
    overlay={%
      \node[partcolor!40, font=\fontsize{40}{40}\selectfont, anchor=north west]
        at ([xshift=2pt,yshift=2pt]frame.north west) {\textbf{``}};
    },
  ]
    #1
  \end{tcolorbox}
}

% =========================================================
%   DIVIDER DECORATIV (foloseste \partdivider in document)
%   Atentie: NU folosi numele "\sectionbreak" -- e rezervat
%   intern de titlesec si va aparea automat intre sectiuni.
% =========================================================
\newcommand{\partdivider}{%
  \par\vspace{0.6em}
  \begin{center}
    \tikz[baseline=-0.5ex]{
      \fill[partcolor] (0,0) circle (2pt);
      \fill[partcolor!60] (10pt,0) circle (2pt);
      \fill[partcolor!30] (20pt,0) circle (2pt);
    }
  \end{center}
  \vspace{0.3em}
}

% =========================================================
%   BOX MODEL DIAGRAM (TikZ)
% =========================================================
\newcommand{\boxmodelfigure}{%
\begin{center}
\begin{tikzpicture}[scale=1, transform shape, every node/.style={font=\sffamily\small}]
  % Margin (outermost)
  \fill[orange!25, rounded corners=4pt] (0,0) rectangle (8,5);
  \node[anchor=north west] at (0.15, 4.85) {\textbf{\color{orange!70!black}\faExpand\ margin}};

  % Border
  \fill[gray!40, rounded corners=3pt] (1,0.7) rectangle (7, 4.3);
  \draw[ultra thick, gray!70!black, rounded corners=3pt] (1,0.7) rectangle (7,4.3);
  \node[anchor=north west] at (1.1, 4.2) {\textbf{\color{gray!50!black}\faSquare\ border}};

  % Padding
  \fill[green!25, rounded corners=2pt] (1.3, 1) rectangle (6.7, 4);
  \node[anchor=north west] at (1.4, 3.9) {\textbf{\color{green!50!black}\faCompress\ padding}};

  % Content (innermost)
  \fill[blue!25, rounded corners=2pt] (2.2, 1.7) rectangle (5.8, 3.3);
  \node[anchor=north west] at (2.3, 3.2) {\textbf{\color{blue!60!black}\faFile*\ content}};
  \node at (4, 2.5) {\textit{textul, imaginea, ce e inauntru}};
\end{tikzpicture}
\end{center}
}

% =========================================================
%   CARD "ICON + TEXT" (utilitate pentru handout-uri)
% =========================================================
\newcommand{\iconcard}[3]{%
  \begin{tcolorbox}[
    enhanced,
    colback=partcolor!8,
    colframe=partcolor!30,
    boxrule=0.5pt,
    arc=6pt,
    left=12pt, right=12pt, top=8pt, bottom=8pt,
  ]
  \begin{minipage}[t]{0.08\linewidth}
    \color{partcolor}\fontsize{20}{20}\selectfont #1
  \end{minipage}%
  \begin{minipage}[t]{0.9\linewidth}
    \textbf{\color{partcolor!85!black}#2}\\
    #3
  \end{minipage}
  \end{tcolorbox}%
}

\setpartlabel{themeHtml}{\faStar}{Training Web}

\begin{document}

% =========================================================
%   COPERTA
% =========================================================
\begin{titlepage}
  \thispagestyle{empty}
  \vspace*{0.5cm}

  \begin{tcolorbox}[
    enhanced,
    colback=themeHtml,
    colframe=themeHtml,
    boxrule=0pt,
    arc=14pt,
    rounded corners,
    drop shadow={themeHtml!40!black, opacity=0.4},
    left=28pt, right=28pt, top=32pt, bottom=32pt,
    interior style={
      top color=themeHtml!85!white,
      bottom color=themeHtml,
    },
    overlay={%
      \begin{tcbclipinterior}%
        \fill[white, opacity=0.10] ([xshift=-3cm,yshift=-2cm]frame.north east) circle (2.5cm);
        \fill[white, opacity=0.06] ([xshift=-1cm,yshift=-0.5cm]frame.north east) circle (1.5cm);
        \fill[white, opacity=0.08] ([xshift=2cm,yshift=2cm]frame.south west) circle (2cm);
        \node[white, opacity=0.12, anchor=east,
              font=\fontsize{180}{180}\selectfont]
          at ([xshift=-1cm]frame.east) {\faCode};
      \end{tcbclipinterior}%
    },
  ]
    \begin{flushleft}
    {\color{white!90}\sffamily\large\textbf{\faSquare\ BEST Ia\textcommabelow{s}i \,\textbullet\, Workshop 2026}}\\[18pt]
    {\color{white}\sffamily\fontsize{52}{60}\selectfont\textbf{Training Web}}\\[8pt]
    {\color{white!92}\sffamily\Huge De la zero la site live in 2 ore (hopefully)}\\[20pt]
    {\color{white!85}\sffamily\large\itshape HTML \,\textbullet\, CSS \,\textbullet\, JavaScript \,\textbullet\, Deploy \,\textbullet\, SEO \,\textbullet\, Lighthouse}
    \end{flushleft}
  \end{tcolorbox}

  \vspace{2em}

  \begin{center}
    \sffamily\large\color{black!75}
    \textbf{Suport pentru participanti \textendash{} 6 faze, fiecare cu aplicatie practica}
  \end{center}

  \vspace{1.5em}

  % Cardurile celor 7 parti
  \noindent
  \begin{tcolorbox}[
    enhanced, colback=themeHtml!8, colframe=themeHtml, boxrule=0pt, leftrule=5pt,
    arc=4pt, sharp corners=downhill, rounded corners=northwest, left=10pt, right=10pt, top=6pt, bottom=6pt,
  ]
    \begin{minipage}[c]{0.06\linewidth}\color{themeHtml}\Large\faCode\end{minipage}\hfill
    \begin{minipage}[c]{0.92\linewidth}\textbf{\color{themeHtml}Faza 1 \textbar{} HTML} \,\textendash\, structura, semantica, accesibilitate\end{minipage}
  \end{tcolorbox}

  \noindent
  \begin{tcolorbox}[
    enhanced, colback=themeCss!8, colframe=themeCss, boxrule=0pt, leftrule=5pt,
    arc=4pt, sharp corners=downhill, rounded corners=northwest, left=10pt, right=10pt, top=6pt, bottom=6pt,
  ]
    \begin{minipage}[c]{0.06\linewidth}\color{themeCss}\Large\faPaintBrush\end{minipage}\hfill
    \begin{minipage}[c]{0.92\linewidth}\textbf{\color{themeCss}Faza 2 \textbar{} CSS} \,\textendash\, stilizare, layout, responsivitate\end{minipage}
  \end{tcolorbox}

  \noindent
  \begin{tcolorbox}[
    enhanced, colback=themeJs!8, colframe=themeJs, boxrule=0pt, leftrule=5pt,
    arc=4pt, sharp corners=downhill, rounded corners=northwest, left=10pt, right=10pt, top=6pt, bottom=6pt,
  ]
    \begin{minipage}[c]{0.06\linewidth}\color{themeJs}\Large\faBolt\end{minipage}\hfill
    \begin{minipage}[c]{0.92\linewidth}\textbf{\color{themeJs}Faza 3 \textbar{} JavaScript} \,\textendash\, interactivitate, manipulare DOM\end{minipage}
  \end{tcolorbox}

  \noindent
  \begin{tcolorbox}[
    enhanced, colback=themeDeploy!8, colframe=themeDeploy, boxrule=0pt, leftrule=5pt,
    arc=4pt, sharp corners=downhill, rounded corners=northwest, left=10pt, right=10pt, top=6pt, bottom=6pt,
  ]
    \begin{minipage}[c]{0.06\linewidth}\color{themeDeploy}\Large\faRocket\end{minipage}\hfill
    \begin{minipage}[c]{0.92\linewidth}\textbf{\color{themeDeploy}Faza 4 \textbar{} Deploy} \,\textendash\, versionare git, hosting GitHub Pages\end{minipage}
  \end{tcolorbox}

  \noindent
  \begin{tcolorbox}[
    enhanced, colback=themeSeo!8, colframe=themeSeo, boxrule=0pt, leftrule=5pt,
    arc=4pt, sharp corners=downhill, rounded corners=northwest, left=10pt, right=10pt, top=6pt, bottom=6pt,
  ]
    \begin{minipage}[c]{0.06\linewidth}\color{themeSeo}\Large\faShareSquare\end{minipage}\hfill
    \begin{minipage}[c]{0.92\linewidth}\textbf{\color{themeSeo}Faza 5 \textbar{} SEO + Open Graph} \,\textendash\, metadate pentru cautare si retele sociale\end{minipage}
  \end{tcolorbox}

  \noindent
  \begin{tcolorbox}[
    enhanced, colback=themeLighthouse!8, colframe=themeLighthouse, boxrule=0pt, leftrule=5pt,
    arc=4pt, sharp corners=downhill, rounded corners=northwest, left=10pt, right=10pt, top=6pt, bottom=6pt,
  ]
    \begin{minipage}[c]{0.06\linewidth}\color{themeLighthouse}\Large\faChartBar\end{minipage}\hfill
    \begin{minipage}[c]{0.92\linewidth}\textbf{\color{themeLighthouse}Faza 6 \textbar{} Lighthouse} \,\textendash\, audit automatizat, Core Web Vitals\end{minipage}
  \end{tcolorbox}

  \noindent
  \begin{tcolorbox}[
    enhanced, colback=themeResurse!8, colframe=themeResurse, boxrule=0pt, leftrule=5pt,
    arc=4pt, sharp corners=downhill, rounded corners=northwest, left=10pt, right=10pt, top=6pt, bottom=6pt,
  ]
    \begin{minipage}[c]{0.06\linewidth}\color{themeResurse}\Large\faGraduationCap\end{minipage}\hfill
    \begin{minipage}[c]{0.92\linewidth}\textbf{\color{themeResurse}Resurse \& urmatorii pasi} \,\textendash\, aplicatii interactive, tutoriale, comunitati\end{minipage}
  \end{tcolorbox}

  \vfill
  \begin{center}
    \small\color{black!60}\sffamily
    Pachetul serveste atat ca suport in timpul training-ului, cat si ca referinta de aprofundare ulterioara.
  \end{center}
\end{titlepage}

% =========================================================
%   PAGINA 2 -- CUM SE CITESC MATERIALELE
% =========================================================

\section*{\faBookOpen\ \ Structura manualului}

Manualul prezent contine 7 capitole tematice (sase faze ale training-ului plus un capitol final cu resurse), structurate uniform pentru consultare facila. Fiecare capitol primeste o culoare distincta, vizibila in bara verticala stanga a blocurilor de cod si in numerotarea sectiunilor.

\bigskip

\begin{nota}[Structura unui capitol]
\begin{itemize}
  \item \textbf{Teorie sintetica} \textendash{} concepte, sintaxa, referinte la standarde.
  \item \textbf{Bloc de cod} \textendash{} fragmente reutilizabile direct prin copiere.
  \item \textbf{\faGamepad\ Aplicatie practica} \textendash{} platforma online cu sarcina de exersare (5\textendash{}10 minute).
  \item \textbf{\faGraduationCap\ Resurse pentru aprofundare} \textendash{} bibliografie cu linkuri.
\end{itemize}
\end{nota}

\begin{ponturi}[Conventii vizuale]
\begin{itemize}
  \item Blocurile de cod sunt marcate printr-o bara verticala colorata in stanga, in culoarea capitolului.
  \item Toate URL-urile sunt clicabile in vizualizatorul PDF.
  \item Iconitele \faCode\ \faPaintBrush\ \faBolt\ \faRocket\ \faShareSquare\ \faChartBar\ identifica vizual fiecare capitol.
\end{itemize}
\end{ponturi}

% =========================================================
%   PROIECTUL FIR ROSU
% =========================================================
\clearpage
\setpartlabel{themeHtml}{\faMap}{Proiectul fir ro\textcommabelow{s}u}

\noindent\begin{minipage}{0.04\linewidth}\color{themeHtml}\rule{4pt}{2.6em}\end{minipage}%
\begin{minipage}{0.94\linewidth}
  {\sffamily\fontsize{30}{34}\selectfont\textbf{\color{themeText}Proiectul fir ro\textcommabelow{s}u}}\\[4pt]
  {\sffamily\large\color{black!55}Construim impreuna portofoliul Voluntarului BEST.}
\end{minipage}
\par\vspace{14pt}
{\color{themeHtml!40}\hrule height 0.5pt}
\par\vspace{14pt}

Manualul nu prezinta concepte izolate. Pe parcursul celor 6 faze, fiecare strat se adauga la o pagina concreta \textendash{} portofoliul Voluntarului BEST, studenta CS in Iasi. La final, fiecare participant aplica acelasi pattern pentru pagina sa proprie.

\bigskip

\textbf{\sffamily Evolu\textcommabelow{t}ia paginii pe parcursul training-ului:}

\bigskip

\begin{tcolorbox}[
  enhanced, colback=themeHtml!8, colframe=themeHtml, boxrule=0pt, leftrule=4pt,
  arc=4pt, sharp corners=downhill, rounded corners=northwest,
  left=12pt, right=12pt, top=8pt, bottom=8pt, before skip=4pt, after skip=4pt,
]
\textbf{\color{themeHtml}Faza 1 \textendash{} HTML.}\quad Text negru pe fundal alb, ca un document Word minimal. Structura semantica corecta in spate.\\
{\small\color{textDim} Pagina \emph{exista} la nivel de informa\textcommabelow{t}ie, dar arata generic.}
\end{tcolorbox}

\begin{tcolorbox}[
  enhanced, colback=themeCss!8, colframe=themeCss, boxrule=0pt, leftrule=4pt,
  arc=4pt, sharp corners=downhill, rounded corners=northwest,
  left=12pt, right=12pt, top=8pt, bottom=8pt, before skip=4pt, after skip=4pt,
]
\textbf{\color{themeCss}+ Faza 2 \textendash{} CSS.}\quad Pagina capata identitate vizuala: paleta de culori, layout cu flexbox, header cu logo \textcommabelow{s}i navigare, hero cu CTA, dark mode.\\
{\small\color{textDim} Pagina \emph{arata} ca un site profesional. Nu raspunde la interac\textcommabelow{t}iuni.}
\end{tcolorbox}

\begin{tcolorbox}[
  enhanced, colback=themeJs!8, colframe=themeJs, boxrule=0pt, leftrule=4pt,
  arc=4pt, sharp corners=downhill, rounded corners=northwest,
  left=12pt, right=12pt, top=8pt, bottom=8pt, before skip=4pt, after skip=4pt,
]
\textbf{\color{themeJs}+ Faza 3 \textendash{} JavaScript.}\quad Hamburger menu func\textcommabelow{t}ional pe mobil, smooth scroll la link-uri, formular care nu reincarca pagina.\\
{\small\color{textDim} Pagina \emph{raspunde} la utilizator. Ramane offline, accesibila doar local.}
\end{tcolorbox}

\begin{tcolorbox}[
  enhanced, colback=themeDeploy!8, colframe=themeDeploy, boxrule=0pt, leftrule=4pt,
  arc=4pt, sharp corners=downhill, rounded corners=northwest,
  left=12pt, right=12pt, top=8pt, bottom=8pt, before skip=4pt, after skip=4pt,
]
\textbf{\color{themeDeploy}+ Faza 4 \textendash{} Deploy.}\quad Site-ul devine accesibil la \code{voluntar-best.github.io/portfolio}. URL real, share-uibil oriunde.\\
{\small\color{textDim} Pagina \emph{exista in lume}. Dar oamenii inca nu o pot gasi.}
\end{tcolorbox}

\begin{tcolorbox}[
  enhanced, colback=themeSeo!8, colframe=themeSeo, boxrule=0pt, leftrule=4pt,
  arc=4pt, sharp corners=downhill, rounded corners=northwest,
  left=12pt, right=12pt, top=8pt, bottom=8pt, before skip=4pt, after skip=4pt,
]
\textbf{\color{themeSeo}+ Faza 5 \textendash{} SEO.}\quad Cand cineva paste-uieste link-ul pe WhatsApp, apare card cu titlu, descriere si imagine. Google indexeaza pagina corect.\\
{\small\color{textDim} Pagina \emph{e descoperibila}. Tehnic poate fi inca rafinata.}
\end{tcolorbox}

\begin{tcolorbox}[
  enhanced, colback=themeLighthouse!8, colframe=themeLighthouse, boxrule=0pt, leftrule=4pt,
  arc=4pt, sharp corners=downhill, rounded corners=northwest,
  left=12pt, right=12pt, top=8pt, bottom=8pt, before skip=4pt, after skip=4pt,
]
\textbf{\color{themeLighthouse}+ Faza 6 \textendash{} Lighthouse.}\quad Audit pe 4 categorii. Scor 95+ pe Performance, Accessibility, Best Practices, SEO.\\
{\small\color{textDim} Pagina \emph{e profesionala}. Standardul pe care firmele il cer la angajare.}
\end{tcolorbox}

\bigskip

\begin{recap}[Pe scurt]
Fiecare faza adauga o capacitate noua paginii. La final, paginia ta personala (cu numele tau in URL) trece toate verificarile pe care un dezvoltator profesional le-ar face.
\end{recap}

% =========================================================
%   CUPRINS
% =========================================================
\clearpage
\setpartlabel{themePrereq}{\faBookmark}{Cuprins}

\noindent\begin{minipage}{0.04\linewidth}\color{themePrereq}\rule{4pt}{2.6em}\end{minipage}%
\begin{minipage}{0.94\linewidth}
  {\sffamily\fontsize{30}{34}\selectfont\textbf{\color{themeText}Cuprins}}\\[4pt]
  {\sffamily\large\color{black!55}Capitole, in ordinea recomandata.}
\end{minipage}
\par\vspace{14pt}
{\color{themePrereq!40}\hrule height 0.5pt}
\par\vspace{20pt}

\begingroup
\renewcommand{\baselinestretch}{1.4}\selectfont
\makeatletter
\@starttoc{toc}
\makeatother
\endgroup

% =========================================================
%   CAPITOLE
% =========================================================
% Fragment -- inclus de main.tex prin \input
\clearpage
\setpartlabel{themeHtml}{\faCode}{Faza 1: HTML}
\setcounter{section}{0}
\phantomsection
\addcontentsline{toc}{part}{\protect\textbf{Faza 1 \textendash{} HTML}}
\handout{Faza 1 \textendash{} HTML}{Structura documentului. Tag-uri semantice. Accesibilitate.}

HTML (\emph{HyperText Markup Language}) este limbajul de marcare folosit pentru descrierea \textbf{structurii} unui document web. Browserul citeste documentul de sus in jos, construieste in memorie un arbore de noduri (DOM-ul), si afiseaza paginea pe ecran.

In arhitectura unei pagini web, responsabilitatile sunt impartite intre trei tehnologii:

\begin{itemize}[itemsep=0pt]
  \item \textbf{HTML} \textendash{} ce contine pagina (titluri, paragrafe, imagini, formulare).
  \item \textbf{CSS} \textendash{} cum arata pagina (culori, dimensiuni, spatiere, layout).
  \item \textbf{JavaScript} \textendash{} cum se comporta pagina (raspuns la click, animatii, validare).
\end{itemize}

Capitolul defineste anatomia minimala, tag-urile semantice principale, atributele obligatorii pentru accesibilitate, si patternurile uzuale pentru linkuri si formulare.

\bigskip

\begin{outcomes}
Dupa capitolul asta vei \textcommabelow{s}ti:
\begin{itemize}[itemsep=1pt]
  \item Scrie boilerplate-ul HTML5 corect, cu cele 3 meta tags obligatorii.
  \item Distinge tag-urile semantice de \code{<div>} \textcommabelow{s}i alege tag-ul potrivit per zona.
  \item Construie\textcommabelow{s}te formulare accesibile cu \code{<label for>} \textcommabelow{s}i \code{<input id>}.
  \item Aplica atributele \code{alt}, \code{lang}, \code{href} cu valori corecte conform WCAG.
  \item Recunoa\textcommabelow{s}te pattern-urile pentru linkuri externe sigure (\code{rel="noopener"}).
  \item \faRocket\ \emph{Bonus self-study:} \code{<picture>} cu fallback de format, \code{loading="lazy"} + \code{decoding="async"}, ARIA (\code{aria-label}, \code{aria-current}).
\end{itemize}
\end{outcomes}

\begin{whymatters}
HTML este temelia oricarei pagini web. Browserele, cititoarele de ecran \textcommabelow{s}i motoarele de cautare interpreteaza pagina pornind de la structura HTML. O structura semantica gre\textcommabelow{s}ita conduce la trei consecin\textcommabelow{t}e directe: SEO penalizat, persoanele cu dizabilita\textcommabelow{t}i nu pot folosi pagina, codul devine ilizibil dupa cateva luni. Toate fazele urmatoare (CSS, JS, deploy, audit) presupun ca structura HTML este corecta.
\end{whymatters}

% =========================================================
\section{Anatomia documentului}

Un document HTML5 contine cinci elemente structurale obligatorii. Pe linia 1, declaratia \code{<!DOCTYPE html>} indica browserului ca documentul respecta specificatia HTML5; in absenta sa, browserul intra in modul ,,quirks'' care reproduce comportamentul browserelor din anii 2000.

\begin{lstlisting}[style=html, caption={Boilerplate HTML5 conform specificatiei W3C}]
<!DOCTYPE html>
<html lang="ro">
<head>
  <meta charset="UTF-8">
  <meta name="viewport" content="width=device-width, initial-scale=1.0">
  <title>Numele Prenumele -- Student CS, Iasi</title>
</head>
<body>
  <!-- continut vizibil -->
</body>
</html>
\end{lstlisting}

\begin{itemize}
  \item \code{<html lang="ro">} \textendash{} elementul radacina. Atributul \code{lang} influenteaza pronuntia in cititoare de ecran, ranking-ul in cautari localizate, si detectia limbii sursa pentru traducere automata.
  \item \code{<head>} \textendash{} contine metadatele paginii. Continutul nu apare pe ecran, dar este consumat de browser, motoare de cautare si platforme sociale.
  \item \code{<body>} \textendash{} contine elementele vizibile pe ecran. Tot ce vede utilizatorul se afla in interiorul acestui tag.
\end{itemize}

% =========================================================
\section{Sintaxa tag-urilor}

Un element HTML are una sau mai multe parti: tag de deschidere, continut, tag de inchidere. Atributele se declara in tag-ul de deschidere, in formatul \code{nume="valoare"}.

\begin{lstlisting}[style=html, caption={Tipuri de tag-uri}]
<!-- Tag container: are continut intre deschidere si inchidere -->
<p>Acesta este un paragraf.</p>

<!-- Tag void: nu are continut, nu se inchide -->
<img src="poza.jpg" alt="Descriere">

<!-- Tag cu atribute multiple -->
<a href="https://best-is.ro" target="_blank" rel="noopener">BEST</a>
\end{lstlisting}

Tag-urile void cuprind: \code{<img>}, \code{<br>}, \code{<hr>}, \code{<input>}, \code{<meta>}, \code{<link>}.

% =========================================================
\section{Metadate obligatorii in <head>}

Trei elemente sunt considerate obligatorii in orice pagina moderna:

\begin{enumerate}
  \item \code{<meta charset="UTF-8">} \textendash{} setul de caractere. UTF-8 acopera toate alfabetele moderne, inclusiv diacriticele romanesti.
  \item \code{<meta name="viewport"...>} \textendash{} controleaza redarea pe dispozitive mobile. Valoarea \code{width=device-width, initial-scale=1.0} instructeaza browserul sa pastreze latimea paginii egala cu cea a ecranului si sa nu zoom-eze initial.
  \item \code{<title>} \textendash{} apare in tab-ul browserului, in bookmark-uri, si ca titlu in rezultatele cautarii.
\end{enumerate}

\begin{atentie}
Conform datelor StatCounter (2024), peste 60\% din traficul web global provine de pe dispozitive mobile. Lipsa meta tag-ului \code{viewport} afecteaza direct experienta majoritatii utilizatorilor, generand o pagina nescalata, cu zoom necesar pentru lectura.
\end{atentie}

% =========================================================
\section{Tag-uri semantice vs <div> soup}

Tag-urile semantice descriu \emph{rolul} continutului, nu doar aspectul vizual. Browserul, motoarele de cautare si tehnologiile asistive (cititoare de ecran, navigatoare prin tastatura) pot interpreta automat structura paginii cand tag-urile semantice sunt folosite.

\begin{dano}[Asa nu \textendash{} ,,div soup'']
\begin{lstlisting}[style=html, numbers=none]
<div class="header">
  <div class="logo">Site</div>
  <div class="menu">
    <div class="link"><a href="/">Acasa</a></div>
  </div>
</div>
<div class="content">
  <div class="hero">
    <div class="title">Salut</div>
  </div>
</div>
<div class="footer">2026</div>
\end{lstlisting}
\textbf{Probleme.} Browserul si cititoarele de ecran nu stiu ce reprezinta fiecare \code{<div>}. Navigatia rapida prin landmarks este imposibila. SEO penalizat pentru lipsa de structura.
\end{dano}

\begin{dado}[Asa da \textendash{} semantic]
\begin{lstlisting}[style=html, numbers=none]
<header>
  <a class="logo" href="/">Site</a>
  <nav>
    <a href="/">Acasa</a>
  </nav>
</header>
<main>
  <section class="hero">
    <h1>Salut</h1>
  </section>
</main>
<footer>2026</footer>
\end{lstlisting}
\textbf{Beneficii.} Cititoarele de ecran ofera comenzi rapide pentru sarit la \code{<header>}, \code{<nav>}, \code{<main>}, \code{<footer>}. Google identifica zonele principale fara analiza euristica. Codul este mai usor de citit dupa 6 luni.
\end{dado}

\subsection{Repertoriul tag-urilor semantice}

\begin{itemize}[itemsep=2pt]
  \item \code{<header>} \textendash{} antetul unei pagini sau al unei sectiuni.
  \item \code{<nav>} \textendash{} bara de navigatie principala.
  \item \code{<main>} \textendash{} continutul unic al paginii (un singur \code{<main>} per pagina).
  \item \code{<section>} \textendash{} grup tematic de continut, de obicei cu un \code{<h2>}.
  \item \code{<article>} \textendash{} unitate de continut auto-suficienta (post de blog, comentariu, card produs).
  \item \code{<aside>} \textendash{} continut tangential (sidebar, note de subsol vizuale).
  \item \code{<footer>} \textendash{} subsolul unei pagini sau al unei sectiuni.
  \item \code{<figure>} + \code{<figcaption>} \textendash{} imagine cu legenda.
\end{itemize}

% =========================================================
\section{Ierarhia titlurilor}

Specificatia HTML5 defineste sase niveluri de titluri, de la \code{<h1>} la \code{<h6>}. Cititoarele de ecran genereaza un \emph{outline} al paginii pe baza acestei ierarhii; un utilizator nevazator parcurge documentul saltand intre titluri exact ca un cititor obisnuit cand parcurge un cuprins.

\begin{itemize}
  \item Un singur \code{<h1>} per pagina (conven\textcommabelow{t}ia industriala pentru accesibilitate \textcommabelow{s}i SEO; HTML Living Standard permite tehnic mai multe, dar Lighthouse \textcommabelow{s}i cititoarele de ecran a\textcommabelow{s}teapta unul singur), descriind tema centrala.
  \item Subsectiunile principale folosesc \code{<h2>}; sub-subsectiunile, \code{<h3>}.
  \item Nivelurile nu se sar (\code{<h1>} urmat direct de \code{<h3>} este o eroare structurala).
  \item Marimea vizuala se controleaza prin CSS, nu prin alegerea unui nivel diferit.
\end{itemize}

\begin{dano}[Asa nu \textendash{} stilul ales prin nivel]
\begin{lstlisting}[style=html, numbers=none]
<h1>Titlul mare</h1>
<h4>Subtitlu, am ales h4 ca arata mai mic</h4>
<h2>Sectiune</h2>
\end{lstlisting}
\end{dano}

\begin{dado}[Asa da \textendash{} stilul ales prin CSS]
\begin{lstlisting}[style=html, numbers=none]
<h1>Titlul mare</h1>
<h2 class="subtitle">Subtitlu</h2>
<h2>Sectiune</h2>
\end{lstlisting}
\code{.subtitle \{ font-size: 1rem; font-weight: normal; \}}
\end{dado}

% =========================================================
\section{Linkuri si navigare}

Elementul \code{<a>} (anchor) creeaza un hyperlink. Atributul \code{href} accepta mai multe forme:

\begin{lstlisting}[style=html, caption={Tipuri de URL-uri}]
<!-- URL absolut -->
<a href="https://best-is.ro">BEST Iasi</a>

<!-- URL relativ -->
<a href="contact.html">Contact</a>
<a href="../images/logo.png">Logo</a>

<!-- Ancora interna -->
<a href="#contact">Sari la contact</a>

<!-- Adresa email / numar de telefon -->
<a href="mailto:hello@best-is.ro">Email</a>
<a href="tel:+40123456789">Telefon</a>
\end{lstlisting}

\subsection{Linkuri externe \textendash{} target si rel}

Deschiderea unui link extern intr-un tab nou se face cu \code{target="\_blank"}. Fara protectie, pagina destinatie poate manipula tab-ul sursa prin \code{window.opener} (tabnabbing $\to$ phishing). Adauga mereu \code{rel="noopener noreferrer"}:

\begin{lstlisting}[style=html, numbers=none]
<a href="https://extern.ro" target="_blank" rel="noopener noreferrer">Extern</a>
\end{lstlisting}

\code{noopener} blocheaza accesul la \code{window.opener}; \code{noreferrer} ascunde URL-ul sursa.

% =========================================================
\section{Formulare accesibile}

Fiecare element \code{<input>} trebuie asociat cu un \code{<label>}. Asocierea se realizeaza prin perechea atribut-valoare \code{for}/\code{id}: valoarea atributului \code{for} de pe \code{<label>} este identica cu valoarea atributului \code{id} de pe \code{<input>}.

\begin{dano}[Asa nu \textendash{} placeholder ca substitut pentru label]
\begin{lstlisting}[style=html, numbers=none]
<input type="email" placeholder="Email">
\end{lstlisting}
\textbf{Probleme.} (1) Cititoarele de ecran nu anunta placeholder-ul ca eticheta. (2) Placeholder-ul dispare la tastare, lasand utilizatorul fara context. (3) Contrast scazut pe textul de placeholder.
\end{dano}

\begin{dado}[Asa da]
\begin{lstlisting}[style=html, numbers=none]
<label for="email">Email</label>
<input type="email" id="email" name="email" required>
\end{lstlisting}
\textbf{Beneficii.} Cititoarele anunta corect; clic-ul pe text focuseaza inputul; eticheta ramane vizibila.
\end{dado}

\subsection{Tipuri de input frecvente}

Atributul \code{type} controleaza interfata oferita de browser si validarea. Pe mobil, tipul determina tastatura afisata.

\begin{itemize}[itemsep=2pt]
  \item \code{type="text"} \textendash{} text generic.
  \item \code{type="email"} \textendash{} validare format email; tastatura cu \code{@} pe mobil.
  \item \code{type="tel"} \textendash{} tastatura numerica pe mobil; nu valideaza formatul (variaza pe tari).
  \item \code{type="number"} \textendash{} accepta doar cifre; arrows pe desktop.
  \item \code{type="date"} \textendash{} date picker nativ.
  \item \code{type="password"} \textendash{} ascunde caracterele.
  \item \code{type="url"} \textendash{} validare format URL.
  \item \code{type="search"} \textendash{} buton clear in unele browsere.
\end{itemize}

\begin{ponturi}
Browserul valideaza nativ inputurile: campurile cu atributul \code{required} blocheaza submit-ul daca sunt goale; \code{type="email"} respinge valori care nu contin \code{@}. Aceasta validare \textbf{nu inlocuieste} validarea pe server, dar imbunatateste experienta utilizatorului inainte de a trimite cererea.
\end{ponturi}

\subsection{Element HTML \textendash{} rol UI generic}

Daca ai mai folosit o interfata grafica (QtCreator, Figma, Visual Studio, Android Studio), recunosti aceste pattern-uri sub alte nume. Tabelul de mai jos pune in oglinda terminologia:

\begin{tcolorbox}[
  enhanced, breakable, sharp corners, arc=0pt,
  colback=accentSoft, colframe=partcolor, boxrule=0pt, leftrule=3pt,
  left=12pt, right=12pt, top=8pt, bottom=8pt,
  fontupper=\small,
]
\renewcommand{\arraystretch}{1.35}
\begin{tabular}{@{}p{0.42\linewidth}@{\hspace{8pt}}p{0.52\linewidth}@{}}
{\sffamily\bfseries\color{partcolor} Element HTML} & {\sffamily\bfseries\color{partcolor} Rolul UI generic} \\
\code{<label>}                       & text static (eticheta unui camp) \\
\code{<input type="text">}           & lineedit -- input pentru text scurt \\
\code{<input type="email">}          & lineedit cu validare email \\
\code{<input type="password">}       & lineedit ascuns (parola) \\
\code{<input type="number">}         & spinbox -- numar cu sageti \\
\code{<input type="date">}           & date picker \\
\code{<input type="range">}          & slider -- valoare pe interval \\
\code{<textarea>}                    & text multi-linie \\
\code{<select>} + \code{<option>}    & combobox -- o alegere dintr-o lista \\
\code{<input type="radio">}          & radio button -- o singura alegere dintr-un grup \\
\code{<input type="checkbox">}       & checkbox -- alegeri multiple (sau pornit/oprit) \\
\code{<input type="file">}           & file dialog -- selectie fisier \\
\code{<input type="color">}          & color picker \\
\code{<button>}                      & buton de actiune \\
\code{<a href="...">}                & link -- navigare catre URL \\
\code{<img>}                         & imagine statica \\
\code{<video>} / \code{<audio>}      & player media nativ \\
\code{<details>} + \code{<summary>}  & accordeon -- arata/ascunde \\
\code{<dialog>}                      & modal -- fereastra suprapusa \\
\code{<progress>}                    & bara de progres \\
\end{tabular}
\end{tcolorbox}

Regula: cand alegi un element, te intrebi mai intai ,,ce rol UI joaca?'' \textendash{} nu ,,cum ar arata cu CSS''. Tag-ul corect rezolva accesibilitatea, validarea \textcommabelow{s}i comportamentul nativ \emph{gratis}.

% =========================================================
\section{Atribute esentiale}

\begin{itemize}
  \item \code{href} pe \code{<a>} \textendash{} URL-ul tinta.
  \item \code{src} pe \code{<img>} \textendash{} calea catre imagine.
  \item \code{alt} pe \code{<img>} \textendash{} text descriptiv. Obligatoriu conform WCAG 2.1 (criteriul 1.1.1 \emph{Non-text Content}).
  \item \code{type} pe \code{<input>} \textendash{} determina validarea si interfata.
  \item \code{required} pe \code{<input>} \textendash{} activeaza validarea nativa.
  \item \code{lang} pe \code{<html>} \textendash{} limba documentului.
  \item \code{title} pe orice element \textendash{} tooltip la hover (uz limitat, nu inlocuieste \code{aria-label}).
\end{itemize}

\subsection{Calitatea textului \texttt{alt}}

Textul \code{alt} descrie ce \emph{vede} imaginea, in contextul paginii. Scopul: utilizatorul fara acces vizual primeste informatia echivalenta.

\begin{dano}[Asa nu \textendash{} alt nesemnificativ sau spam SEO]
\begin{lstlisting}[style=html, numbers=none]
<img src="echipa.jpg" alt="">
<img src="echipa.jpg" alt="poza">
<img src="echipa.jpg" alt="best best best iasi training web html css">
\end{lstlisting}
\end{dano}

\begin{dado}[Asa da \textendash{} alt descriptiv si specific]
\begin{lstlisting}[style=html, numbers=none]
<img src="echipa.jpg" alt="Echipa BEST Iasi in fata sediului UAIC, primavara 2026">
\end{lstlisting}
\textbf{Exceptie:} pentru imagini pur decorative (fundal, separator, ornament), \code{alt=""} este corect \textendash{} cititorul de ecran le ignora.
\end{dado}

% =========================================================
\section{Tag-uri suplimentare \textendash{} men\textcommabelow{t}iuni rapide}

Pentru proiecte mai complexe, HTML5 ofera elemente specializate care nu intra in scopul fazei curente, dar merita cunoscute la nivel de existen\textcommabelow{t}a:

\begin{itemize}[itemsep=2pt]
  \item \code{<table>} cu \code{<thead>}, \code{<tbody>}, \code{<tfoot>}, \code{<caption>} \textendash{} pentru date tabulare reale (nu pentru layout). Atributele \code{scope}, \code{headers} asigura accesibilitatea.
  \item \code{<video>} si \code{<audio>} \textendash{} mediafile native, fara nevoie de plugin-uri (Flash, Silverlight au murit). Atribute: \code{controls}, \code{autoplay}, \code{muted}, \code{loop}, \code{poster} (pentru video).
  \item \code{<canvas>} \textendash{} suprafa\textcommabelow{t}a 2D pentru desenare programatica cu JavaScript. Folosit pentru jocuri, vizualizari, anima\textcommabelow{t}ii complexe.
  \item \code{<details>} si \code{<summary>} \textendash{} accordeon nativ, fara JavaScript.
  \item \code{<dialog>} \textendash{} modal nativ, suportat in toate browserele moderne.
  \item \code{<picture>} \textendash{} container pentru \code{<img>} cu surse alternative pe formate sau dimensiuni.
\end{itemize}

Documenta\textcommabelow{t}ia exhaustiva: MDN \textendash{} HTML elements reference (link in resursele de la finalul capitolului).

% =========================================================
\section{Comentarii si caractere speciale}

Comentariile HTML nu apar in pagina dar raman vizibile in \emph{View Source}. Sintaxa: \code{<!-- text -->}.

Caracterele rezervate (\code{<}, \code{>}, \code{\&}) trebuie scrise prin entitati pentru a nu fi interpretate ca markup:

\begin{itemize}[itemsep=1pt]
  \item \code{\&lt;} \textendash{} \code{<}
  \item \code{\&gt;} \textendash{} \code{>}
  \item \code{\&amp;} \textendash{} \code{\&}
  \item \code{\&copy;} \textendash{} \code{(c)}
  \item \code{\&nbsp;} \textendash{} spatiu nedespartibil
\end{itemize}

% =========================================================
\begin{bonusbox}
\textbf{Acest box e optional.} Daca esti la inceput, treci direct la \emph{Verificare}. Pattern-urile de mai jos preg\u atesc terenul pentru Lighthouse \textcommabelow{s}i accesibilitate avansata.

\medskip

\textbf{\faImage\ \texttt{<picture>} \textendash{} imagini responsive cu fallback}

\code{<img>} obi\textcommabelow{s}nuit serve\textcommabelow{s}te o singura sursa. \code{<picture>} alege intre mai multe surse: format modern (AVIF $\to$ WebP $\to$ JPEG fallback) sau dimensiune (mobil vs desktop):

\begin{lstlisting}[style=html, numbers=none]
<picture>
  <source srcset="hero.avif" type="image/avif">
  <source srcset="hero.webp" type="image/webp">
  <img src="hero.jpg" alt="Sediul BEST Iasi" width="800" height="450">
</picture>
\end{lstlisting}

Browserul incarca \textbf{prima sursa pe care o sus\textcommabelow{t}ine}. Chrome modern $\to$ AVIF (cel mai mic), Safari vechi $\to$ JPEG. Zero JavaScript, zero detec\textcommabelow{t}ie de feature.

\medskip

\textbf{\faBolt\ \texttt{loading="lazy"} \textcommabelow{s}i \texttt{decoding="async"}}

Doua atribute care imbunata\textcommabelow{t}esc semnificativ Lighthouse fara cost:

\begin{lstlisting}[style=html, numbers=none]
<img src="profile.jpg" alt="Voluntar BEST"
     loading="lazy"           <!-- amana download pana cand intra in viewport -->
     decoding="async"         <!-- decodifica pe alt thread, nu blocheaza UI -->
     width="200" height="200">
\end{lstlisting}

\textbf{Excep\textcommabelow{t}ie:} imaginea hero (prima vizibila) \textbf{nu} primeste \code{loading="lazy"} \textendash{} ar intarzia LCP. Pentru ea: \code{fetchpriority="high"}.

\medskip

\textbf{\faUniversalAccess\ ARIA \textendash{} c\u and semantica nu e suficienta}

ARIA (\emph{Accessible Rich Internet Applications}) adauga semnificatie pentru cititoarele de ecran cand HTML-ul semantic nu acopera contextul. \textbf{Regula de aur:} prima alegere e mereu un tag semantic. ARIA il completeaza.

Cele 3 atribute pe care le folose\textcommabelow{s}ti des:

\begin{itemize}[itemsep=2pt]
  \item \code{aria-label="..."} \textendash{} eticheta accesibila cand textul vizibil lipseste sau e ambiguu. Exemplu clasic: butonul hamburger.
\begin{lstlisting}[style=html, numbers=none]
<button class="menu-toggle" aria-label="Deschide meniul">
  <svg>...</svg>
</button>
\end{lstlisting}

  \item \code{aria-current="page"} \textendash{} indica link-ul activ in navigare:
\begin{lstlisting}[style=html, numbers=none]
<nav>
  <a href="/" aria-current="page">Acasa</a>
  <a href="/projects">Proiecte</a>
</nav>
\end{lstlisting}

  \item \code{role="..."} \textendash{} \textbf{rar necesar}. Folose\textcommabelow{s}te-l doar cand transformi un \code{<div>} in ceva semantic (ce \textbf{nu ar trebui sa faci}). Excep\textcommabelow{t}ie utila: \code{role="alert"} pentru notificari live.
\end{itemize}

\begin{atentie}
NU adauga \code{role="button"} pe un \code{<button>}, sau \code{role="navigation"} pe un \code{<nav>}. Tag-ul semantic are deja rolul corect. ARIA duplicat = bruiaj pentru cititoarele de ecran.
\end{atentie}

\medskip

\textbf{\faRoute\ Unde mergi de aici}

\begin{itemize}[itemsep=1pt]
  \item \emph{MDN \texttt{<picture>} reference:} \url{https://developer.mozilla.org/en-US/docs/Web/HTML/Element/picture}
  \item \emph{web.dev Lazy load images:} \url{https://web.dev/articles/browser-level-image-lazy-loading}
  \item \emph{WAI-ARIA Authoring Practices:} \url{https://www.w3.org/WAI/ARIA/apg/}
\end{itemize}
\end{bonusbox}

% =========================================================
\section*{Verificare cuno\textcommabelow{s}tin\textcommabelow{t}e}

\begin{selfcheck}
\begin{enumerate}[itemsep=2pt]
  \item Care este diferen\textcommabelow{t}a func\textcommabelow{t}ionala intre \code{<div class="header">} \textcommabelow{s}i \code{<header>}?
  \item De ce este nevoie de \code{<meta name="viewport">} pentru afi\textcommabelow{s}area pe mobil?
  \item \emph{Spot the issue:} \code{<input type="email" placeholder="Email">}.
  \item Cate elemente \code{<h1>} recomanda conven\textcommabelow{t}ia industriala (accesibilitate + SEO) sa existe intr-o pagina?
  \item Care sunt cele doua atribute necesare cand un link extern se deschide intr-un tab nou?
\end{enumerate}
\end{selfcheck}

% =========================================================
\section*{Recapitulare}

\begin{recap}[Faza 1 -- HTML]
\begin{itemize}[itemsep=1pt]
  \item Anatomia: \code{<!DOCTYPE html>} $\rightarrow$ \code{<html lang>} $\rightarrow$ \code{<head>} (metadate) $\rightarrow$ \code{<body>} (continut).
  \item Trei meta tags obligatorii: \code{charset}, \code{viewport}, \code{<title>}.
  \item Tag-uri semantice in locul \code{<div>}: \code{<header>}, \code{<main>}, \code{<section>}, \code{<footer>}, \code{<nav>}, \code{<article>}, \code{<aside>}.
  \item Un singur \code{<h1>} per pagina; ierarhia titlurilor nu admite salturi de nivel.
  \item Forms accesibile: \code{<label for="X">} corespunde cu \code{<input id="X">}.
  \item Atribute obligatorii pentru accesibilitate: \code{alt} pe \code{<img>}, \code{lang} pe \code{<html>}.
  \item Linkuri externe in tab nou: combinatia \code{target="\_blank"} + \code{rel="noopener noreferrer"}.
  \item \faRocket\ Bonus: \code{<picture>} pentru fallback AVIF/WebP/JPEG, \code{loading="lazy"} + \code{decoding="async"} pe imagini, ARIA (\code{aria-label}, \code{aria-current}).
\end{itemize}
\end{recap}

\partdivider

% =========================================================
\begin{joc}
\textbf{Platforma:} CodePen (\url{codepen.io}). \quad \textit{URL-ul exact este distribuit de trainer in chat.}

\medskip

\textbf{Activitatea.} Documentul de pornire contine o structura non-semantica (\code{<div>} folosit excesiv), atribute \code{alt} lipsa pe imagini si elemente \code{<input>} fara \code{<label>} asociat. Sarcina este de a transforma documentul intr-o varianta semantica si conformanta cu standardele de accesibilitate.

\medskip

\textbf{Criterii de validare:}
\begin{itemize}[itemsep=1pt]
  \item Folosirea tag-urilor \code{<header>}, \code{<main>}, \code{<section>}, \code{<footer>} in locul \code{<div>}.
  \item Prezenta atributului \code{alt} pe toate elementele \code{<img>}.
  \item Asocierea \code{<label for>} \textendash{} \code{<input id>} pentru fiecare camp de formular.
  \item Existenta unui singur element \code{<h1>} per pagina.
  \item Linkuri externe au \code{target="\_blank"} + \code{rel="noopener noreferrer"}.
\end{itemize}

\medskip

\textbf{Durata estimata:} 5 minute.
\end{joc}

% =========================================================
\section*{\faGraduationCap\ \ Resurse pentru aprofundare}

\begin{itemize}[itemsep=2pt]
  \item Mozilla Developer Network. \emph{HTML elements reference}. \\\url{https://developer.mozilla.org/en-US/docs/Web/HTML/Element}
  \item The Odin Project. \emph{HTML Foundations}. \\\url{https://www.theodinproject.com/paths/foundations/courses/foundations\#html-foundations}
  \item Google. \emph{Learn HTML}. \url{https://web.dev/learn/html}
  \item HTML Reference. \emph{Vizual cheatsheet de tag-uri}. \url{https://htmlreference.io}
  \item Google. \emph{Learn Forms}. \url{https://web.dev/learn/forms}
  \item W3C. \emph{Web Accessibility Initiative \textendash{} WCAG}. \url{https://www.w3.org/WAI/standards-guidelines/wcag/}
\end{itemize}

\bigskip

\bridge{Capitolul 2 (CSS) imbraca structura HTML cu stil vizual. Aspectul, layout-ul \textcommabelow{s}i responsivitatea se construiesc plecand de la pagina semantica generata aici.}

% Fragment -- inclus de main.tex prin \input
\clearpage
\setpartlabel{themeCss}{\faPaintBrush}{Faza 2: CSS}
\setcounter{section}{0}
\phantomsection
\addcontentsline{toc}{part}{\protect\textbf{Faza 2 \textendash{} CSS}}
\handout{Faza 2 \textendash{} CSS}{Stilizare. Layout. Responsivitate.}

CSS (\emph{Cascading Style Sheets}) controleaza aspectul vizual al documentelor HTML. Numele provine de la \emph{cascada}: cand mai multe reguli vizeaza acelasi element, browserul aplica un algoritm de prioritate (cascada) pentru a alege valoarea finala.

Capitolul prezinta selectorii, modelul de cutie, unitatile de masura, variabilele CSS, sistemele de layout flexbox si grid, abordarea \emph{mobile-first}, si suportul nativ pentru tema dark/light.

\bigskip

\begin{outcomes}
Dupa capitolul asta vei \textcommabelow{s}ti:
\begin{itemize}[itemsep=1pt]
  \item Folose\textcommabelow{s}te selectorii CSS \textcommabelow{s}i in\textcommabelow{t}elege specificitatea ca triplet \code{(id, clasa, tag)}.
  \item Aplica modelul de cutie cu \code{box-sizing: border-box} pentru calcule predictibile.
  \item Define\textcommabelow{s}te variabile CSS in \code{:root} si le refolose\textcommabelow{s}te prin \code{var()}.
  \item Construie\textcommabelow{s}te layout-uri cu flexbox (6 proprieta\textcommabelow{t}i cheie).
  \item Scrie reguli mobile-first cu \code{@media (min-width: ...)}.
  \item Pozitioneaza elemente cu \code{position} \textcommabelow{s}i creeaza header sticky.
  \item Aplica anima\textcommabelow{t}ii cu \code{transition} si \code{@keyframes}, respectand \code{prefers-reduced-motion}.
  \item Combina \code{transform}, \code{transition} si \code{cubic-bezier()} pentru micro-interactiuni profesionale (hover-lift, fade-in, spinner).
  \item \faRocket\ \emph{Bonus self-study:} transformari 3D (flip-card), \code{animation-timeline} pentru scroll-driven, \code{filter}/\code{backdrop-filter}, \code{clip-path}, container queries.
\end{itemize}
\end{outcomes}

\begin{whymatters}
CSS face diferen\textcommabelow{t}a intre pagina-document \textcommabelow{s}i pagina-produs. Aceea\textcommabelow{s}i structura HTML poate parea neprofesionala sau profesionala, in func\textcommabelow{t}ie de stilizare. Mai mult: CSS-ul prost arhitecturat (cu \code{!important}, specificitate ridicata, valori absolute) genereaza ore de depanare. Conventiile prezentate aici (variabile, mobile-first, flexbox) sunt cele aplicate in practica industriala.
\end{whymatters}

\begin{nota}[Pre-rechizite]
Capitolul presupune in\textcommabelow{t}elegerea tag-urilor HTML semantice \textcommabelow{s}i a structurii \code{<head>}/\code{<body>} (Faza 1).
\end{nota}

% =========================================================
\section{Conectarea fisierului CSS la HTML}

CSS poate fi inclus in trei moduri. Conventia recomandata este externalizarea intr-un fisier separat:

\begin{lstlisting}[style=html, caption={Stilurile externe (recomandat)}]
<head>
  <link rel="stylesheet" href="style.css">
</head>
\end{lstlisting}

\begin{dano}[Asa nu \textendash{} stiluri inline imprastiate]
\begin{lstlisting}[style=html, numbers=none]
<button style="background: blue; color: white; padding: 10px;">Click</button>
\end{lstlisting}
\textbf{Probleme.} Imposibil de reutilizat; stilul nu poate fi schimbat global; specificitatea ridicata creeaza conflicte; pagina HTML devine ilizibila.
\end{dano}

\begin{dado}[Asa da \textendash{} clasa cu stil extern]
\begin{lstlisting}[style=html, numbers=none]
<button class="cta">Click</button>
\end{lstlisting}
In \code{style.css}: \code{.cta \{ background: blue; color: white; padding: 10px; \}}
\end{dado}

% =========================================================
\section{Selectori}

Selectorii determina elementele asupra carora se aplica un set de reguli. Categoriile principale:

\begin{lstlisting}[style=css, caption={Selectori uzuali}]
/* Selector de tip: vizeaza toate elementele <h1> */
h1 { color: navy; }

/* Selector de clasa: reutilizabil, prefixat cu . */
.cta { background: orange; }

/* Selector de id: unic pe pagina, prefixat cu # */
#contact { padding: 2rem; }

/* Selector de atribut */
input[type="email"] { border-color: blue; }

/* Combinator descendent: <a> oriunde sub <nav> */
nav a { text-decoration: none; }

/* Combinator copil direct: <li> direct sub <ul> */
ul > li { margin-bottom: 0.5rem; }

/* Combinator frate adiacent: <p> imediat dupa <h2> */
h2 + p { margin-top: 0; }

/* Pseudo-clase: stari ale elementului */
.cta:hover { background: red; }
.cta:focus-visible { outline: 2px solid blue; }
.cta:disabled { opacity: 0.5; }

/* Pseudo-elemente: parti generate */
p::first-letter { font-size: 2em; }
\end{lstlisting}

\subsection{Specificitate \textendash{} cum se rezolva conflictele}

Cand mai multe reguli vizeaza acelasi element, browserul calculeaza o specificitate pentru fiecare si o aplica pe cea cu valoare maxima. Calculul foloseste un triplet \code{(a, b, c)}:

\begin{itemize}[itemsep=1pt]
  \item \code{a} \textendash{} numarul de id-uri.
  \item \code{b} \textendash{} numarul de clase, atribute, pseudo-clase.
  \item \code{c} \textendash{} numarul de tag-uri si pseudo-elemente.
\end{itemize}

Comparatia se face lexicografic: \code{(0,1,0)} pierde fata de \code{(1,0,0)} indiferent de \code{b} si \code{c}. Stilurile inline au specificitate \code{(1,0,0,0)}, peste orice selector. Cuvantul cheie \code{!important} suprascrie tot, dar este considerat un \emph{antipattern} \textendash{} indica faptul ca arhitectura CSS este compromisa.

\begin{dano}[Asa nu \textendash{} folosirea \texttt{!important} ca solutie]
\code{.btn \{ background: red !important; \}}\\
\textbf{Problema.} Cand alta regula are nevoie sa schimbe culoarea, va trebui si ea sa foloseasca \code{!important}, escaladand conflictul.
\end{dano}

\begin{dado}[Asa da \textendash{} specificitate proportionala]
\code{.btn \{ background: blue; \} .btn.btn-danger \{ background: red; \}}\\
\textbf{Solutie:} clase modificator (\code{.btn-danger}) ridica specificitatea natural cu \code{(0,2,0)} fata de \code{(0,1,0)}.
\end{dado}

% =========================================================
\section{Modelul de cutie}

Fiecare element redat de browser este o cutie compusa din patru straturi concentrice: \emph{content}, \emph{padding}, \emph{border}, \emph{margin}.

\boxmodelfigure

\begin{itemize}
  \item \textbf{content} \textendash{} zona ocupata de continutul efectiv (text, imagine).
  \item \textbf{padding} \textendash{} spatiul dintre continut si chenar.
  \item \textbf{border} \textendash{} chenarul.
  \item \textbf{margin} \textendash{} spatiul exterior, intre cutia curenta si elementele adiacente.
\end{itemize}

\subsection{Implicit vs \texttt{border-box}}

Comportamentul implicit (\code{box-sizing: content-box}) include in \code{width} doar continutul. Adaugarea de \code{padding} sau \code{border} mareste latimea totala. Acest comportament face calculele predispuse la erori.

\begin{dano}[Asa nu \textendash{} content-box, latimi imprevizibile]
\begin{lstlisting}[style=css, numbers=none]
.card {
  width: 300px;
  padding: 20px;
  border: 1px solid #ccc;
  /* Latime totala: 300 + 20*2 + 1*2 = 342px (oversize) */
}
\end{lstlisting}
\end{dano}

\begin{dado}[Asa da \textendash{} border-box global, latimi predictibile]
\begin{lstlisting}[style=css, numbers=none]
*, *::before, *::after { box-sizing: border-box; }

.card {
  width: 300px;
  padding: 20px;
  border: 1px solid #ccc;
  /* Latime totala: 300px (padding si border absorbite in width) */
}
\end{lstlisting}
\end{dado}

% =========================================================
\section{Unitati de masura}

CSS suporta zeci de unitati. In practica, doar 5\textendash{}6 acopera majoritatea cazurilor:

\begin{itemize}
  \item \code{px} \textendash{} pixel absolut. Pentru border-uri si dimensiuni de detaliu.
  \item \code{rem} \textendash{} relativ la font-size-ul \code{<html>} (implicit 16px). Pentru font-size si spatiere uniforma.
  \item \code{em} \textendash{} relativ la font-size-ul elementului parinte. Util pentru padding-uri scaland cu textul.
  \item \code{\%} \textendash{} relativ la dimensiunea parintelui pe acea axa.
  \item \code{vw} / \code{vh} \textendash{} 1\% din latimea / inaltimea viewport-ului.
  \item \code{ch} \textendash{} latimea unui caracter ,,0''. Util pentru \code{max-width} pe paragrafe.
\end{itemize}

\begin{ponturi}
Pentru font-size si spatiere generala, \code{rem} este alegerea recomandata. Cand utilizatorul mareste font-size-ul implicit din browser (Settings $\to$ Appearance), tot site-ul scaleaza proportional. \code{px} pe font-size ignora aceasta preferinta.
\end{ponturi}

\begin{lstlisting}[style=css, caption={Compunere uzuala}]
:root {
  font-size: 100%;          /* echivalent 16px in browserele standard */
}

body { font-size: 1rem; }   /* 16px */
h1   { font-size: 2.5rem; } /* 40px */
.container { max-width: 64rem; }    /* 1024px */

p { max-width: 65ch; }      /* lungime optima de citire */

.hero {
  min-height: 80vh;         /* 80% din inaltimea ferestrei */
  padding: 2rem;
}
\end{lstlisting}

% =========================================================
\section{Variabile CSS (custom properties)}

Variabilele CSS permit definirea valorilor o singura data si reutilizarea lor pe parcursul intregului document. Modificarea unei singure declaratii in \code{:root} se propaga automat in toate locurile in care variabila este folosita.

\begin{lstlisting}[style=css]
:root {
  /* Paleta */
  --primary: #0673BA;
  --accent:  #FAA51E;
  --bg:      #FFFCF6;
  --text:    #1A1A1A;

  /* Spatiere consistenta */
  --space-1: 0.5rem;
  --space-2: 1rem;
  --space-3: 2rem;

  /* Layout */
  --max-width: 64rem;
  --radius: 0.5rem;
}

body {
  background: var(--bg);
  color: var(--text);
}

h1 { color: var(--primary); }
.cta { background: var(--accent); }
\end{lstlisting}

Variabilele CSS sunt \emph{scoped}: redefinirea intr-un selector descendent suprascrie valoarea pentru subarbore. Acest comportament permite teme contextuale fara duplicarea regulilor:

\begin{lstlisting}[style=css, caption={Tema schimbata local}]
.card-dark {
  --bg:   #1F2937;
  --text: #F1F5F9;
  /* toate copiii care folosesc var(--bg) si var(--text) preiau noile valori */
}
\end{lstlisting}

% =========================================================
\section{Proprietatea \texttt{display} \textendash{} valori}

\code{display} controleaza modul fundamental in care elementul se aseaza in pagina. Cele 6 valori pe care le intalnesti in 99\% din cazuri:

\begin{tcolorbox}[
  enhanced, breakable, sharp corners, arc=0pt,
  colback=accentSoft, colframe=partcolor, boxrule=0pt, leftrule=3pt,
  left=12pt, right=12pt, top=8pt, bottom=8pt, fontupper=\small,
]
\renewcommand{\arraystretch}{1.4}
\begin{tabular}{@{}p{0.22\linewidth}@{\hspace{8pt}}p{0.74\linewidth}@{}}
{\sffamily\bfseries\color{partcolor} Valoare} & {\sffamily\bfseries\color{partcolor} Efect} \\
\code{block}        & Element pe linie noua, ocupa toata latimea disponibila. Default pentru \code{<div>}, \code{<p>}, \code{<h1>...<h6>}, \code{<section>}. \\
\code{inline}       & Curge in textul din jur, ignora \code{width}/\code{height}. Default pentru \code{<span>}, \code{<a>}, \code{<strong>}, \code{<em>}. \\
\code{inline-block} & Curge in text ca inline, dar accepta \code{width}/\code{height} ca un block. Util pentru butoane in text. \\
\code{flex}         & Activeaza modelul Flexbox (uni-dimensional) pe \emph{copii}i elementului. Cel mai folosit pentru header-uri, navigatie, alinieri. \\
\code{grid}         & Activeaza modelul CSS Grid (bi-dimensional) pe copii. Pentru galerii, dashboard-uri, layout-uri complexe. \\
\code{none}         & Ascunde complet elementul (nu ocupa spatiu). Util pentru meniuri mobile inchise, modale ascunse. \\
\end{tabular}
\end{tcolorbox}

\begin{atentie}
\code{display: none} \textbf{ascunde complet} elementul (cititoarele de ecran il sar). Daca vrei doar sa-l faci invizibil dar accesibil, foloseste \code{visibility: hidden} sau \code{opacity: 0; pointer-events: none}.
\end{atentie}

% =========================================================
\section{Flexbox}

Flexbox (\emph{Flexible Box Layout}) este un model de layout uni-dimensional, util pentru distribuirea spatiului si alinierea continutului pe o axa (orizontala sau verticala).

\begin{lstlisting}[style=css, caption={Header cu logo si navigare aliniate}]
header {
  display: flex;                   /* activeaza flexbox */
  flex-direction: row;             /* axa principala = orizontala (default) */
  justify-content: space-between;  /* distribuire pe axa principala */
  align-items: center;             /* aliniere pe axa secundara */
  flex-wrap: wrap;                 /* trecere pe randul urmator daca nu incape */
  gap: var(--space-2);             /* spatiu intre copii */
}
\end{lstlisting}

\subsection{Cele 6 proprietati pe care le folosesti}

\begin{itemize}[itemsep=2pt]
  \item \code{display: flex} \textendash{} activeaza modelul (pe parinte).
  \item \code{flex-direction} \textendash{} axa principala: \code{row} (default, orizontal), \code{column} (vertical), \code{row-reverse}, \code{column-reverse}.
  \item \code{justify-content} \textendash{} distributie pe axa principala. Vezi tabelul de mai jos.
  \item \code{align-items} \textendash{} aliniere pe axa secundara. Vezi tabelul de mai jos.
  \item \code{gap} \textendash{} spatiu uniform intre copii (in unitati: \code{1rem}, \code{16px}). Inlocuieste \code{margin}.
  \item \code{flex-wrap} \textendash{} \code{nowrap} (default, totul pe o linie), \code{wrap} (trece pe linia urmatoare).
  \item \code{flex: 1} (pe copil) \textendash{} cre\textcommabelow{s}te pentru a umple spatiul disponibil.
\end{itemize}

\subsubsection{Valorile pentru \texttt{justify-content} (axa principala)}

\begin{tcolorbox}[
  enhanced, breakable, sharp corners, arc=0pt,
  colback=accentSoft, colframe=partcolor, boxrule=0pt, leftrule=3pt,
  left=12pt, right=12pt, top=8pt, bottom=8pt, fontupper=\small,
]
\renewcommand{\arraystretch}{1.35}
\begin{tabular}{@{}p{0.30\linewidth}@{\hspace{8pt}}p{0.66\linewidth}@{}}
\code{flex-start} (default) & Toate copii lipiti la inceputul axei. \\
\code{center}               & Grupate la mijlocul axei. \\
\code{flex-end}             & Lipiti la finalul axei. \\
\code{space-between}        & Primul/ultimul la margini, spa\textcommabelow{t}iu egal intre restul. \\
\code{space-around}         & Spa\textcommabelow{t}iu egal in jurul fiecaruia (margine = jumatate). \\
\code{space-evenly}         & Spa\textcommabelow{t}iu absolut egal peste tot (inclusiv la margini). \\
\end{tabular}
\end{tcolorbox}

\subsubsection{Valorile pentru \texttt{align-items} (axa secundara)}

\begin{tcolorbox}[
  enhanced, breakable, sharp corners, arc=0pt,
  colback=accentSoft, colframe=partcolor, boxrule=0pt, leftrule=3pt,
  left=12pt, right=12pt, top=8pt, bottom=8pt, fontupper=\small,
]
\renewcommand{\arraystretch}{1.35}
\begin{tabular}{@{}p{0.30\linewidth}@{\hspace{8pt}}p{0.66\linewidth}@{}}
\code{stretch} (default) & Copii se intind pe toata latimea axei secundare. \\
\code{center}            & Centrate pe axa secundara. \\
\code{flex-start}        & Lipiti la inceputul axei (sus pentru row, stanga pentru column). \\
\code{flex-end}          & Lipiti la final (jos / dreapta). \\
\code{baseline}          & Aliniati dupa linia de baza a textului (util pentru fonturi de marimi diferite). \\
\end{tabular}
\end{tcolorbox}

\begin{ponturi}
\textbf{Trick uzual:} \code{display: flex; justify-content: center; align-items: center;} centreaza orice element in containerul sau, atat orizontal cat si vertical, in 3 linii.
\end{ponturi}

% =========================================================
\section{CSS Grid (introducere scurta)}

Grid este sistemul de layout bidimensional al CSS-ului. Spre deosebire de flexbox (uni-dimensional), grid permite controlul simultan al randurilor si coloanelor.

\begin{lstlisting}[style=css, caption={Galerie cu coloane responsive automat}]
.gallery {
  display: grid;
  grid-template-columns: repeat(auto-fit, minmax(240px, 1fr));
  gap: var(--space-2);
}
\end{lstlisting}

Mecanism: \code{auto-fit} creeaza cate coloane incap; \code{minmax(240px, 1fr)} cere o latime minima de 240px si o expansiune egala (\code{1fr}) cand ramane spatiu. Rezultatul: galerie responsive fara media queries.

% =========================================================
\section{Mobile-first}

Strategia \emph{mobile-first} consta in scrierea regulilor pentru ecrane mici ca valori implicite, urmata de extinderea acestora pentru ecrane mai mari prin \emph{media queries}:

\begin{lstlisting}[style=css]
/* Stiluri implicite, pentru mobil */
.nav { display: none; }
.nav-toggle { display: flex; }

/* Adaugiri pentru ecrane >= 768px */
@media (min-width: 768px) {
  .nav { display: flex; }
  .nav-toggle { display: none; }
}
\end{lstlisting}

\begin{dano}[Asa nu \textendash{} desktop-first cu \texttt{max-width}]
\begin{lstlisting}[style=css, numbers=none]
.nav { display: flex; }
@media (max-width: 767px) {
  .nav { display: none; }
  .nav-toggle { display: flex; }
}
\end{lstlisting}
\textbf{Problema.} Mobilul descarca si proceseaza CSS conceput pentru desktop, apoi il anuleaza. Pierdere de performanta pe contextul cu cele mai limitate resurse.
\end{dano}

\begin{dado}[Asa da \textendash{} mobile-first cu \texttt{min-width}]
\begin{lstlisting}[style=css, numbers=none]
.nav { display: none; }
.nav-toggle { display: flex; }
@media (min-width: 768px) {
  .nav { display: flex; }
  .nav-toggle { display: none; }
}
\end{lstlisting}
\textbf{Beneficiu.} Mobilul (>60\% din trafic) primeste varianta optimizata implicit; desktopul adauga incremental.
\end{dado}

% =========================================================
\section{Suport pentru tema sistem (light / dark)}

Functia \code{light-dark()} si proprietatea \code{color-scheme} permit afisarea conditionata a culorilor in functie de preferinta sistemului de operare al utilizatorului:

\begin{lstlisting}[style=css]
html { color-scheme: light dark; }

:root {
  --bg:   light-dark(#FFFCF6, #0F172A);
  --text: light-dark(#1A1A1A, #F1F5F9);
}
\end{lstlisting}

Browserul selecteaza prima sau a doua valoare in functie de tema activa. Suportul nativ in browserele moderne este disponibil incepand cu 2024 (Chrome 123+, Firefox 120+, Safari 17.5+).

% =========================================================
\section{Pozi\textcommabelow{t}ionarea elementelor}

Proprietatea \code{position} controleaza modul in care un element se aseaza in fluxul paginii. Cinci valori, cu reper de cand le folose\textcommabelow{s}ti:

\begin{tcolorbox}[
  enhanced, breakable, sharp corners, arc=0pt,
  colback=accentSoft, colframe=partcolor, boxrule=0pt, leftrule=3pt,
  left=12pt, right=12pt, top=8pt, bottom=8pt, fontupper=\small,
]
\renewcommand{\arraystretch}{1.4}
\begin{tabular}{@{}p{0.20\linewidth}@{\hspace{6pt}}p{0.50\linewidth}@{\hspace{6pt}}p{0.24\linewidth}@{}}
{\sffamily\bfseries\color{partcolor} Valoare} & {\sffamily\bfseries\color{partcolor} Comportament} & {\sffamily\bfseries\color{partcolor} Cand foloseste} \\
\code{static} (default) & Respecta fluxul natural; \code{top}/\code{left} ignorate. & implicit, nu il setezi \\
\code{relative}         & Stay-in-place pana il mu\textcommabelow{t}i cu \code{top}/\code{left}; pastreaza spa\textcommabelow{t}iul. & ancora pentru \code{absolute} \\
\code{absolute}         & Scos din flux; pozi\textcommabelow{t}ionat fa\textcommabelow{t}a de primul parinte \code{relative}. & tooltip-uri, badge-uri \\
\code{fixed}            & Fixat pe viewport; ramane la scroll. & modal, buton ,,scroll-to-top'' \\
\code{sticky}           & \code{relative} pana ajunge la o margine, apoi \code{fixed}. & header sticky, sidebar \\
\end{tabular}
\end{tcolorbox}

\begin{atentie}
\code{absolute} are nevoie de un parinte \code{relative} (sau cade pana la viewport). Daca elementul ,,sare in lume'', verifica daca containerul are \code{position: relative}.
\end{atentie}

\begin{lstlisting}[style=css, caption={Header sticky modern}]
header {
  position: sticky;
  top: 0;
  background: var(--bg);
  z-index: 10;
}
\end{lstlisting}

% =========================================================
\section{Cheatsheet \textendash{} alte proprieta\textcommabelow{t}i cu enum-uri}

Cele patru cele mai uzuale dupa cele de mai sus. Le intalne\textcommabelow{s}ti zilnic.

\begin{tcolorbox}[
  enhanced, breakable, sharp corners, arc=0pt,
  colback=accentSoft, colframe=partcolor, boxrule=0pt, leftrule=3pt,
  left=12pt, right=12pt, top=8pt, bottom=8pt, fontupper=\small,
]
\renewcommand{\arraystretch}{1.35}
\begin{tabular}{@{}p{0.22\linewidth}@{\hspace{6pt}}p{0.34\linewidth}@{\hspace{6pt}}p{0.40\linewidth}@{}}
{\sffamily\bfseries\color{partcolor} Proprietate} & {\sffamily\bfseries\color{partcolor} Valori} & {\sffamily\bfseries\color{partcolor} Cand foloseste} \\
\code{box-sizing}     & \code{content-box} (default), \code{border-box} & \emph{Mereu border-box} \textendash{} aplicat global. \\
\code{overflow}       & \code{visible} (default), \code{hidden}, \code{scroll}, \code{auto} & \code{hidden} pentru a tunde, \code{auto} pentru scroll doar la nevoie. \\
\code{text-align}     & \code{left} (default), \code{center}, \code{right}, \code{justify} & Aliniere text in container. \\
\code{cursor}         & \code{auto}, \code{pointer}, \code{default}, \code{text}, \code{grab}, \code{not-allowed}, \code{wait} & \code{pointer} pe elemente clicabile non-link. \\
\code{font-weight}    & \code{normal} (400), \code{bold} (700), 100\textendash{}900 in pa\textcommabelow{s}i de 100 & \code{500} pentru \emph{medium}, \code{600} pentru \emph{semibold}. \\
\code{text-transform} & \code{none} (default), \code{uppercase}, \code{lowercase}, \code{capitalize} & In CSS, nu in HTML \textendash{} pastrezi semantica. \\
\code{visibility}     & \code{visible} (default), \code{hidden} & Ascunde dar pastreaza spa\textcommabelow{t}iul; \code{display: none} il scoate complet. \\
\code{pointer-events} & \code{auto} (default), \code{none} & \code{none} dezactiveaza click \textcommabelow{s}i hover (overlay-uri inactive). \\
\end{tabular}
\end{tcolorbox}

\begin{ponturi}
\textbf{Regula generala:} valori unitare ca \code{cursor: pointer}, \code{text-align: center}, \code{overflow: hidden} arata exact ca in tooltip-ul autocomplete al VS Code. \code{Ctrl+Space} dupa numele proprieta\textcommabelow{t}ii afi\textcommabelow{s}eaza lista completa.
\end{ponturi}

% =========================================================
\section{Anima\textcommabelow{t}ii CSS}

Animatiile imbunatatesc perceptia calitatii produsului \textcommabelow{s}i orienteaza aten\textcommabelow{t}ia utilizatorului. CSS suporta nativ doua mecanisme: \emph{transitions} pentru schimbari intre stari, si \emph{keyframes} pentru animatii arbitrare.

\subsection{Transitions \textendash{} schimbari intre stari}

\code{transition} interpoleaza fluent o proprietate cand valoarea sa se schimba (de exemplu, la \code{:hover}).

\begin{lstlisting}[style=css, caption={Buton cu hover smooth}]
.cta {
  background: var(--primary);
  color: white;
  padding: var(--space-1) var(--space-3);
  border-radius: var(--radius);
  transition: background 0.2s ease, transform 0.2s ease;
}

.cta:hover {
  background: var(--accent);
  transform: translateY(-2px);
}
\end{lstlisting}

Sintaxa: \code{transition: <proprietate> <durata> <timing-function> <delay>}.

\subsubsection{Timing function \textendash{} cum se simte miscarea}

Durata spune \emph{cat} dureaza animatia; \emph{timing function} spune \emph{cum} se distribuie miscarea in timp. Valorile predefinite sunt suficiente pentru 90\% din cazuri:

\begin{itemize}[itemsep=1pt]
  \item \code{linear} \textendash{} viteza constanta. Util pentru spinner-e si bare de progres.
  \item \code{ease} (default) \textendash{} accelereaza la inceput, decelereaza la final.
  \item \code{ease-in} \textendash{} porneste lent, se incheie rapid. Folosit pentru iesiri (\emph{exit}).
  \item \code{ease-out} \textendash{} porneste rapid, se incheie lent. Folosit pentru intrari (\emph{enter}); senzatie de aterizare lina.
  \item \code{ease-in-out} \textendash{} simetric. Folosit pentru tranzitii reversibile (toggle).
\end{itemize}

\begin{ponturi}[Regula UX rapida]
Pentru micro-interactiuni (hover, click): \code{ease-out 150ms\textendash{}250ms}. Sub 150ms se simte abrupt; peste 300ms se simte lent.
\end{ponturi}

\subsubsection{Cubic-bezier \textendash{} curba personalizata}

Cele 5 valori predefinite sunt scurtaturi pentru curbe Bezier de gradul 3. Cand vrei o curba personalizata (ex. \emph{bounce}, \emph{overshoot}), o defines\textcommabelow{t}i prin 4 numere care reprezinta cele doua puncte de control:

\begin{lstlisting}[style=css, caption={Easing personalizat cu overshoot subtil}]
.cta {
  transition: transform 0.3s cubic-bezier(0.34, 1.56, 0.64, 1);
}
.cta:hover { transform: scale(1.05); }
\end{lstlisting}

\begin{ponturi}
\textbf{Tool recomandat:} \url{https://cubic-bezier.com} \textendash{} editor vizual cu preview. Copiezi cele 4 valori si le lipesti in CSS.
\end{ponturi}

\subsection{Transform \textendash{} scalare, rota\textcommabelow{t}ie, transla\textcommabelow{t}ie}

\code{transform} aplica o transformare geometrica fara a afecta layout-ul (alte elemente nu se mut\u a). Operatiile uzuale:

\begin{itemize}
  \item \code{translate(x, y)} / \code{translateX} / \code{translateY} \textendash{} mutare pe axe (px, \%, rem). Mai performant decat \code{top}/\code{left}.
  \item \code{scale(n)} \textendash{} marire/micsorare proportionala. \code{scaleX}/\code{scaleY} pe o singura axa.
  \item \code{rotate(deg)} \textendash{} rotire in plan; pozitiv = sensul acelor de ceas.
  \item \code{skew(deg)} \textendash{} forfecare (rar folosit, util pentru efecte oblique).
  \item Combinare \emph{ordonata}: \code{transform: translateY(-2px) scale(1.05) rotate(2deg)} \textendash{} ordinea conteaza, se aplica de la dreapta la stanga.
\end{itemize}

\subsubsection{Transform-origin \textendash{} punctul de pivot}

Implicit, transformarile se aplica fata de centrul elementului. \code{transform-origin} schimba acest pivot:

\begin{lstlisting}[style=css, caption={Card care se ,,deschide'' din coltul stanga-sus}]
.card {
  transform-origin: top left;
  transition: transform 0.3s ease-out;
}
.card:hover { transform: scale(1.05); }
\end{lstlisting}

Valori uzuale: \code{center} (default), \code{top}, \code{bottom}, \code{left}, \code{right}, combinatii (\code{top left}, \code{50\% 0\%}), sau coordonate explicite in px/\%.

\begin{ponturi}[Performance]
Browserul accelereaza pe GPU doar \code{transform} si \code{opacity}. Anima\textcommabelow{t}ia altor proprieta\textcommabelow{t}i (ex. \code{width}, \code{margin}, \code{top}) declan\textcommabelow{s}eaza reflow \textcommabelow{s}i scade FPS. Regula: pentru mi\textcommabelow{s}care, folose\textcommabelow{s}te \code{transform: translate(...)}, nu \code{top}/\code{left}.
\end{ponturi}

\subsection{Keyframes \textendash{} animatii arbitrare}

Pentru anima\textcommabelow{t}ii care nu sunt declan\textcommabelow{s}ate de o stare (loading spinner, fade-in, parallax), \code{@keyframes} defineste o secven\textcommabelow{t}a de stari:

\begin{lstlisting}[style=css, caption={Spinner infinit}]
@keyframes rotate {
  from { transform: rotate(0deg); }
  to   { transform: rotate(360deg); }
}

.spinner {
  width: 32px;
  height: 32px;
  border: 4px solid #E5E7EB;
  border-top-color: var(--primary);
  border-radius: 50%;
  animation: rotate 1s linear infinite;
}
\end{lstlisting}

\begin{lstlisting}[style=css, caption={Fade-in pentru titlu hero}]
@keyframes fadeInUp {
  from { opacity: 0; transform: translateY(20px); }
  to   { opacity: 1; transform: translateY(0); }
}

.hero h1 {
  animation: fadeInUp 0.6s ease-out;
}
\end{lstlisting}

Sintaxa shorthand pentru animation: \code{animation: <nume> <durata> <timing> <delay> <iteratii> <direction> <fill-mode>}.

\subsubsection{Keyframes multi-step}

\code{@keyframes} accepta procente intermediare, nu doar \code{from} si \code{to}. Astfel po\textcommabelow{t}i descrie traiectorii complexe intr-o singura anima\textcommabelow{t}ie:

\begin{lstlisting}[style=css, caption={Buton care ,,pulseaz\u a'' atragator}]
@keyframes pulse {
  0%   { transform: scale(1);    box-shadow: 0 0 0 0 rgba(250, 165, 30, 0.7); }
  50%  { transform: scale(1.05); box-shadow: 0 0 0 12px rgba(250, 165, 30, 0); }
  100% { transform: scale(1);    box-shadow: 0 0 0 0 rgba(250, 165, 30, 0); }
}

.cta-featured { animation: pulse 2s ease-in-out infinite; }
\end{lstlisting}

\subsubsection{Animation-fill-mode \textendash{} ce ramane dupa anima\textcommabelow{t}ie}

Implicit, dupa ce animatia se termina, elementul revine la starea CSS originala. \code{animation-fill-mode} controleaza acest comportament:

\begin{itemize}[itemsep=1pt]
  \item \code{none} (default) \textendash{} elementul revine la stilul original.
  \item \code{forwards} \textendash{} pastreaza valorile din ultimul keyframe. \textbf{Cel mai folosit} pentru fade-in / slide-in (vrei sa ramana la \code{opacity: 1}).
  \item \code{backwards} \textendash{} aplica primul keyframe inainte de start (in timpul \code{animation-delay}).
  \item \code{both} \textendash{} combinatie \code{forwards + backwards}.
\end{itemize}

\begin{lstlisting}[style=css, caption={Fade-in care r\u amane vizibil}]
.hero h1 {
  opacity: 0;
  animation: fadeInUp 0.6s ease-out forwards;
}
\end{lstlisting}

\subsubsection{Animation-play-state \textendash{} pauz\u a \textcommabelow{s}i reluare}

\code{animation-play-state} permite punerea pe pauza a unei animatii \textendash{} util pentru carusele care se opresc la hover:

\begin{lstlisting}[style=css, numbers=none]
.carousel { animation: scroll 30s linear infinite; }
.carousel:hover { animation-play-state: paused; }
\end{lstlisting}

\subsection{Pattern complet: card cu hover-lift}

Combinatia transition + transform + box-shadow produce efectul de ,,ridicare'' pe care il vezi pe orice landing page modern:

\begin{lstlisting}[style=css, caption={Hover-lift, pattern de baz\u a pentru carduri}]
.card {
  background: var(--bg);
  padding: var(--space-3);
  border-radius: var(--radius);
  box-shadow: 0 1px 3px rgba(0, 0, 0, 0.08);
  transition: transform 0.25s ease-out, box-shadow 0.25s ease-out;
  cursor: pointer;
}

.card:hover {
  transform: translateY(-4px);
  box-shadow: 0 12px 24px rgba(0, 0, 0, 0.12);
}
\end{lstlisting}

Detalii care fac diferenta:
\begin{itemize}[itemsep=1pt]
  \item \code{ease-out} pe enter \textendash{} senzatie de aterizare lina.
  \item Umbra cre\textcommabelow{s}te si se ,,rasfira'' (offset Y mai mare, blur mai mare) \textendash{} simuleaza distan\textcommabelow{t}a fa\textcommabelow{t}a de fundal.
  \item Tranzi\textcommabelow{t}ie pe ambele proprieta\textcommabelow{t}i simultan, nu doar \code{transform}.
  \item Cursor \code{pointer} \textendash{} afi\textcommabelow{s}eaza inten\textcommabelow{t}ia ca elementul e clicabil.
\end{itemize}

\subsection{Accesibilitate: \texttt{prefers-reduced-motion}}

Unii utilizatori (afec\textcommabelow{t}iuni vestibulare, dependente de aten\textcommabelow{t}ie) au setat in OS preferin\textcommabelow{t}a pentru anima\textcommabelow{t}ii reduse. CSS detecteaza aceasta cerere:

\begin{lstlisting}[style=css]
@media (prefers-reduced-motion: reduce) {
  *, *::before, *::after {
    animation-duration: 0.01ms !important;
    transition-duration: 0.01ms !important;
  }
}
\end{lstlisting}

Conformitatea este o cerin\textcommabelow{t}a WCAG 2.1 (criteriul 2.3.3) \textcommabelow{s}i o practica recomandata pentru orice site cu animatii.

% =========================================================
\begin{bonusbox}
\textbf{Acest box e optional.} Daca esti la inceput, treci direct la \emph{Verificare cuno\textcommabelow{s}tin\textcommabelow{t}e}. Sec\textcommabelow{t}iunile de mai jos sunt pentru cand vrei sa ridici stacheta.

\medskip

\textbf{\faCube\ Transform\u ari 3D \textendash{} flip-card}

CSS suporta nativ transformari in spatiu, nu doar in plan. Trei piese:

\begin{itemize}[itemsep=1pt]
  \item \code{perspective} (pe parinte) \textendash{} ,,distan\textcommabelow{t}a'' fata de plan. Valori uzuale 600\textendash{}1500px.
  \item \code{transform-style: preserve-3d} (pe parinte) \textendash{} permite copiilor sa existe in spatiu 3D.
  \item \code{backface-visibility: hidden} (pe fete) \textendash{} ascunde fa\textcommabelow{t}a din spate.
\end{itemize}

\begin{lstlisting}[style=css, caption={Card care se intoarce la hover (flip-card)}]
.flip-card { perspective: 1000px; }
.flip-card-inner {
  position: relative;
  width: 100%; height: 100%;
  transition: transform 0.6s;
  transform-style: preserve-3d;
}
.flip-card:hover .flip-card-inner { transform: rotateY(180deg); }
.flip-card-front, .flip-card-back {
  position: absolute; inset: 0;
  backface-visibility: hidden;
}
.flip-card-back { transform: rotateY(180deg); }
\end{lstlisting}

\medskip

\textbf{\faScroll\ Anima\textcommabelow{t}ii legate de scroll \textendash{} \texttt{animation-timeline}}

CSS modern (Chrome 115+, 2023) permite legarea unei animatii de pozitia scroll-ului, fara JavaScript:

\begin{lstlisting}[style=css, caption={Bara de progres care urmare\textcommabelow{s}te scroll-ul}]
@keyframes growProgress {
  from { transform: scaleX(0); }
  to   { transform: scaleX(1); }
}

.progress-bar {
  position: fixed; top: 0; left: 0;
  height: 4px; width: 100%;
  background: var(--accent);
  transform-origin: left;
  animation: growProgress linear;
  animation-timeline: scroll(root);  /* legat de scroll-ul paginii */
}
\end{lstlisting}

\code{view-timeline} este varianta complementar\u a: anima\textcommabelow{t}ia avanseaza pe masura ce un element intra in viewport. Util pentru fade-in la scroll, dar fara cost JavaScript.

\begin{nota}
Suportul pentru \code{animation-timeline} este partial in 2026: Chrome/Edge da, Firefox in spatele unui flag, Safari in pregatire. Foloseste cu fallback (animatia pur si simplu nu ruleaza in browserele vechi \textendash{} progressive enhancement OK).
\end{nota}

\medskip

\textbf{\faMagic\ \texttt{filter} \textcommabelow{s}i \texttt{backdrop-filter} \textendash{} efecte de imagine}

\code{filter} aplica un efect vizual pe element. \code{backdrop-filter} aplica acela\textcommabelow{s}i efect pe ce este \emph{in spate}: a\textcommabelow{s}a se face glassmorphism-ul (efect ,,sticla mata'').

\begin{lstlisting}[style=css, caption={Header transparent cu blur pe continutul dedesubt}]
header {
  background: rgba(255, 252, 246, 0.7);
  backdrop-filter: blur(12px) saturate(180%);
  -webkit-backdrop-filter: blur(12px) saturate(180%);  /* Safari */
}
\end{lstlisting}

Functii utile: \code{blur(8px)}, \code{brightness(1.2)}, \code{contrast(1.1)}, \code{grayscale(80\%)}, \code{hue-rotate(15deg)}, \code{drop-shadow(0 4px 8px rgba(0,0,0,0.2))}. Pot fi combinate prin spatiere.

\medskip

\textbf{\faCut\ \texttt{clip-path} \textendash{} forme custom \textcommabelow{s}i reveal effects}

\code{clip-path} ,,decupeaza'' elementul dupa o forma. Util pentru hero-uri oblice, butoane in forma de chevron, anima\textcommabelow{t}ii de tip ,,revelare'':

\begin{lstlisting}[style=css, caption={Imagine cu margine diagonala}]
.hero-image {
  clip-path: polygon(0 0, 100% 0, 100% 85%, 0 100%);
}
\end{lstlisting}

\begin{lstlisting}[style=css, caption={Reveal animat: pagina se ,,deschide'' la load}]
@keyframes reveal {
  from { clip-path: inset(0 100% 0 0); }
  to   { clip-path: inset(0 0 0 0); }
}
.hero h1 { animation: reveal 0.8s ease-out forwards; }
\end{lstlisting}

\begin{ponturi}
\textbf{Tool vizual:} \url{https://bennettfeely.com/clippy/} \textendash{} editor SVG-like care genereaza valoarea \code{clip-path} pentru orice forma desenata.
\end{ponturi}

\medskip

\textbf{\faMobile*\ Container queries \textendash{} ,,media queries'' pentru componente}

Problema cu \code{@media}: depinde de viewport, nu de containerul real al componentei. Un card de \code{300px} se comporta diferit intr-o coloana ingusta vs pe ecran lat, dar \code{@media} nu vede asta.

\code{@container} (CSS 2023, suport stabil) rezolva exact asta:

\begin{lstlisting}[style=css, caption={Card responsive in functie de container, nu de viewport}]
.card-wrapper { container-type: inline-size; }

.card { display: block; }

@container (min-width: 400px) {
  .card { display: flex; gap: var(--space-2); }
}
\end{lstlisting}

Acela\textcommabelow{s}i card devine ,,horizontal'' cand are spatiu, ,,vertical'' cand nu \textendash{} fara sa-i pese de marimea ecranului. Standardul modern pentru design system-uri.

\medskip

\textbf{\faRoute\ Unde mergi de aici}

\begin{itemize}[itemsep=1pt]
  \item \emph{Animatable properties:} \url{https://developer.mozilla.org/en-US/docs/Web/CSS/CSS_animated_properties}
  \item \emph{Scroll-driven animations:} \url{https://developer.chrome.com/docs/css-ui/scroll-driven-animations}
  \item \emph{Container queries:} \url{https://web.dev/learn/design/container-queries}
\end{itemize}
\end{bonusbox}

% =========================================================
\section*{Verificare cuno\textcommabelow{s}tin\textcommabelow{t}e}

\begin{selfcheck}
\begin{enumerate}[itemsep=2pt]
  \item Care este specificitatea pentru selectorul \code{.btn.primary}? Comparativ, ce specificitate are \code{button\#submit}?
  \item De ce strategia mobile-first este preferabila pentru contextul actual al traficului web?
  \item \emph{Spot the issue:} \code{.btn \{ background: red !important; \}} repetat in trei locuri.
  \item Care este diferen\textcommabelow{t}a func\textcommabelow{t}ionala intre \code{rem} si \code{px} pentru \code{font-size}?
  \item Ce face \code{box-sizing: border-box} fa\textcommabelow{t}a de comportamentul implicit?
  \item Pentru un hover smooth pe un buton, ce duration si timing-function alegi? Argumenteaza.
  \item De ce \code{transform: translateY(-2px)} pe hover este preferabil fa\textcommabelow{t}a de \code{top: -2px}?
  \item Cand folose\textcommabelow{s}ti \code{animation-fill-mode: forwards}? Da un exemplu concret.
\end{enumerate}
\end{selfcheck}

% =========================================================
\section*{Recapitulare}

\begin{recap}[Faza 2 -- CSS]
\begin{itemize}[itemsep=1pt]
  \item Stiluri externe in \code{<link rel="stylesheet">}; nu inline.
  \item Specificitate: \code{(id, clasa, tag)}; \code{!important} este \emph{antipattern}.
  \item \code{box-sizing: border-box} aplicat global pentru calcul predictibil.
  \item Unitati: \code{rem} pentru fonturi/spatiere; \code{px} pentru border; \code{\%}/\code{vw}/\code{vh} pentru layout responsiv.
  \item Variabilele CSS in \code{:root}; redefinire locala pentru teme contextuale.
  \item Flexbox: 6 proprietati cheie (\code{display}, \code{flex-direction}, \code{justify-content}, \code{align-items}, \code{gap}, \code{flex: 1}).
  \item Mobile-first cu \code{@media (min-width: ...)}, nu desktop-first cu \code{max-width}.
  \item Pozitionare: \code{static} (default), \code{relative}, \code{absolute}, \code{fixed}, \code{sticky}.
  \item Anima\textcommabelow{t}ii: \code{transition} pentru schimbari de stare (hover/focus), \code{@keyframes} pentru animatii arbitrare.
  \item Timing function: \code{ease-out} pentru intrari, \code{ease-in} pentru iesiri, \code{cubic-bezier()} pentru curbe custom.
  \item Performance: anima doar \code{transform} \textcommabelow{s}i \code{opacity} (accelerate GPU). Evita \code{top}/\code{left}/\code{width} \textendash{} declan\textcommabelow{s}eaza reflow.
  \item Pattern hover-lift: \code{transform: translateY(-4px) + box-shadow} cu \code{transition} pe ambele.
  \item \code{animation-fill-mode: forwards} pastreaza ultimul keyframe (esential pentru fade-in care ramane vizibil).
  \item Accesibilitate: respecta \code{prefers-reduced-motion}.
  \item \faRocket\ Bonus: 3D transforms, scroll-driven animations, \code{filter}, \code{clip-path}, container queries.
\end{itemize}
\end{recap}

\partdivider

% =========================================================
\begin{joc}
\textbf{Platforma:} Flexbox Froggy (\url{https://flexboxfroggy.com}).

\medskip

\textbf{Activitatea.} Aplicatia este un set de 24 de provocari care necesita pozitionarea unor obiecte pe o grila prin folosirea proprietatilor \code{justify-content}, \code{align-items}, \code{flex-direction}, \code{order}, \code{flex-wrap}.

\medskip

\textbf{Obiectiv:} primele 6 niveluri (acopera \code{justify-content} si \code{align-items}, cele mai folosite proprietati flexbox in practica).

\medskip

\textbf{Durata estimata:} 5 minute.
\end{joc}

% =========================================================
\section*{\faGraduationCap\ \ Resurse pentru aprofundare}

\begin{itemize}[itemsep=2pt]
  \item Mozilla Developer Network. \emph{CSS reference}. \url{https://developer.mozilla.org/en-US/docs/Web/CSS}
  \item Google. \emph{Learn CSS}. \url{https://web.dev/learn/css}
  \item Coyier C. \emph{A Complete Guide to Flexbox}. CSS-Tricks. \\\url{https://css-tricks.com/snippets/css/a-guide-to-flexbox/}
  \item Coyier C. \emph{A Complete Guide to Grid}. CSS-Tricks. \\\url{https://css-tricks.com/snippets/css/complete-guide-grid/}
  \item The Odin Project. \emph{CSS Foundations}. \\\url{https://www.theodinproject.com/paths/foundations/courses/foundations\#css-foundations}
  \item Aplicatii similare: \url{https://flukeout.github.io} (selectori), \url{https://cssgridgarden.com} (grid), \url{https://cssbattle.dev} (speedrun).
\end{itemize}

\bigskip

\bridge{Capitolul 3 (JavaScript) adauga interactivitate. Pagina stilizata in Faza 2 va raspunde la click-uri, va schimba dark mode, va procesa formulare \textendash{} fara reincarcare.}

% Fragment -- inclus de main.tex prin \input
\clearpage
\setpartlabel{themeJs}{\faBolt}{Faza 3: JavaScript}
\setcounter{section}{0}
\phantomsection
\addcontentsline{toc}{part}{\protect\textbf{Faza 3 \textendash{} JavaScript}}
\handout{Faza 3 \textendash{} JavaScript}{Variabile. Functii. Pattern-ul SELECT-LISTEN-CHANGE. Manipulare DOM.}

JavaScript este limbajul de programare standardizat ECMAScript (ECMA-262), interpretat nativ de toate browserele moderne. Spre deosebire de HTML (structura) si CSS (aspect), JavaScript adauga \textbf{comportament}: raspunde la actiunile utilizatorului, modifica continutul paginii dinamic, comunica cu servere prin retea.

Capitolul prezinta sintaxa de baza, modelul de programare orientat-eveniment specific paginilor web statice, si interactiunea cu arborele DOM (\emph{Document Object Model}).

\bigskip

\begin{outcomes}
Dupa capitolul asta vei \textcommabelow{s}ti:
\begin{itemize}[itemsep=1pt]
  \item Declara variabile cu \code{const} \textcommabelow{s}i \code{let}, recunoscand cand fiecare se aplica.
  \item Aplica pattern-ul \textbf{SELECT $\to$ LISTEN $\to$ CHANGE} pentru orice interac\textcommabelow{t}iune.
  \item Manipuleaza DOM-ul cu \code{querySelector}, \code{classList}, \code{addEventListener}.
  \item Intercepteaza submit-ul unui formular \textcommabelow{s}i extrage valorile cu \code{FormData}.
  \item Folose\textcommabelow{s}te \code{IntersectionObserver} pentru anima\textcommabelow{t}ii \textcommabelow{s}i lazy loading legate de scroll.
  \item Diagnostiheaza erorile JavaScript folosind console-ul browserului.
  \item \faRocket\ \emph{Bonus self-study:} debounce/throttle, event delegation, \code{requestAnimationFrame}, \code{fetch} cu async/await, metode pe array (\code{map}/\code{filter}), \code{localStorage}.
\end{itemize}
\end{outcomes}

\begin{whymatters}
HTML \textcommabelow{s}i CSS construiesc o pagina statica. JavaScript transforma pagina intr-o aplica\textcommabelow{t}ie: poate reac\textcommabelow{t}iona la click-uri, poate schimba con\textcommabelow{t}inutul fara reincarcare, poate trimite cereri catre servere. Pattern-ul SELECT-LISTEN-CHANGE prezentat aici acopera 80\% din scenariile de interactivitate \textcommabelow{s}i este aceea\textcommabelow{s}i abstrac\textcommabelow{t}ie pe care React, Vue \textcommabelow{s}i Svelte o ambala in API-uri declarative.
\end{whymatters}

\begin{nota}[Pre-rechizite]
Capitolul presupune in\textcommabelow{t}elegerea structurii HTML (tag-uri, atribute) \textcommabelow{s}i a selectorilor CSS (clase, id-uri). Faze 1 \textcommabelow{s}i 2.
\end{nota}

% =========================================================
\section{Conectarea fisierului JavaScript la HTML}

Includerea unui fisier JavaScript se realizeaza prin elementul \code{<script>}, plasat in \code{<head>} cu atributul \code{defer}:

\begin{lstlisting}[style=html]
<script src="script.js" defer></script>
\end{lstlisting}

\subsection{\texttt{defer} vs \texttt{async} vs nimic}

\begin{itemize}
  \item \textbf{Fara atribut.} Browserul opreste parsarea HTML-ului, descarca si ruleaza scriptul, apoi continua. Blocheaza randarea \textendash{} antipattern.
  \item \code{async} \textendash{} descarcare in paralel cu parsarea, dar executie imediat ce scriptul este disponibil (poate intrerupe parsarea). Util doar pentru scripturi independente (analytics).
  \item \code{defer} \textendash{} descarcare in paralel, executie \textbf{dupa} parsarea HTML-ului. Garantat executat in ordinea declararii. Optiunea recomandata pentru codul propriu.
\end{itemize}

\begin{atentie}
Fara \code{defer}, scriptul ruleaza inainte de parsarea \code{<body>}: \code{document.querySelector('.cta')} returneaza \code{null}. Cu \code{defer}, intregul DOM exista cand ruleaza scriptul.
\end{atentie}

% =========================================================
\section{Declararea variabilelor}

ECMAScript 2015 introduce doua cuvinte cheie pentru declararea variabilelor: \code{const} (valoare imuabila) si \code{let} (valoare mutabila). Cuvantul cheie \code{var} este considerat depasit din cauza regulilor de scope.

\begin{lstlisting}[style=js, caption={Reguli de declarare}]
const PI = 3.14;            // valoare imutabila
let counter = 0;            // valoare mutabila

counter = counter + 1;      // permis
// PI = 3.15;               // TypeError la rulare
\end{lstlisting}

\subsection{Diferenta dintre \texttt{var} si \texttt{let} / \texttt{const}}

\code{var} are \emph{function scope}: o variabila declarata in interiorul unui \code{if} sau \code{for} este accesibila in tot blocul functiei. \code{let} si \code{const} au \emph{block scope}: vizibile doar in interiorul \code{\{...\}} in care au fost declarate.

\begin{dano}[Asa nu \textendash{} \texttt{var}, scope confuz]
\begin{lstlisting}[style=js, numbers=none]
function exemplu() {
  if (true) {
    var x = 10;
  }
  console.log(x);  // 10 -- scapa din if!
}
\end{lstlisting}
\end{dano}

\begin{dado}[Asa da \textendash{} \texttt{let}, scope previzibil]
\begin{lstlisting}[style=js, numbers=none]
function exemplu() {
  if (true) {
    let x = 10;
  }
  console.log(x);  // ReferenceError -- nu exista in afara if-ului
}
\end{lstlisting}
\end{dado}

\textbf{Conventia practica:} \code{const} este valoarea implicita; \code{let} se foloseste doar cand reasignarea este necesara; \code{var} se evita complet in cod nou.

% =========================================================
\section{Tipuri primitive si compuse}

\begin{lstlisting}[style=js]
const nume      = "Voluntar";              // string
const varsta    = 21;                   // number (intregi si zecimale)
const eStudent  = true;                 // boolean
const limbaje   = ["HTML", "CSS"];      // array (lista ordonata)
const persoana  = {                     // object (pereche cheie-valoare)
  nume:   "Voluntar",
  varsta: 21,
};

console.log(persoana.nume);             // "Voluntar"
console.log(limbaje[0]);                // "HTML"
console.log(typeof varsta);             // "number"
\end{lstlisting}

JavaScript este un limbaj cu \emph{tipare dinamica}: tipul unei variabile este determinat de valoarea atribuita si poate fi modificat in timpul executiei.

\subsection{Verificarea egalitatii: \texttt{===} vs \texttt{==}}

Operatorul \code{==} efectueaza \emph{coercitie de tip}: converteste operanzii inainte de comparatie. Rezultatul este des contraintuitiv. Operatorul \code{===} compara strict valoare si tip.

\begin{dano}[Asa nu \textendash{} \texttt{==}, coercitie surprinzatoare]
\begin{lstlisting}[style=js, numbers=none]
0 == "0"             // true
0 == ""              // true
"" == false          // true
null == undefined    // true
[] == false          // true
\end{lstlisting}
\end{dano}

\begin{dado}[Asa da \textendash{} \texttt{===} mereu]
\begin{lstlisting}[style=js, numbers=none]
0 === "0"            // false (number vs string)
0 === ""             // false
null === undefined   // false
\end{lstlisting}
\end{dado}

% =========================================================
\section{Functii: declaratie, expresie, arrow}

JavaScript ofera trei sintaxe pentru definirea unei functii. Cele trei comportamente difera la nivel de \code{this} si hoisting, dar pentru cazul uzual (handler de eveniment) sunt interschimbabile.

\begin{lstlisting}[style=js, caption={Trei moduri de a scrie aceeasi functie}]
// 1. Declaratie de functie (function declaration)
function saluta(nume) {
  return `Salut, ${nume}!`;
}

// 2. Expresie de functie (function expression)
const saluta2 = function(nume) {
  return `Salut, ${nume}!`;
};

// 3. Arrow function (sintaxa concisa)
const saluta3 = (nume) => `Salut, ${nume}!`;

// Apel: identic in toate cazurile
console.log(saluta("Voluntar"));
\end{lstlisting}

\textbf{Conventia practica:} arrow functions pentru handlere si callbacks; declaratii pentru functii cu nume reutilizate, definite la nivelul fisierului.

% =========================================================
\section{Pattern-ul SELECT $\to$ LISTEN $\to$ CHANGE}

Majoritatea interactiunilor pe paginile statice respecta o schema in trei pasi: selectarea elementului DOM, atasarea unui ascultator de eveniment, modificarea elementului in raspuns la eveniment.

\begin{lstlisting}[style=js, caption={Toggle pentru tema dark/light}]
// 1. SELECT
const btn = document.querySelector('.toggle-theme');

// 2. LISTEN
btn.addEventListener('click', () => {

  // 3. CHANGE
  document.body.classList.toggle('dark');

});
\end{lstlisting}

\begin{ponturi}
Pattern-ul descris reprezinta abstractizarea fundamentala pe care framework-urile moderne (React, Vue, Svelte, Angular) o ascund in spatele propriilor API-uri declarative. Intelegerea variantei \emph{vanilla} faciliteaza tranzitia ulterioara catre framework-uri.
\end{ponturi}

% =========================================================
\section{Manipularea DOM-ului}

API-ul \code{document} expune metode pentru selectia si modificarea nodurilor:

\begin{lstlisting}[style=js, caption={Operatii uzuale}]
// SELECT -- un singur element
const el = document.querySelector('.cta');           // primul .cta
const id = document.getElementById('contact');       // alternativa pt id

// SELECT -- mai multe elemente
const linkuri = document.querySelectorAll('nav a');
linkuri.forEach((link) => {
  link.addEventListener('click', () => console.log('click'));
});

// CHANGE -- continut
el.textContent = 'Salut!';                  // text simplu (sigur)
el.innerHTML  = '<strong>Salut!</strong>';  // HTML (atentie XSS)

// CHANGE -- atribute
el.setAttribute('aria-expanded', 'true');
el.classList.add('active');
el.classList.remove('hidden');
el.classList.toggle('open');

// CHANGE -- creare element nou
const nou = document.createElement('li');
nou.textContent = 'Element nou';
document.querySelector('ul').appendChild(nou);
\end{lstlisting}

\begin{dano}[Asa nu \textendash{} querySelector la fiecare apel]
\begin{lstlisting}[style=js, numbers=none]
btn.addEventListener('click', () => {
  document.querySelector('.modal').classList.add('open');
  document.querySelector('.modal').focus();
  document.querySelector('.modal').scrollIntoView();
});
\end{lstlisting}
\textbf{Problema.} Trei traversari ale DOM-ului pentru aceeasi referinta.
\end{dano}

\begin{dado}[Asa da \textendash{} cache local]
\begin{lstlisting}[style=js, numbers=none]
const modal = document.querySelector('.modal');

btn.addEventListener('click', () => {
  modal.classList.add('open');
  modal.focus();
  modal.scrollIntoView();
});
\end{lstlisting}
\end{dado}

% =========================================================
\section{Procesarea formularelor}

Submit-ul nativ al unui formular trimite o cerere HTTP catre URL-ul declarat in \code{action} si reincarca pagina. In majoritatea cazurilor, fluxul dorit este interceptarea evenimentului si manipularea datelor in JavaScript:

\begin{lstlisting}[style=js, caption={Captura datelor de formular fara reincarcarea paginii}]
const form = document.querySelector('#contact-form');

form.addEventListener('submit', (e) => {
  e.preventDefault();                                   // blocheaza reincarcarea

  const data = Object.fromEntries(new FormData(form));  // {nume, email, mesaj}
  console.log('Date primite:', data);

  form.reset();                                         // goleste campurile
  alert(`Multumesc, ${data.nume}!`);
});
\end{lstlisting}

Fluxul: \code{e.preventDefault()} suspenda comportamentul implicit; \code{new FormData(form)} colecteaza valorile dupa atributul \code{name}; \code{Object.fromEntries(...)} le converteste intr-un obiect uzual; \code{form.reset()} reseteaza valorile in interfata.

\begin{ponturi}
Pentru transmiterea efectiva a datelor catre un endpoint extern (fara backend propriu), serviciul \href{https://formspree.io}{Formspree} ofera un endpoint gratuit. Configurarea consta intr-o singura modificare a atributului \code{action}.
\end{ponturi}

% =========================================================
\section{Debugging si console}

Console-ul browserului (F12 $\to$ tab \emph{Console}) este principalul instrument de diagnostic. Erorile JavaScript apar automat aici cu mesaj, fisier si linie.

\begin{lstlisting}[style=js, caption={Trei moduri de a inspecta valorile}]
// Cel mai uzual: log valoare cu eticheta
console.log('Date primite:', data);

// Tabel pentru obiecte / array-uri (formatare imbunatatita)
console.table(linkuri);

// Pauza in executie -- DevTools deschis
debugger;
\end{lstlisting}

\textbf{Sfaturi practice.} Inainte de a cere ajutor, citeste mesajul de eroare din console. Numele functiei si numarul liniei sunt aproape intotdeauna corecte. \code{console.log(typeof x)} verifica tipul cand comportamentul este surprinzator.

% =========================================================
\section{IntersectionObserver \textendash{} anima\textcommabelow{t}ii la scroll}

Cea mai uzuala cerere ,,fa o anima\textcommabelow{t}ie cand sectiunea devine vizibila'' (fade-in pe scroll, reveal pe carduri, lazy loading pentru imagini) NU se rezolva cu \code{window.addEventListener('scroll', ...)}. Browserul ofera un API dedicat, mult mai eficient: \code{IntersectionObserver}.

Idea: in loc sa intrebi de 60 de ori pe secunda ,,e vizibila sectiunea?'', tu declari ,,anun\textcommabelow{t}a-ma cand starea de vizibilitate se schimba''. Browserul face restul, fara cost de performance.

\begin{lstlisting}[style=js, caption={Fade-in pe scroll: pattern canonic in 10 linii}]
const observer = new IntersectionObserver((entries) => {
  entries.forEach((entry) => {
    if (entry.isIntersecting) {
      entry.target.classList.add('visible');
      observer.unobserve(entry.target);  // o singura data, nu repeti
    }
  });
}, { threshold: 0.15 });  // declanseaza cand 15% e in viewport

document.querySelectorAll('.reveal').forEach((el) => observer.observe(el));
\end{lstlisting}

Pe CSS, partea sora:

\begin{lstlisting}[style=css, numbers=none]
.reveal { opacity: 0; transform: translateY(20px); transition: all 0.6s ease-out; }
.reveal.visible { opacity: 1; transform: translateY(0); }
\end{lstlisting}

Cum se cite\textcommabelow{s}te API-ul:
\begin{itemize}[itemsep=1pt]
  \item Constructorul: \code{new IntersectionObserver(callback, optiuni)}.
  \item \code{callback(entries)} \textendash{} se ape\-leaza cand starea de vizibilitate se schimba.
  \item \code{entry.isIntersecting} \textendash{} \code{true} cand elementul intra in viewport.
  \item \code{entry.target} \textendash{} elementul DOM (acela\textcommabelow{s}i pe care l-ai dat la \code{observe}).
  \item \code{threshold} \textendash{} procent din element vizibil pentru a declan\textcommabelow{s}a (\code{0} = orice pixel, \code{1} = complet vizibil, \code{0.15} = 15\%).
  \item \code{rootMargin} \textendash{} extinde/contracta viewport-ul (ex. \code{"0px 0px -100px 0px"} declan\textcommabelow{s}eaza cu 100px inainte sa intre real in cadru).
  \item \code{observer.unobserve(el)} \textendash{} opre\textcommabelow{s}te monitorizarea (perfect dupa prima activare).
\end{itemize}

\begin{dano}[Asa nu \textendash{} scroll handler cu \texttt{getBoundingClientRect}]
\begin{lstlisting}[style=js, numbers=none]
window.addEventListener('scroll', () => {
  document.querySelectorAll('.reveal').forEach((el) => {
    if (el.getBoundingClientRect().top < window.innerHeight) {
      el.classList.add('visible');
    }
  });
});
\end{lstlisting}
\textbf{Problema.} Ruleaza la fiecare pixel de scroll (zeci \textendash{} sute de ori pe secunda), forteaza recalculare layout (\emph{layout thrashing}), scade FPS pe mobile.
\end{dano}

\begin{dado}[Asa da \textendash{} IntersectionObserver]
Browserul observa eficient (compus pe thread separat). Callback-ul ruleaza doar cand starea se schimba. Zero impact pe performance.
\end{dado}

\begin{ponturi}[Aplica\textcommabelow{t}ii directe]
\begin{itemize}[itemsep=1pt]
  \item Lazy loading custom pentru con\textcommabelow{t}inut dinamic (cand nu poti folosi \code{loading="lazy"}).
  \item Active state pentru navigare \emph{scrollspy} (eviden\textcommabelow{t}iaza link-ul activ in nav).
  \item Infinite scroll \textendash{} cand userul ajunge aproape de jos, incarci urmatoarea pagina.
  \item Analytics ,,a vazut sec\textcommabelow{t}iunea X'' (tracking eveniment doar cand sec\textcommabelow{t}iunea devine vizibila).
\end{itemize}
\end{ponturi}

% =========================================================
\section{Tratarea erorilor (\texttt{try}/\texttt{catch})}

Codul care poate gen\-era erori la rulare (parsare JSON, apeluri \code{fetch}, acces la API-uri externe) trebuie incadrat intr-un bloc \code{try}/\code{catch} pentru ca o eroare sa nu opreasca executia restului paginii:

\begin{lstlisting}[style=js, caption={Parsare JSON cu fallback}]
const raw = localStorage.getItem('preferinte');

try {
  const config = JSON.parse(raw);
  console.log('Preferinte incarcate:', config);
} catch (err) {
  console.warn('JSON invalid in localStorage; folosesc valori implicite');
  // continuare cu valori default
}
\end{lstlisting}

Cand eroarea este previzibila si recuperabila (utilizatorul poate continua), \code{try}/\code{catch} este abordarea corecta. Pentru erori neprevazute, evenimentul global \code{window.addEventListener('error', ...)} permite logarea catre un serviciu de monitoring.

% =========================================================
\section{Erori frecvente}

\begin{itemize}
  \item \textbf{\code{Cannot read property '...' of null}.} Cauza: \code{querySelector} a returnat \code{null} si codul a incercat sa acceseze o proprietate. Diagnostic: \code{console.log(el)} inainte de a folosi.
  \item \textbf{\code{addEventListener is not a function}.} Aceeasi cauza: rezultatul \code{querySelector} este \code{null}.
  \item \textbf{Hoisting.} Variabilele \code{let} si \code{const} declarate intr-un bloc nu sunt accesibile inainte de declaratie. Accesul anticipat genereaza \emph{ReferenceError: Cannot access before initialization}.
  \item \textbf{Comparatie cu \code{undefined}.} Foloseste \code{x === undefined}, nu \code{!x} \textendash{} ultimul prinde si \code{0}, \code{""}, \code{false}, \code{null}.
\end{itemize}

% =========================================================
\begin{bonusbox}
\textbf{Acest box e optional.} Daca esti la inceput, treci direct la \emph{Verificare cuno\textcommabelow{s}tin\textcommabelow{t}e}. Mai jos sunt 8 pattern-uri pe care le folose\textcommabelow{s}ti zilnic in proiecte reale.

\medskip

\textbf{\faClock\ Debounce \textcommabelow{s}i throttle \textendash{} st\u apane\textcommabelow{s}te evenimentele frecvente}

Evenimentele \code{scroll}, \code{resize}, \code{input}, \code{mousemove} se declan\textcommabelow{s}eaza de zeci \textendash{} sute de ori pe secunda. Daca handler-ul face ceva costisitor (fetch, recalcul layout), pagina ,,inghea\textcommabelow{t}a''.

\emph{Debounce} \textendash{} amana executia pana cand evenimentul nu mai apare \emph{ms} milisecunde. Bun pentru autocomplete (a\textcommabelow{s}tepta sa termini de scris).

\emph{Throttle} \textendash{} executa cel mult o data la fiecare \emph{ms}. Bun pentru scroll, resize.

\begin{lstlisting}[style=js, caption={Utilitar reutilizabil}]
function debounce(fn, ms = 300) {
  let timer;
  return function(...args) {
    clearTimeout(timer);
    timer = setTimeout(() => fn.apply(this, args), ms);
  };
}

const search = document.querySelector('#search');
search.addEventListener('input', debounce((e) => {
  console.log('Cautare:', e.target.value);  // ruleaza dupa 300ms de tacere
}, 300));
\end{lstlisting}

\begin{lstlisting}[style=js, caption={Buton ,,scroll to top'' care apare dupa 400px}]
function throttle(fn, ms = 100) {
  let last = 0;
  return function(...args) {
    const now = Date.now();
    if (now - last >= ms) { last = now; fn.apply(this, args); }
  };
}

const btnTop = document.querySelector('.to-top');
window.addEventListener('scroll', throttle(() => {
  btnTop.classList.toggle('visible', window.scrollY > 400);
}, 100));
\end{lstlisting}

\medskip

\textbf{\faSitemap\ Event delegation \textendash{} un singur listener pentru N elemente}

Daca ai 50 de butoane \code{.delete-btn}, NU pune 50 de \code{addEventListener}. Pune unul pe parinte si folose\textcommabelow{s}te \code{event.target.closest()} pentru a verifica ce a fost clicat. Bonus: func\textcommabelow{t}ioneaza \textcommabelow{s}i pentru elemente adaugate dinamic dupa atasare.

\begin{lstlisting}[style=js, caption={Lista de task-uri cu \u stergere}]
const list = document.querySelector('#task-list');

list.addEventListener('click', (e) => {
  const deleteBtn = e.target.closest('.delete-btn');
  if (!deleteBtn) return;                          // click in alta parte
  deleteBtn.closest('li').remove();
});
\end{lstlisting}

Cum func\textcommabelow{t}ioneaza: evenimentul \emph{bubbles} (urca de la \code{target} catre parinti). Listener-ul pe parinte vede toate click-urile copiilor.

\medskip

\textbf{\faCog\ Op\textcommabelow{t}iuni la \texttt{addEventListener}}

Parametrul al treilea poate fi un obiect cu op\textcommabelow{t}iuni:

\begin{itemize}[itemsep=1pt]
  \item \code{\{ once: true \}} \textendash{} executa o singura data, apoi se auto-elimina. Util pentru splash/intro animation.
  \item \code{\{ passive: true \}} \textendash{} promiti browserului ca NU vei face \code{preventDefault()}. Pentru \code{scroll}/\code{touchmove} mai ales \textendash{} cre\textcommabelow{s}te FPS pe mobile.
  \item \code{\{ capture: true \}} \textendash{} prinde evenimentul in faza de \emph{capture} (de sus in jos), nu \emph{bubble} (rar folosit).
\end{itemize}

\begin{lstlisting}[style=js, numbers=none]
window.addEventListener('scroll', onScroll, { passive: true });
btn.addEventListener('click', showWelcome, { once: true });
\end{lstlisting}

\medskip

\textbf{\faPlay\ \texttt{requestAnimationFrame} \textendash{} anima\textcommabelow{t}ii fluide din JS}

Cand vrei sa animezi ceva din JavaScript (nu doar prin CSS), NU folosi \code{setTimeout(..., 16)}. Folose\textcommabelow{s}te \code{requestAnimationFrame} \textendash{} se sincronizeaza cu refresh-ul ecranului (60fps natural, 120fps pe ecrane rapide).

\begin{lstlisting}[style=js, caption={Counter animat de la 0 la o valoare \textcommabelow{t}int\u a}]
function animateNumber(el, target, duration = 1000) {
  const start = performance.now();

  function tick(now) {
    const elapsed = now - start;
    const progress = Math.min(elapsed / duration, 1);   // 0..1
    el.textContent = Math.floor(progress * target);
    if (progress < 1) requestAnimationFrame(tick);
  }

  requestAnimationFrame(tick);
}

animateNumber(document.querySelector('.stat-projects'), 42);
\end{lstlisting}

\medskip

\textbf{\faGlobe\ Async / await + fetch \textendash{} apeluri spre API-uri}

Functiile \code{async} permit scrierea codului asincron in stil sincron. \code{await} pauzeaza pana cand un \code{Promise} se rezolva.

\begin{lstlisting}[style=js, caption={Voluntar afi\textcommabelow{s}eaz\u a stats GitHub live pe portofoliu}]
async function loadGitHubStats(username) {
  try {
    const response = await fetch(`https://api.github.com/users/${username}`);
    if (!response.ok) throw new Error(`HTTP ${response.status}`);

    const data = await response.json();

    document.querySelector('.gh-repos').textContent     = data.public_repos;
    document.querySelector('.gh-followers').textContent = data.followers;
    document.querySelector('.gh-avatar').src            = data.avatar_url;
  } catch (err) {
    console.warn('Nu am putut incarca stats GitHub:', err);
  }
}

loadGitHubStats('voluntar-best');
\end{lstlisting}

Trei reguli:
\begin{itemize}[itemsep=1pt]
  \item \code{async} inainte de \code{function} \textendash{} func\textcommabelow{t}ia returneaza un \code{Promise}.
  \item \code{await} func\textcommabelow{t}ioneaza doar in interiorul unei func\textcommabelow{t}ii \code{async}.
  \item Erorile asincrone se prind cu \code{try}/\code{catch} la fel ca cele sincrone.
\end{itemize}

\medskip

\textbf{\faList\ Metode pe array \textendash{} \texttt{map}, \texttt{filter}, \texttt{find}}

In loc de \code{for} cu indice si \code{array.push()}, JavaScript ofera metode care returneaza un array nou (imuabilitate). Mult mai citibile pentru transformari de date.

\begin{lstlisting}[style=js, caption={Lista de proiecte randat\u a din date}]
const proiecte = [
  { titlu: 'Portfolio',   link: '/portfolio',   tag: 'web' },
  { titlu: 'Game of Life', link: '/gol',         tag: 'algo' },
  { titlu: 'Calculator',   link: '/calc',        tag: 'web' },
];

// map -- transforma fiecare element
const html = proiecte
  .filter((p) => p.tag === 'web')                      // doar proiectele web
  .map((p) => `<li><a href="${p.link}">${p.titlu}</a></li>`)
  .join('');

document.querySelector('.projects').innerHTML = html;

// find -- gase\textcommabelow{s}te primul element care satisface
const portfolio = proiecte.find((p) => p.titlu === 'Portfolio');
\end{lstlisting}

\code{reduce} (mai rar in front-end de inceput) acumuleaza un singur rezultat din array (suma, grup\u ari, statistici).

\medskip

\textbf{\faPuzzlePiece\ Destructuring + template literals}

Doua sintaxe pe care le vezi peste tot. Destructuring extrage valori din obiecte/array-uri intr-un singur pas:

\begin{lstlisting}[style=js, numbers=none]
const { nume, varsta } = persoana;       // in loc de persoana.nume, persoana.varsta
const [primul, al_doilea] = limbaje;     // in loc de limbaje[0], limbaje[1]

// Aplicat in parametri de functie
function saluta({ nume, varsta }) {
  return `${nume}, ${varsta} ani`;
}
saluta({ nume: 'Voluntar', varsta: 21 });
\end{lstlisting}

Template literals (backtick) permit interpolare \code{\$\{...\}} \textcommabelow{s}i string-uri multi-linie. Le-ai folosit deja in exemplul GitHub si in HTML-ul generat din \code{map}.

\medskip

\textbf{\faSave\ \texttt{localStorage} \textendash{} salveaza preferin\textcommabelow{t}ele user-ului}

\code{localStorage} este un dictionar persistent intre vizite. Stocheaza doar string-uri \textendash{} pentru obiecte folose\textcommabelow{s}ti \code{JSON.stringify}/\code{JSON.parse}.

\begin{lstlisting}[style=js, caption={Tema dark se p\u astreaz\u a intre vizite}]
const toggle = document.querySelector('.toggle-theme');

// La incarcare: aplica tema salvata
if (localStorage.getItem('theme') === 'dark') {
  document.body.classList.add('dark');
}

// La click: schimba si salveaza
toggle.addEventListener('click', () => {
  document.body.classList.toggle('dark');
  const tema = document.body.classList.contains('dark') ? 'dark' : 'light';
  localStorage.setItem('theme', tema);
});
\end{lstlisting}

\code{sessionStorage} e identic, dar se sterge cand userul inchide tab-ul. Ambele au limita ~5MB per origin.

\medskip

\textbf{\faRoute\ Unde mergi de aici}

\begin{itemize}[itemsep=1pt]
  \item \emph{MDN IntersectionObserver:} \url{https://developer.mozilla.org/en-US/docs/Web/API/Intersection_Observer_API}
  \item \emph{javascript.info \textendash{} Promises, async/await:} \url{https://javascript.info/async}
  \item \emph{web.dev Event handling best practices:} \url{https://web.dev/articles/optimize-input-delay}
  \item \emph{Free public APIs pentru exersat:} \url{https://github.com/public-apis/public-apis}
\end{itemize}
\end{bonusbox}

% =========================================================
\section*{Verificare cuno\textcommabelow{s}tin\textcommabelow{t}e}

\begin{selfcheck}
\begin{enumerate}[itemsep=2pt]
  \item Ce face atributul \code{defer} pe un \code{<script>} \textcommabelow{s}i de ce este necesar?
  \item Care sunt cei trei pa\textcommabelow{s}i ai pattern-ului \textbf{SELECT $\to$ LISTEN $\to$ CHANGE}?
  \item De ce \code{===} este preferabil lui \code{==}? Da\textcommabelow{t}i un exemplu de coerci\textcommabelow{t}ie surprinzatoare.
  \item Ce face \code{e.preventDefault()} apelat intr-un handler de \code{submit}?
  \item \emph{Spot the issue:} \code{const btn = document.querySelector('.cta'); btn.addEventListener(...);} \textendash{} \code{btn} este \code{null}.
  \item De ce este \code{IntersectionObserver} preferabil lui \code{scroll} listener pentru anima\textcommabelow{t}ii la scroll?
  \item Ce face \code{observer.unobserve(el)} \textcommabelow{s}i cand il folose\textcommabelow{s}ti?
\end{enumerate}
\end{selfcheck}

% =========================================================
\section*{Recapitulare}

\begin{recap}[Faza 3 -- JavaScript]
\begin{itemize}[itemsep=1pt]
  \item Includere: \code{<script src="..." defer></script>} in \code{<head>}.
  \item \code{const} este declararea implicita; \code{let} doar cand reasignarea este necesara; nu \code{var}.
  \item Comparatie: exclusiv \code{===}, niciodata \code{==}.
  \item Pattern-cheie: \textbf{SELECT $\to$ LISTEN $\to$ CHANGE} acopera majoritatea interactiunilor.
  \item DOM: \code{querySelector} (un element); \code{querySelectorAll} (NodeList iterabil cu \code{.forEach}).
  \item Form submit: \code{e.preventDefault()} + \code{new FormData(form)} + \code{Object.fromEntries(...)}.
  \item Cache referintele DOM in variabile; nu repeta \code{querySelector} pentru acelasi element.
  \item \code{IntersectionObserver} pentru anima\textcommabelow{t}ii la scroll: zero cost de performance, callback doar cand vizibilitatea se schimba.
  \item Debug: \code{F12 $\to$ Console}, \code{console.log()}, \code{debugger}.
  \item \faRocket\ Bonus: debounce/throttle, event delegation, \code{requestAnimationFrame}, \code{fetch} cu async/await, \code{map}/\code{filter}/\code{find}, destructuring, \code{localStorage}.
\end{itemize}
\end{recap}

\partdivider

% =========================================================
\begin{joc}
\textbf{Platforma:} CodePen (\url{codepen.io}). \quad \textit{URL-ul exact este distribuit de trainer in chat.}

\medskip

\textbf{Activitatea.} Documentul de pornire contine HTML si CSS complet, dar fisierul JavaScript este gol. Sarcina este implementarea pattern-ului SELECT-LISTEN-CHANGE: la apasarea butonului prezent in pagina, clasa \code{dark} se aplica si se scoate alternativ pe elementul \code{<body>}.

\medskip

\textbf{Solutia asteptata:}

\begin{lstlisting}[style=js, numbers=none]
const btn = document.querySelector('button');
btn.addEventListener('click', () => {
  document.body.classList.toggle('dark');
});
\end{lstlisting}

\textbf{Durata estimata:} 5 minute.
\end{joc}

% =========================================================
\section*{\faGraduationCap\ \ Resurse pentru aprofundare}

\begin{itemize}[itemsep=2pt]
  \item Kantor I. \emph{The Modern JavaScript Tutorial}. \url{https://javascript.info}
  \item Mozilla Developer Network. \emph{JavaScript reference}. \\\url{https://developer.mozilla.org/en-US/docs/Web/JavaScript}
  \item The Odin Project. \emph{JavaScript Foundations}. \\\url{https://www.theodinproject.com/paths/foundations/courses/foundations\#javascript-basics}
  \item Mozilla Developer Network. \emph{Document Object Model}. \\\url{https://developer.mozilla.org/en-US/docs/Web/API/Document_Object_Model}
  \item Google. \emph{Learn JavaScript}. \url{https://web.dev/learn/javascript}
\end{itemize}

\bigskip

\bridge{Capitolul 4 (Deploy) muta site-ul de pe ma\textcommabelow{s}ina locala pe un URL public. Toate fazele 1\textendash{}3 sunt locale; abia dupa deploy pagina exista in lume.}

% Fragment -- inclus de main.tex prin \input
\clearpage
\setpartlabel{themeDeploy}{\faRocket}{Faza 4: Deploy}
\setcounter{section}{0}
\phantomsection
\addcontentsline{toc}{part}{\protect\textbf{Faza 4 \textendash{} Deploy}}
\handout{Faza 4 \textendash{} Deploy}{Versionare cu git. Publicare prin GitHub Pages.}

Capitolul descrie procesul de publicare a unui site web static prin combinatia git \textendash{} GitHub \textendash{} GitHub Pages. Acoperirea include comenzile esentiale ale fluxului de lucru, configurarea unui repository remote, activarea serviciului de hosting integrat, si convenii de scriere a mesajelor de commit.

\bigskip

\begin{outcomes}
Dupa capitolul asta vei \textcommabelow{s}ti:
\begin{itemize}[itemsep=1pt]
  \item Executa cele 5 comenzi git pentru ini\textcommabelow{t}ializarea unui repository \textcommabelow{s}i primul push.
  \item Configureaza GitHub Pages pe un repository pentru hosting static automat.
  \item Folose\textcommabelow{s}te \code{git status}, \code{git diff}, \code{git log} pentru inspectarea starii.
  \item Scrie mesaje de commit la imperativ, descriptive, sub 72 caractere.
  \item Excludre fi\textcommabelow{s}iere sensibile prin \code{.gitignore}.
\end{itemize}
\end{outcomes}

\begin{whymatters}
Site-ul construit local pe ma\textcommabelow{s}ina dezvoltatorului nu este accesibil nimanui altcuiva. Deploy-ul transforma cod local in URL public. In plus, git este standardul industrial pentru versionare \textendash{} cuno\textcommabelow{s}tin\textcommabelow{t}ele dobandite aici se aplica identic in orice job de dezvoltare software. Mesajele de commit bune sunt diferen\textcommabelow{t}a intre un istoric care ajuta la debug \textcommabelow{s}i unul inutil.
\end{whymatters}

\begin{nota}[Pre-rechizite]
Capitolul presupune un site func\textcommabelow{t}ional in localhost (Faza 1\textendash{}3) \textcommabelow{s}i un cont GitHub creat in prealabil.
\end{nota}

% =========================================================
\section{De ce git?}

Git este sistemul de control al versiunilor (\emph{Version Control System}) folosit de aproape toata industria software. Beneficiile fata de o solutie de tip ,,upload manual prin FTP'':

\begin{itemize}
  \item \textbf{Istoric.} Fiecare commit reprezinta un \emph{snapshot} al codului, identificat printr-un hash unic. Reverenta la o versiune anterioara este o operatie atomica.
  \item \textbf{Backup distribuit.} Repository-ul exista local, pe GitHub, si pe orice copie clonata. Pierderea unei singure copii nu afecteaza istoricul.
  \item \textbf{Colaborare.} Mai multi autori pot lucra simultan pe ramuri (\emph{branches}) separate, sincronizate prin merge.
  \item \textbf{Standard industrial.} Cunoasterea git este o cerinta de baza in posturile de dezvoltare software.
\end{itemize}

% =========================================================
\section{Modelul mental: trei zone}

Git mentine trei zone in care fisierele pot exista la un moment dat:

\begin{lstlisting}[style=base, basicstyle=\ttfamily\small, numbers=none]
  Working dir         Staging area         Repository
  (fisiere editate)   (pregatite pt.       (snapshot-uri
                       commit)              salvate permanent)

       |                   |                    |
       |  git add .        |  git commit -m     |
       |------------------>|------------------->|
       |                   |                    |
       |   git restore     |    git reset       |
       |<------------------|<-------------------|
       |                   |                    |
                                                |
                                                |  git push
                                                |--------->  GitHub
                                                |             (remote)
                                                |<---------
                                                |  git pull
\end{lstlisting}

\begin{itemize}
  \item \textbf{Working directory} \textendash{} fisierele asa cum apar pe disc.
  \item \textbf{Staging area} (sau \emph{index}) \textendash{} multimea fisierelor pregatite pentru urmatorul commit.
  \item \textbf{Repository} \textendash{} istoricul snapshot-urilor confirmate.
\end{itemize}

Comenzile \code{add}, \code{commit}, \code{push} muta starea intre zone in directia stanga $\to$ dreapta.

% =========================================================
\section{Cele cinci comenzi de baza}

Fluxul minim pentru publicarea unui director nou pe GitHub:

\begin{lstlisting}[style=bash]
git init -b main                               # 1
git add .                                      # 2
git commit -m "Initial commit"                 # 3
git remote add origin URL_DE_PE_GITHUB         # 4
git push -u origin main                        # 5
\end{lstlisting}

\begin{itemize}
  \item \code{git init -b main} \textendash{} initializeaza un repository local cu ramura implicita \code{main}.
  \item \code{git add .} \textendash{} include toate fisierele modificate in zona de pregatire (\emph{staging area}).
  \item \code{git commit -m "..."} \textendash{} creeaza un snapshot al starii pregatite, cu un mesaj descriptiv.
  \item \code{git remote add origin ...} \textendash{} asociaza repository-ul local cu un remote pe GitHub.
  \item \code{git push -u origin main} \textendash{} trimite commit-urile locale catre remote, stabilind ramura \code{main} ca \emph{upstream}.
\end{itemize}

% =========================================================
\section{Comenzi de inspectie (read-only)}

Inainte de \code{commit}, doua comenzi confirma ce urmeaza sa fie salvat:

\begin{lstlisting}[style=bash, caption={Verificare inainte de commit}]
git status                  # Lista fisierelor modificate / pregatite
git diff                    # Diferenta linie cu linie (working dir)
git diff --staged           # Diferenta in staging area
git log --oneline -10       # Ultimele 10 commit-uri, format compact
\end{lstlisting}

Aceste comenzi nu modifica nimic; pot fi rulate oricand fara risc.

% =========================================================
\section{Configurarea GitHub Pages}

Procedura standard pentru activarea hosting-ului:

\begin{enumerate}
  \item Creare a unui repository public la \url{https://github.com/new}, fara fisiere implicite (README, gitignore).
  \item Copierea URL-ului HTTPS afisat de GitHub si executarea comenzii \code{git remote add origin <URL>} urmata de \code{git push -u origin main}.
  \item Navigare la \emph{Settings} $\to$ \emph{Pages} (in panoul lateral al repository-ului).
  \item Selectarea ramurii \code{main} si a folder-ului \code{/ (root)} in sectiunea \emph{Branch}, urmata de \emph{Save}.
  \item Asteptarea de 30 \textendash{} 60 secunde, dupa care URL-ul public devine accesibil in caseta verde din partea de sus.
\end{enumerate}

\textbf{Anatomia URL-ului public:} \code{https://USERNAME.github.io/REPO}

% =========================================================
\section{Workflow zilnic dupa primul push}

Actualizarea continutului dupa deploy-ul initial necesita doar trei comenzi:

\begin{lstlisting}[style=bash]
git status                              # ce s-a modificat?
git add .                               # marcare pentru commit
git commit -m "Adauga sectiunea proiecte"
git push                                # trimitere pe GitHub
\end{lstlisting}

GitHub Pages reconstruieste si republica pagina automat, in aproximativ 60 de secunde.

% =========================================================
\section{Convenii pentru mesajele de commit}

Mesajul de commit ramane in istoricul repository-ului permanent. Calitatea acestuia este vizibila la \code{git log} si afecteaza utilitatea istoricului ca documentatie a evolutiei proiectului.

\begin{dano}[Asa nu]
\begin{lstlisting}[style=bash, numbers=none]
git commit -m "update"
git commit -m "fix"
git commit -m "asdf"
git commit -m "Modificari"
git commit -m "."
\end{lstlisting}
\textbf{Probleme.} Niciun context. La revizuire, autorul nu mai stie ce contine fiecare commit. Inutilizabil pentru \code{git revert}, \code{git bisect}, code review.
\end{dano}

\begin{dado}[Asa da]
\begin{lstlisting}[style=bash, numbers=none]
git commit -m "Adauga sectiunea hero cu CTA"
git commit -m "Fix overflow horizontal pe mobil < 360px"
git commit -m "Implementeaza dark mode prin light-dark()"
git commit -m "Reduce dimensiunea hero.jpg de la 2.4MB la 180KB"
\end{lstlisting}
\textbf{Reguli simple.} (1) Verb la imperativ in prima pozitie (\emph{Adauga}, \emph{Fix}, \emph{Implementeaza}, \emph{Reduce}). (2) Maxim 50\textendash{}72 caractere. (3) Descrie \emph{ce} si \emph{de ce}, nu \emph{cum} (codul arata cum).
\end{dado}

% =========================================================
\section{Fisierul .gitignore}

Anumite fisiere si directoare nu trebuie sa ajunga in repository: dependinte instalate (\code{node\_modules/}), fisiere temporare ale editoarelor, secretele (\code{.env}). Lista lor se declara in fisierul \code{.gitignore} la radacina proiectului:

\begin{lstlisting}[style=bash, caption={Exemplu minimal pentru proiect web static}]
# Dependinte
node_modules/

# Editoare
.vscode/
.idea/
*.swp

# Sistem
.DS_Store
Thumbs.db

# Secrete
.env
.env.local

# Build output
dist/
build/
\end{lstlisting}

Fisierul se commit-eaza ca orice alt fisier. Toate caile listate sunt ignorate la \code{git add}.

\begin{atentie}
\textbf{Nu commite niciodata fisiere \code{.env} cu chei API, parole, token-uri.} Odata commit-ate, raman in istoric chiar dupa stergere; pot fi recuperate de oricine cu acces la repository. Daca s-a intamplat: regenerati cheia compromisa imediat la furnizorul respectiv.
\end{atentie}

% =========================================================
\section{Probleme frecvente}

\begin{atentie}
\textbf{Autentificare prin parola.} GitHub a renuntat la autentificarea prin parola pentru operatiunile git (din 2021). Solutia recomandata este generarea unui \emph{Personal Access Token} (PAT) la \url{https://github.com/settings/tokens} cu scope-ul \code{repo}, folosit ulterior in locul parolei. Browserul memoreaza tokenul pentru sesiunile viitoare.
\end{atentie}

\begin{atentie}
\textbf{Eroare 404 pe pagina deployata.} Cauze posibile: (1) intervalul de propagare nu a fost respectat (60 secunde dupa \emph{Save}); (2) fisierul de pornire nu se numeste \code{index.html} (este sensibil la majuscule pe serverul GitHub); (3) cache-ul browserului afiseaza versiunea precedenta \textendash{} se rezolva cu \code{Ctrl+F5}.
\end{atentie}

% =========================================================
\section*{Verificare cuno\textcommabelow{s}tin\textcommabelow{t}e}

\begin{selfcheck}
\begin{enumerate}[itemsep=2pt]
  \item Care sunt cele 5 comenzi git pentru ini\textcommabelow{t}ializare \textcommabelow{s}i primul push?
  \item De ce GitHub a renun\textcommabelow{t}at la autentificarea prin parola? Care este alternativa?
  \item \emph{Spot the issue:} \code{git commit -m "update"}.
  \item Diferen\textcommabelow{t}a intre \code{git status} \textcommabelow{s}i \code{git diff}?
  \item Pentru un repository numit \code{portfolio} al utilizatorului \code{maria}, care este URL-ul GitHub Pages rezultat?
\end{enumerate}
\end{selfcheck}

% =========================================================
\section*{Recapitulare}

\begin{recap}[Faza 4 -- Deploy]
\begin{itemize}[itemsep=1pt]
  \item Modelul mental: working dir $\to$ staging $\to$ repository $\to$ remote.
  \item Cinci comenzi de baza: \code{git init -b main}, \code{git add .}, \code{git commit -m}, \code{git remote add origin}, \code{git push -u origin main}.
  \item Comenzi de inspectie: \code{git status}, \code{git diff}, \code{git log --oneline -10}.
  \item Activare Pages: \emph{Settings} $\to$ \emph{Pages} $\to$ branch \code{main}, folder \code{/ (root)} $\to$ \emph{Save}.
  \item Mesaje de commit: verb la imperativ, $\leq$72 caractere, descrie \emph{ce} si \emph{de ce}.
  \item Workflow zilnic: 3 comenzi (\code{add}, \code{commit}, \code{push}); redeploy automat in $\sim$60s.
  \item Autentificare prin Personal Access Token (PAT), nu prin parola.
  \item \code{.gitignore} exclude \code{node\_modules/}, \code{.env}, fisiere de editor, build artifacts.
\end{itemize}
\end{recap}

\partdivider

% =========================================================
\begin{joc}
\textbf{Platforma:} GitHub (\url{https://github.com}) \textendash{} aplicatia practica este deploy-ul efectiv al site-ului construit in fazele 1\textendash{}3.

\medskip

\textbf{Activitatea.} Procedura completa, executata sincron de toti participantii sub indrumarea trainer-ului:
\begin{enumerate}[itemsep=1pt]
  \item Inchiderea fisierelor in editor; navigare in \emph{Git Bash} la directorul proiectului.
  \item Executarea celor cinci comenzi git din sectiunea ,,Cele cinci comenzi de baza''.
  \item Crearea repository-ului public pe GitHub si conectarea remote-ului.
  \item Activarea GitHub Pages in setari.
  \item Verificarea URL-ului public si publicarea acestuia in canalul Discord al grupului.
\end{enumerate}

\textbf{Durata estimata:} 10 minute.
\end{joc}

% =========================================================
\section*{\faGraduationCap\ \ Resurse pentru aprofundare}

\begin{itemize}[itemsep=2pt]
  \item GitHub Inc. \emph{GitHub Pages documentation}. \url{https://pages.github.com}
  \item Cottle P. \emph{Learn Git Branching}. \url{https://learngitbranching.js.org}
  \item Chacon S., Straub B. \emph{Pro Git} (editia a 2-a). \url{https://git-scm.com/book/en/v2}
  \item GitHub Inc. \emph{Configuring a custom domain for your GitHub Pages site}. \\\url{https://docs.github.com/pages/configuring-a-custom-domain-for-your-github-pages-site}
  \item Beams C. \emph{How to Write a Git Commit Message}. \url{https://cbea.ms/git-commit/}
  \item Alternative: \url{https://www.netlify.com}, \url{https://vercel.com} (servicii similare cu deploy automat din git).
\end{itemize}

\bigskip

\bridge{Capitolul 5 (SEO) face site-ul descoperibil. URL public exista; acum trebuie ca motoarele de cautare \textcommabelow{s}i re\textcommabelow{t}elele sociale sa-l afi\textcommabelow{s}eze corect.}

% Fragment -- inclus de main.tex prin \input
\clearpage
\setpartlabel{themeSeo}{\faShareSquare}{Faza 5: SEO + Open Graph}
\setcounter{section}{0}
\phantomsection
\addcontentsline{toc}{part}{\protect\textbf{Faza 5 \textendash{} SEO + Open Graph}}
\handout{Faza 5 \textendash{} SEO + Open Graph}{Metadate pentru motoarele de cautare si retelele sociale.}

Capitolul prezinta tag-urile meta utilizate de motoarele de cautare (in principal Google) si de platformele sociale (Facebook, X, LinkedIn, WhatsApp, Discord) pentru afisarea corecta a paginilor distribuite. Acoperirea include SEO clasic, protocolul Open Graph (OGP), datele structurate Schema.org, si URL-urile canonice.

\bigskip

\begin{outcomes}
Dupa capitolul asta vei \textcommabelow{s}ti:
\begin{itemize}[itemsep=1pt]
  \item Configureaza \code{<title>} \textcommabelow{s}i \code{<meta description>} pentru rezultatele de cautare.
  \item Adauga cele 5 tag-uri \code{og:*} pentru previzualizare in re\textcommabelow{t}ele sociale.
  \item Genereaza \textcommabelow{s}i valideaza imaginea \code{og:image} la dimensiunile \textcommabelow{s}i formatul corect.
  \item Adauga date structurate (JSON-LD) pentru rich snippets in Google.
  \item Genereaza \code{robots.txt} \textcommabelow{s}i \code{sitemap.xml} pentru a ghida crawler-ele.
  \item Distinge mit-urile SEO de practica actuala.
  \item \faRocket\ \emph{Bonus self-study:} tipuri Schema.org (Article, BreadcrumbList, FAQ), \code{preconnect}/\code{dns-prefetch}, linking intern, mobile-first indexing.
\end{itemize}
\end{outcomes}

\begin{whymatters}
Un site fara SEO este invizibil. Aplica\textcommabelow{t}iile sociale moderne (WhatsApp, Discord, LinkedIn) afi\textcommabelow{s}eaza by default cardul de previzualizare; un site fara meta tags arata neprofesional cand cineva il share-uieste. SEO de baza este \emph{configurare}, nu \emph{magic}: cele 7\textendash{}10 tag-uri prezentate aici acopera 80\% din optimizarea esentiala.
\end{whymatters}

\begin{nota}[Pre-rechizite]
Capitolul presupune un site deployat (Faza 4) cu URL public accesibil. Tag-urile se valideaza pe varianta live, nu pe localhost.
\end{nota}

% =========================================================
\section{Doua audiente, doua seturi de tag-uri}

\begin{itemize}
  \item \textbf{Motoare de cautare} (Google, Bing). Folosesc \code{<title>}, \code{<meta name="description">}, structura semantica, datele Schema.org.
  \item \textbf{Platforme sociale} (Facebook, X, LinkedIn, WhatsApp, Discord). Folosesc protocolul Open Graph: \code{<meta property="og:*">}.
\end{itemize}

% =========================================================
\section{Metadate pentru motoarele de cautare}

Doua elemente determina aspectul paginii in rezultatele cautarii:

\begin{lstlisting}[style=html, caption={Plasare in <head>}]
<title>Voluntar BEST -- Student CS, Iasi</title>
<meta name="description" content="Pagina personala a Voluntarului BEST, studenta la CS, pasionata de web development.">
\end{lstlisting}

\subsection{Calitatea \texttt{<title>}}

\begin{dano}[Asa nu]
\begin{lstlisting}[style=html, numbers=none]
<title>Document</title>
<title>Untitled</title>
<title>Index</title>
<title>Acasa</title>
<title>web web web html portofoliu cv student iasi cs uaic</title>
\end{lstlisting}
\textbf{Probleme.} Generic, nesemnificativ; sau spam de cuvinte-cheie penalizat de Google.
\end{dano}

\begin{dado}[Asa da]
\begin{lstlisting}[style=html, numbers=none]
<title>Voluntar BEST -- Portofoliu Web Developer, Iasi</title>
<title>Despre BEST Iasi -- Asociatie studenteasca tehnica</title>
\end{lstlisting}
\textbf{Reguli.} 50\textendash{}60 caractere; specific paginii (nu site-ului); structura ,,Subiect \textendash{} Context''.
\end{dado}

\subsection{Calitatea \texttt{<meta description>}}

Aceleasi reguli ca pentru \code{<title>}, doar lungimea difera: \textbf{150\textendash{}160 caractere}, descrie con\textcommabelow{t}inutul concret (proiecte, CV, sec\textcommabelow{t}iuni), fara fraze de umplutura tip \emph{,,Bine a\textcommabelow{t}i venit''}.

\begin{lstlisting}[style=html, numbers=none]
<meta name="description" content="Portofoliul Voluntarului BEST: proiecte web (HTML, CSS, JS), CV, contact. Voluntar BEST Iasi si student CS la UAIC.">
\end{lstlisting}

% =========================================================
\section{Protocolul Open Graph}

Open Graph Protocol (OGP), specificat de Facebook in 2010, este standardul de fapt pentru previzualizarea paginilor web in retelele sociale. Implementarea consta in adaugarea a cinci tag-uri \code{<meta>} cu prefixul \code{og:}:

\begin{lstlisting}[style=html, caption={Tag-uri OGP minimale}]
<meta property="og:type" content="website">
<meta property="og:title" content="Voluntar BEST -- Student CS, Iasi">
<meta property="og:description" content="Pagina personala. Web dev, design, open source.">
<meta property="og:image" content="https://voluntar-best.github.io/site/og.png">
<meta property="og:url" content="https://voluntar-best.github.io/site/">
\end{lstlisting}

\begin{itemize}
  \item \code{og:type} \textendash{} categoria continutului. Valoarea \code{website} acopera majoritatea cazurilor; alternative: \code{article}, \code{video}, \code{book}, \code{profile}.
  \item \code{og:title} si \code{og:description} \textendash{} pot diferi de \code{<title>} si \code{<meta description>}, fiind optimizate specific pentru context social (mai punctuale, cu CTA).
  \item \code{og:image} \textendash{} dimensiunea recomandata 1200\,x\,630 pixeli (raport 1.91:1).
  \item \code{og:url} \textendash{} URL-ul canonic al paginii.
\end{itemize}

\subsection{Constrangerea \texttt{og:image}: URL absolut}

\begin{dano}[Asa nu \textendash{} URL relativ]
\code{<meta property="og:image" content="og.png">}\\
\code{<meta property="og:image" content="/images/og.png">}\\
\textbf{Problema.} Crawler-ele platformelor sociale nu pot rezolva caile relative, intrucat nu cunosc contextul de baza al paginii sursa. Cardul afisat ramane fara imagine.
\end{dano}

\begin{dado}[Asa da \textendash{} URL absolut]
\code{<meta property="og:image" content="https://voluntar-best.github.io/site/og.png">}\\
\textbf{Beneficiu.} Imaginea este accesibila direct prin URL si poate fi cache-uita de Facebook, Slack, Discord.
\end{dado}

\subsection{Calitatea imaginii \texttt{og:image}}

\begin{itemize}[itemsep=2pt]
  \item Dimensiune ideala: \textbf{1200\,x\,630 px} (raport 1.91:1, formatul afisat de Facebook si LinkedIn).
  \item Dimensiune minima: 600\,x\,315 px. Sub aceasta limita, unele platforme nu afiseaza cardul mare.
  \item Format: JPG sau PNG; dimensiunea maxima 8 MB (Facebook), in practica $\leq$1 MB pentru incarcare rapida.
  \item Continut lizibil la dimensiune redusa: titlul + 1\textendash{}2 elemente vizuale, nu paragrafe intregi.
\end{itemize}

% =========================================================
\section{Twitter Card}

Platforma X (anterior Twitter) suporta nativ OGP, dar accepta si tag-uri proprii pentru carduri cu imagine mare:

\begin{lstlisting}[style=html]
<meta name="twitter:card" content="summary_large_image">
\end{lstlisting}

Adaugarea acestei singure linii instruieste X sa afiseze cardul in formatul extins, folosind valorile deja prezente in tag-urile \code{og:*}.

% =========================================================
\section{Date structurate (JSON-LD)}

Schema.org defineste un vocabular standardizat pentru descrierea entitatilor (persoane, organizatii, evenimente, articole) intr-un format pe care motoarele de cautare il proceseaza explicit. Implementarea recomandata: JSON-LD plasat intr-un \code{<script>} cu tipul \code{application/ld+json}.

\begin{lstlisting}[style=html, caption={Schema Person pentru o pagina personala}]
<script type="application/ld+json">
{
  "@context": "https://schema.org",
  "@type": "Person",
  "name": "Voluntar BEST",
  "url": "https://voluntar-best.github.io/site/",
  "jobTitle": "Student CS",
  "alumniOf": "Universitatea Alexandru Ioan Cuza",
  "sameAs": [
    "https://github.com/voluntar-best",
    "https://www.linkedin.com/in/voluntar-best"
  ]
}
</script>
\end{lstlisting}

Avantajul principal: Google poate afisa pagina cu \emph{rich snippets} (foto, link-uri sociale, descriere structurata) in rezultatele cautarii. Validarea se face cu Google Rich Results Test (\url{https://search.google.com/test/rich-results}).

\begin{ponturi}
Pentru un eveniment, foloseste \code{"@type": "Event"} cu campurile \code{startDate}, \code{location}, \code{organizer}. Google poate afisa evenimentul direct intr-un panel lateral. Vocabularul complet: \url{https://schema.org}.
\end{ponturi}

% =========================================================
\section{URL-ul canonic}

Cand acelasi continut este accesibil prin mai multe URL-uri (cu si fara \code{www}, cu si fara slash final, cu parametri UTM, etc.), Google trebuie sa stie care URL este versiunea ,,oficiala''. Tag-ul canonical previne penalizarea pentru continut duplicat:

\begin{lstlisting}[style=html]
<link rel="canonical" href="https://voluntar-best.github.io/site/">
\end{lstlisting}

Indicatie: pe paginile de proiect, URL-ul canonic este URL-ul ,,curat'' al paginii, fara parametri de tracking.

% =========================================================
\section{Mituri SEO}

\begin{itemize}
  \item \textbf{Mit:} \code{<meta name="keywords">} influenteaza ranking-ul.
  \item \textbf{Realitate:} Google a anuntat oficial in 2009 ca ignora complet acest tag. Bing si Yandex il folosesc minim. \textbf{A nu se mai folosi.}
\end{itemize}

\begin{itemize}
  \item \textbf{Mit:} repetarea cuvintelor cheie in continut creste ranking-ul.
  \item \textbf{Realitate:} \emph{keyword stuffing} este penalizat. Algoritmii moderni (BERT, MUM) inteleg sensul textului si premiaza scrierea naturala.
\end{itemize}

\begin{itemize}
  \item \textbf{Mit:} site-urile noi sunt penalizate de Google.
  \item \textbf{Realitate:} indexarea dureaza 1\textendash{}7 zile pentru un site nou. Inscrierea in Google Search Console (\url{https://search.google.com/search-console}) accelereaza procesul.
\end{itemize}

% =========================================================
\section{\texttt{robots.txt} \textendash{} ghideaza crawler-ele}

\code{robots.txt} este un fisier text plasat in radacina site-ului (\code{https://voluntar-best.github.io/robots.txt}) care indica robotilor de crawl ce pagini sa indexeze si unde sa caute sitemap-ul. Format strict, dar simplu:

\begin{lstlisting}[caption={\texttt{robots.txt} minimal pentru GitHub Pages}, basicstyle=\ttfamily\small, backgroundcolor=\color{codeBg}, frame=leftline, framerule=2.5pt, rulecolor=\color{partcolor}, framexleftmargin=4pt, xleftmargin=14pt]
User-agent: *
Allow: /

Sitemap: https://voluntar-best.github.io/sitemap.xml
\end{lstlisting}

Cum se cite\textcommabelow{s}te:
\begin{itemize}[itemsep=1pt]
  \item \code{User-agent: *} \textendash{} regulile se aplica oricarui robot (Googlebot, Bingbot, etc.).
  \item \code{Allow: /} \textendash{} permite accesul la intreg site-ul.
  \item \code{Disallow: /admin/} (alternativ) \textendash{} interzice anumite cai. Util pentru pagini in lucru, dashboard-uri private.
  \item \code{Sitemap: ...} \textendash{} URL absolut catre sitemap-ul XML (vezi sectiunea urmatoare).
\end{itemize}

\begin{atentie}
\code{robots.txt} este o \emph{conventie}, nu un mecanism de securitate. Crawler-ele rau-voitoare il ignora. Daca ai date sensibile, foloseste autentificare reala, nu \code{Disallow}.
\end{atentie}

\begin{ponturi}
Verificare: deschide direct in browser \code{https://<site>/robots.txt}. Daca se afiseaza textul, e plasat corect.
\end{ponturi}

% =========================================================
\section{\texttt{sitemap.xml} \textendash{} ce pagini exista}

Sitemap-ul este un fisier XML care listeaza toate paginile importante ale site-ului. Google il foloseste pentru indexare mai rapida si completa, in special pentru site-urile noi sau cu structura adanca.

\begin{lstlisting}[caption={\texttt{sitemap.xml} scris manual pentru un portofoliu mic}, basicstyle=\ttfamily\small, backgroundcolor=\color{codeBg}, frame=leftline, framerule=2.5pt, rulecolor=\color{partcolor}, framexleftmargin=4pt, xleftmargin=14pt]
<?xml version="1.0" encoding="UTF-8"?>
<urlset xmlns="http://www.sitemaps.org/schemas/sitemap/0.9">
  <url>
    <loc>https://voluntar-best.github.io/</loc>
    <lastmod>2026-05-11</lastmod>
    <priority>1.0</priority>
  </url>
  <url>
    <loc>https://voluntar-best.github.io/projects.html</loc>
    <lastmod>2026-05-08</lastmod>
    <priority>0.8</priority>
  </url>
  <url>
    <loc>https://voluntar-best.github.io/contact.html</loc>
    <lastmod>2026-04-22</lastmod>
    <priority>0.5</priority>
  </url>
</urlset>
\end{lstlisting}

Reguli:
\begin{itemize}[itemsep=1pt]
  \item \code{<loc>} \textendash{} URL absolut (obligatoriu).
  \item \code{<lastmod>} \textendash{} data ultimei modificari, format ISO \code{YYYY-MM-DD} (op\textcommabelow{t}ional, recomandat).
  \item \code{<priority>} \textendash{} 0.0 \textendash{} 1.0, relativ in cadrul site-ului. Home = 1.0, pagini interne = 0.5\textendash{}0.8.
  \item Pentru site-uri mici (\textless{}50 pagini), scrierea manuala este OK. Pentru site-uri mari, foloseste generatoare automate.
\end{itemize}

\subsection{Submisie in Google Search Console}

\begin{enumerate}[itemsep=1pt]
  \item Deschide \url{https://search.google.com/search-console}, adauga proprietatea (URL prefix: \code{https://voluntar-best.github.io/}).
  \item Verifica posesiunea \textendash{} pentru GitHub Pages, metoda recomandata e meta tag in \code{<head>}.
  \item Sec\textcommabelow{t}iunea ,,Sitemaps'' $\to$ adauga \code{sitemap.xml} $\to$ Submit.
  \item In 1\textendash{}7 zile, Google va incepe sa indexeze toate URL-urile listate.
\end{enumerate}

\begin{ponturi}
\emph{Beneficiul real:} fara sitemap, Google poate descoperi pagini doar prin link-uri. Cu sitemap, le \textbf{declari explicit}. Pentru un site nou (fara backlink-uri inca), diferen\textcommabelow{t}a e intre indexare in 7 zile vs in 1\textendash{}2 zile.
\end{ponturi}

% =========================================================
\section{Favicon}

Favicon-ul (iconita afisata in tab-ul browserului si in bookmark-uri) se declara prin elementul \code{<link>}:

\begin{lstlisting}[style=html]
<link rel="icon" href="favicon.svg" type="image/svg+xml">
<link rel="apple-touch-icon" href="apple-touch-icon.png">
\end{lstlisting}

Generarea unui favicon simplu se poate face cu serviciul \url{https://favicon.io/favicon-generator/}, care produce fisierele necesare pornind de la o singura litera si o culoare. Pentru un favicon profesional din SVG, conversia automata in toate dimensiunile necesare se face cu \url{https://realfavicongenerator.net}.

% =========================================================
\begin{bonusbox}
\textbf{Acest box e optional.} Daca esti la inceput, treci direct la \emph{Verificare}. Mai jos sunt pattern-urile pe care le folosesc site-urile cu trafic real.

\medskip

\textbf{\faSitemap\ Tipuri Schema.org dincolo de Person}

Vocabularul Schema.org are sute de \code{@type}-uri. Cele mai utile in practica:

\medskip

\emph{Article \textendash{} pentru blog/posts:}
\begin{lstlisting}[style=html, numbers=none]
{
  "@context": "https://schema.org",
  "@type": "Article",
  "headline": "Cum am invatat web development in 2 saptamani",
  "author": { "@type": "Person", "name": "Voluntar BEST" },
  "datePublished": "2026-05-11",
  "dateModified": "2026-05-11",
  "image": "https://voluntar-best.github.io/blog/learning/cover.jpg"
}
\end{lstlisting}

\emph{BreadcrumbList \textendash{} navigare ierarhic\u a in rezultate:}
\begin{lstlisting}[style=html, numbers=none]
{
  "@context": "https://schema.org",
  "@type": "BreadcrumbList",
  "itemListElement": [
    { "@type": "ListItem", "position": 1, "name": "Acasa",     "item": "https://voluntar-best.github.io/" },
    { "@type": "ListItem", "position": 2, "name": "Proiecte",  "item": "https://voluntar-best.github.io/projects/" },
    { "@type": "ListItem", "position": 3, "name": "Portfolio Web" }
  ]
}
\end{lstlisting}
Google afi\textcommabelow{s}eaza aceasta ierarhie direct in rezultatele de cautare \textendash{} click rate semnificativ mai bun.

\emph{FAQPage \textendash{} intreb\u ari frecvente, afi\textcommabelow{s}ate ca rich snippet:}
\begin{lstlisting}[style=html, numbers=none]
{
  "@context": "https://schema.org",
  "@type": "FAQPage",
  "mainEntity": [{
    "@type": "Question",
    "name": "Cum aplic la BEST Iasi?",
    "acceptedAnswer": { "@type": "Answer", "text": "Inscrie-te la recrutarea din toamna pe best.org.ro." }
  }]
}
\end{lstlisting}

\emph{Event \textendash{} cu datele cheie indexate:}
\begin{lstlisting}[style=html, numbers=none]
{
  "@context": "https://schema.org",
  "@type": "Event",
  "name": "BEST Training Web 2026",
  "startDate": "2026-11-15T10:00:00+02:00",
  "location": { "@type": "Place", "name": "TUIASI" },
  "organizer": { "@type": "Organization", "name": "BEST Iasi" }
}
\end{lstlisting}

\medskip

\textbf{\faPlug\ \texttt{preconnect} \textcommabelow{s}i \texttt{dns-prefetch} \textendash{} performance + SEO signal}

Cand pagina ta foloseste resurse de pe alt domeniu (Google Fonts, CDN, Analytics), DNS lookup + TCP handshake adauga 100\textendash{}300ms inainte ca resursa sa inceapa sa se descarce. \code{preconnect} face acest setup in paralel cu HTML parsing, in fundal:

\begin{lstlisting}[style=html, numbers=none]
<link rel="preconnect" href="https://fonts.googleapis.com">
<link rel="preconnect" href="https://fonts.gstatic.com" crossorigin>
<link rel="dns-prefetch" href="https://www.google-analytics.com">
\end{lstlisting}

\code{preconnect} \textendash{} DNS + TCP + TLS (cost mai mare, dar resursa e gata sa fie descarcata).\\
\code{dns-prefetch} \textendash{} doar DNS (mai ieftin, foloseste pentru \textgreater 3 domenii).

Impact: imbunata\textcommabelow{t}e\textcommabelow{s}te LCP (Core Web Vitals) cu 100\textendash{}500ms cand pagina depinde de un CDN.

\medskip

\textbf{\faLink\ Linking intern \textendash{} arhitectur\u a, nu trick}

Google urmare\textcommabelow{s}te link-urile interne pentru a in\textcommabelow{t}elege relevan\textcommabelow{t}a relativa a paginilor. Pagini cu mai multe link-uri intrand sunt considerate ,,mai importante''.

\begin{itemize}[itemsep=1pt]
  \item Folose\textcommabelow{s}te \textbf{anchor text descriptiv} (\code{<a href="/projects">proiectele mele web</a>}), nu \code{click aici}.
  \item Conecteaza paginile inrudite: pe pagina de proiect, link spre alte proiecte. Pe blog post, link catre articole inrudite.
  \item Evita \textbf{orphan pages} (pagini fara niciun link intrand). Google le indexeaza tarziu sau deloc.
\end{itemize}

\medskip

\textbf{\faMobile*\ Mobile-first indexing}

Din 2023, Google indexeaza \textbf{varianta mobila} a paginii, nu desktop. Pagini care arata bine pe desktop dar sunt rupte pe mobile sunt penalizate masiv. Verifica:

\begin{itemize}[itemsep=1pt]
  \item \code{<meta name="viewport" content="width=device-width, initial-scale=1">} \textendash{} obligatoriu.
  \item Textul lizibil fara zoom (\code{font-size} \textgreater{} 14px pe mobil).
  \item Butoane clic\u abile (\textgreater 44x44px, conform Apple HIG).
  \item Con\textcommabelow{t}inutul important \textbf{nu este ascuns} pe mobil (Google considera continutul vizibil ca cel canonic).
  \item Aceea\textcommabelow{s}i pagina pe mobil \textcommabelow{s}i desktop \textendash{} nu varianta separata \code{m.site.ro}.
\end{itemize}

\medskip

\textbf{\faChartBar\ Performance ca semnal de ranking}

Din 2021, Core Web Vitals (LCP, INP, CLS \textendash{} vezi Faza 6) sunt \textbf{semnal de ranking}. Doua pagini cu con\textcommabelow{t}inut similar: cea mai rapida castiga. Nu vei concura cu site-uri lente.

Recap rapid:
\begin{itemize}[itemsep=1pt]
  \item \textbf{LCP} \textless 2.5s \textendash{} cea mai mare imagine/text apare repede.
  \item \textbf{INP} \textless 200ms \textendash{} pagina raspunde rapid la click/tap.
  \item \textbf{CLS} \textless 0.1 \textendash{} layout-ul nu ,,sare'' dupa incarcare.
\end{itemize}

\medskip

\textbf{\faRoute\ Unde mergi de aici}

\begin{itemize}[itemsep=1pt]
  \item \emph{Schema.org vocabular complet:} \url{https://schema.org/docs/full.html}
  \item \emph{Generator JSON-LD vizual:} \url{https://technicalseo.com/tools/schema-markup-generator/}
  \item \emph{Google Mobile-Friendly Test:} \url{https://search.google.com/test/mobile-friendly}
  \item \emph{Sitemap generator pentru site-uri statice:} \url{https://www.xml-sitemaps.com}
\end{itemize}
\end{bonusbox}

% =========================================================
\section*{Verificare cuno\textcommabelow{s}tin\textcommabelow{t}e}

\begin{selfcheck}
\begin{enumerate}[itemsep=2pt]
  \item Cate tag-uri \code{og:*} sunt minim necesare pentru un card valid de previzualizare?
  \item De ce \code{og:image} trebuie sa fie URL absolut, nu relativ?
  \item \emph{Spot the issue:} \code{<title>Document</title>}.
  \item Care este diferen\textcommabelow{t}a func\textcommabelow{t}ionala intre \code{<title>} si \code{og:title}?
  \item De ce \code{<meta name="keywords">} este considerat depasit?
  \item Ce reprezinta directiva \code{Sitemap:} in \code{robots.txt}?
  \item De ce \code{<lastmod>} in \code{sitemap.xml} este util pentru Google?
\end{enumerate}
\end{selfcheck}

% =========================================================
\section*{Recapitulare}

\begin{recap}[Faza 5 -- SEO + Open Graph]
\begin{itemize}[itemsep=1pt]
  \item Doua audiente: motoare de cautare (\code{<title>} + \code{<meta description>}) si retele sociale (5 \code{og:*} tags).
  \item \code{<title>}: 50\textendash{}60 caractere, specific, structura ,,Subiect \textendash{} Context''.
  \item \code{<meta description>}: 150\textendash{}160 caractere, descriere concreta, fara umplutura.
  \item \code{og:image}: URL absolut, dimensiune optima 1200\,x\,630, format JPG/PNG, $\leq$1MB.
  \item \code{<meta name="twitter:card" content="summary\_large\_image">} activeaza cardul mare pe X.
  \item JSON-LD (Schema.org) descrie entitati: Person, Organization, Article, Event.
  \item \code{<link rel="canonical">} previne penalizarea pentru continut duplicat.
  \item \code{robots.txt} ghideaza crawler-ele \textcommabelow{s}i declara sitemap-ul. \code{sitemap.xml} listeaza URL-urile importante.
  \item Submisia in Google Search Console accelereaza indexarea pentru site-uri noi.
  \item \code{<meta name="keywords">} este ignorat de Google din 2009 \textendash{} \textbf{a nu se folosi}.
  \item \faRocket\ Bonus: tipuri Schema.org (Article, BreadcrumbList, FAQ, Event), \code{preconnect}/\code{dns-prefetch}, linking intern, mobile-first indexing.
\end{itemize}
\end{recap}

\partdivider

% =========================================================
\begin{joc}
\textbf{Platforma:} OpenGraph.xyz (\url{https://opengraph.xyz}).

\medskip

\textbf{Activitatea.}
\begin{enumerate}[itemsep=1pt]
  \item Adaugarea celor cinci tag-uri \code{og:*} in \code{<head>}, cu valori specifice paginii personale.
  \item Inlocuirea valorii pentru \code{og:image} cu un URL absolut catre o imagine 1200\,x\,630, accesibila public.
  \item Commit si push: \code{git add .; git commit -m "Adauga meta tags Open Graph"; git push}.
  \item Accesarea OpenGraph.xyz si lipirea URL-ului public al paginii.
  \item Verificarea previzualizarilor pentru Facebook, X, LinkedIn, WhatsApp, iMessage.
\end{enumerate}

\textbf{Diagnoza problemelor:} (1) URL-ul \code{og:image} nu este absolut; (2) imaginea returneaza HTTP 404; (3) dimensiunea imaginii este sub 600 pixeli pe oricare axa.

\medskip

\textbf{Durata estimata:} 5 minute.
\end{joc}

% =========================================================
\section*{\faGraduationCap\ \ Resurse pentru aprofundare}

\begin{itemize}[itemsep=2pt]
  \item Open Graph Protocol. \emph{Specificatia oficiala}. \url{https://ogp.me}
  \item Meta Tags. \emph{Generator vizual}. \url{https://metatags.io}
  \item Google. \emph{Rich Results Test}. \url{https://search.google.com/test/rich-results}
  \item Google. \emph{Search Console} (indexare manuala). \url{https://search.google.com/search-console}
  \item Google. \emph{Document metadata}. \url{https://web.dev/learn/html/metadata}
  \item Schema.org. \emph{Structured data documentation}. \url{https://schema.org/docs/gs.html}
  \item X Inc. \emph{Twitter Cards documentation}. \url{https://developer.x.com/en/docs/twitter-for-websites/cards}
\end{itemize}

\bigskip

\bridge{Capitolul 6 (Lighthouse) ofera o masuratoare obiectiva a calita\textcommabelow{t}ii. Toate fix-urile aplicate in fazele 1\textendash{}5 sunt acum cuantificate intr-un scor 0\textendash{}100.}

% Fragment -- inclus de main.tex prin \input
\clearpage
\setpartlabel{themeLighthouse}{\faChartBar}{Faza 6: Lighthouse}
\setcounter{section}{0}
\phantomsection
\addcontentsline{toc}{part}{\protect\textbf{Faza 6 \textendash{} Lighthouse}}
\handout{Faza 6 \textendash{} Lighthouse}{Audit automatizat. Indicatori de performanta. Core Web Vitals.}

Lighthouse este un instrument de audit automatizat dezvoltat de Google, integrat in Chrome DevTools si in serviciul web PageSpeed Insights. Capitolul prezinta cele patru categorii evaluate, modul de interpretare a raportului, indicatorii \emph{Core Web Vitals} folositi de Google in algoritmul de ranking, si remedierile pentru problemele frecvente.

\bigskip

\begin{outcomes}
Dupa capitolul asta vei \textcommabelow{s}ti:
\begin{itemize}[itemsep=1pt]
  \item Ruleaza un audit Lighthouse pe propria pagina prin PageSpeed Insights sau DevTools.
  \item Cite\textcommabelow{s}te raportul \textcommabelow{s}i identifica problemele cu impact maxim pe scor.
  \item Aplica fix-uri uzuale: \code{alt}, \code{lang}, \code{width}/\code{height} pe imagini, \code{defer} pe script.
  \item Optimizeaza imaginile cu format modern (WebP), \code{srcset}, \code{loading="lazy"}.
  \item Interpreteaza Core Web Vitals (LCP, INP, CLS) \textcommabelow{s}i pragurile lor.
\end{itemize}
\end{outcomes}

\begin{whymatters}
Lighthouse este standardul industrial pentru evaluarea unei pagini web. Scorul 95+ pe toate cele 4 categorii este criteriul implicit la code review in firmele moderne. Mai mult, Core Web Vitals influenteaza direct ranking-ul Google. Un audit de 2 minute identifica 80\% din problemele care ar costa ore de cautare manuala.
\end{whymatters}

\begin{nota}[Pre-rechizite]
Capitolul presupune un site deployat (Faza 4) si configurat cu meta tags (Faza 5). Auditul ruleaza pe URL-ul public.
\end{nota}

% =========================================================
\section{Categoriile de evaluare}

Lighthouse genereaza un scor 0\textendash{}100 pe fiecare dintre cele patru axe:

\begin{itemize}
  \item \textbf{Performance} \textendash{} timpul de incarcare si receptivitatea paginii.
  \item \textbf{Accessibility} \textendash{} conformitatea cu standardele WCAG (in principal contrast, structura semantica, asocieri label\textendash{}input, atribute \code{alt}).
  \item \textbf{Best Practices} \textendash{} folosirea de protocoale si API-uri moderne, absenta erorilor in consola, lipsa codului depasit.
  \item \textbf{SEO} \textendash{} prezenta meta tag-urilor obligatorii, structura semantica, indexabilitatea.
\end{itemize}

Scorurile se interpreteaza astfel: peste 90 = bun; peste 95 = nivel profesional; sub 50 = necesita interventie.

\subsection{Lab data vs field data}

Lighthouse foloseste \emph{lab data}: un scenariu sintetic in conditii standardizate (CPU lent, retea simulata 4G slow). Scorul reflecta capacitatea pe un dispozitiv mediu, nu situatia reala a utilizatorilor. Pentru date din productie, PageSpeed Insights afiseaza si \emph{field data} extrase din Chrome User Experience Report (CrUX), agregate pe ultimele 28 de zile.

% =========================================================
\section{Modalitati de executie}

\subsection{PageSpeed Insights (interfata web)}

Cea mai accesibila metoda. Pasi: accesarea \url{https://pagespeed.web.dev}, lipirea URL-ului public al paginii, executarea analizei. Avantaje: nu necesita instalare; include datele reale de la utilizatori (CrUX) atunci cand acestea exista; raport partajabil prin URL.

\subsection{Chrome DevTools (local)}

Procedura: deschiderea paginii in Chrome, \code{F12} $\to$ tab \emph{Lighthouse}, selectarea categoriilor, executarea \emph{Analyze page load}. Avantaje: functioneaza pe paginile servite de pe \code{localhost}, fara necesitatea de deploy; permite testare rapida intre modificari.

\subsection{CLI (pentru integrare CI)}

Lighthouse este disponibil ca pachet npm. Folosit in pipeline-uri CI/CD pentru a bloca deploy-urile care scad scorul sub un prag stabilit.

% =========================================================
\section{Structura raportului}

Fiecare categorie include:

\begin{itemize}
  \item Scorul numeric, evidentiat color (rosu / portocaliu / verde).
  \item Lista problemelor identificate, grupate pe tip. Selectarea unei probleme afiseaza descrierea, impactul, si pasii de remediere.
  \item Lista verificarilor trecute cu succes.
  \item Sectiunea \emph{Opportunities} (pentru Performance) listeaza imbunatatirile cu economia estimata in milisecunde.
\end{itemize}

% =========================================================
\section{Probleme frecvente si remediere}

\subsection{Performance}

\begin{itemize}
  \item \textbf{Imagini fara dimensiuni explicite.} Setarea atributelor \code{width} si \code{height} pe \code{<img>} permite browserului sa rezerve spatiu, prevenind \emph{Cumulative Layout Shift}.
  \item \textbf{Resurse care blocheaza randarea.} Adaugarea atributului \code{defer} pe \code{<script>} permite parsarea HTML in paralel.
  \item \textbf{Imagini supradimensionate.} O imagine 2400\,x\,1600 redata la 800 pixeli iroseste banda. Optimizare: \url{https://squoosh.app}.
  \item \textbf{Format imagini neoptimizat.} JPEG/PNG in locul WebP / AVIF.
  \item \textbf{Cache lipsa.} Resursele statice (CSS, JS, imagini) ar trebui servite cu header \code{Cache-Control: max-age=31536000} (1 an). Pe GitHub Pages: cache-ul este configurat automat, fara interventie.
\end{itemize}

\subsection{Accessibility}

\begin{itemize}
  \item \textbf{Atribut \code{alt} lipsa pe \code{<img>}.} Toate imaginile relevante semantic necesita \code{alt}; imaginile pur decorative folosesc \code{alt=""}.
  \item \textbf{Buton fara nume accesibil.} Includerea de text in interior sau adaugarea \code{aria-label}.
  \item \textbf{Etichete lipsa pe formulare.} Asocierea \code{<label for>} cu \code{<input id>}.
  \item \textbf{Contrast insuficient.} Conform WCAG AA, raportul de contrast minim este 4.5:1 pentru text normal si 3:1 pentru text mare. Verificare: \url{https://webaim.org/resources/contrastchecker/}.
\end{itemize}

\subsection{Best Practices}

\begin{itemize}
  \item \textbf{Erori in consola.} Lighthouse penalizeaza orice eroare \code{console.error} la incarcare. Verificare: \code{F12 $\to$ Console}.
  \item \textbf{Resurse incarcate prin HTTP (nu HTTPS).} Mixed content; rezolvat automat pe GitHub Pages.
\end{itemize}

\subsection{SEO}

\begin{itemize}
  \item \textbf{Document fara \code{<title>}.}
  \item \textbf{Lipsa \code{<meta name="description">}.}
  \item \textbf{Element \code{<html>} fara \code{lang}.}
  \item \textbf{Linkuri fara text descriptiv} (\code{<a href="X">click aici</a>} este penalizat).
\end{itemize}

% =========================================================
\section{Core Web Vitals}

Cei trei indicatori folositi explicit de Google in algoritmul de ranking, incepand cu 2021:

\begin{itemize}
  \item \textbf{LCP} (\emph{Largest Contentful Paint}) \textendash{} timpul pana la afisarea celui mai mare element vizibil. Prag: $<$2.5 secunde.
  \item \textbf{INP} (\emph{Interaction to Next Paint}) \textendash{} latenta intre o actiune a utilizatorului si raspunsul vizual. Prag: $<$200 milisecunde.
  \item \textbf{CLS} (\emph{Cumulative Layout Shift}) \textendash{} masura instabilitatii vizuale in timpul incarcarii. Prag: $<$0.1.
\end{itemize}

% =========================================================
\section{Optimizarea imaginilor}

Imaginile reprezinta, in medie, 50\% din dimensiunea unei pagini web. Optimizarea lor are impact direct asupra scorului Performance.

\subsection{Format modern}

WebP reduce dimensiunea cu 25\textendash{}35\% fata de JPEG, AVIF cu 50\%. Conversia se face cu \url{https://squoosh.app} (interfata web, fara instalare).

\subsection{Atribute necesare pe \texttt{<img>}}

\begin{dano}[Asa nu]
\begin{lstlisting}[style=html, numbers=none]
<img src="hero.jpg">
\end{lstlisting}
\textbf{Probleme.} Lipseste \code{alt} (a11y); lipsesc \code{width}/\code{height} (CLS); lipseste \code{loading} (download eager); o singura sursa pentru toate ecranele (banda risipita pe mobil).
\end{dano}

\begin{dado}[Asa da \textendash{} optimizat complet]
\begin{lstlisting}[style=html, numbers=none]
<img src="hero-800.webp"
     srcset="hero-400.webp 400w,
             hero-800.webp 800w,
             hero-1600.webp 1600w"
     sizes="(min-width: 768px) 50vw, 100vw"
     width="800" height="450"
     loading="lazy"
     alt="Echipa BEST Iasi in fata sediului UAIC">
\end{lstlisting}
\textbf{Beneficii.} Browserul alege rezolutia potrivita din \code{srcset}; imaginile sub fold sunt incarcate doar la nevoie; rezerva spatiu prin \code{width}/\code{height}; alt asigura accesibilitatea.
\end{dado}

\subsection{Lazy loading}

Atributul \code{loading="lazy"} amana descarcarea imaginilor aflate sub \emph{fold} (zona vizibila initial). Suportul browser este universal incepand cu 2020.

\textbf{Exceptie:} pentru imaginea \emph{hero} (prima imagine din \emph{viewport}), se foloseste \code{loading="eager"} sau atributul lipseste; \code{loading="lazy"} aici intarzie LCP-ul.

% =========================================================
\section{Imbunatatiri suplimentare pentru Performance}

\subsection{Preload pentru resurse critice}

Resursele critice (font principal, imaginea hero) pot fi declarate cu \code{<link rel="preload">} pentru ca browserul sa le descarce inainte de a parsa CSS-ul.

\begin{lstlisting}[style=html]
<link rel="preload" href="fonts/lato.woff2" as="font" type="font/woff2" crossorigin>
<link rel="preload" href="hero.webp" as="image">
\end{lstlisting}

\subsection{CSS critic inline}

Pentru paginile cu prima vizita prioritara (landing pages), CSS-ul critic (necesar pentru randarea zonei vizibile) poate fi inclus direct in \code{<head>} prin \code{<style>}, restul fiind incarcat asincron. Tehnica nu este necesara pentru pagini personale; merita atentie pentru site-uri comerciale.

% =========================================================
\section*{Verificare cuno\textcommabelow{s}tin\textcommabelow{t}e}

\begin{selfcheck}
\begin{enumerate}[itemsep=2pt]
  \item Care sunt cele 4 categorii evaluate de Lighthouse?
  \item Diferen\textcommabelow{t}a intre \emph{lab data} \textcommabelow{s}i \emph{field data}?
  \item Care sunt cele 3 metrici Core Web Vitals \textcommabelow{s}i pragurile lor?
  \item \emph{Spot the issue:} \code{<img src="hero.jpg">} pe o pagina de productie.
  \item Cand este corect sa folose\textcommabelow{s}ti \code{loading="lazy"} \textcommabelow{s}i cand nu?
\end{enumerate}
\end{selfcheck}

% =========================================================
\section*{Recapitulare}

\begin{recap}[Faza 6 -- Lighthouse]
\begin{itemize}[itemsep=1pt]
  \item Patru categorii: Performance, Accessibility, Best Practices, SEO.
  \item Praguri: 90+ acceptabil, 95+ profesional, sub 50 necesita interventie.
  \item Lab data (Lighthouse) vs field data (CrUX); PageSpeed Insights le combina.
  \item Executie: \url{https://pagespeed.web.dev} (web) sau Chrome DevTools $\to$ \emph{Lighthouse}.
  \item Top fix-uri: \code{alt} pe imagini, \code{lang} pe \code{<html>}, contrast $\geq$4.5:1, \code{width}/\code{height} pe \code{<img>}, \code{defer} pe \code{<script>}.
  \item Core Web Vitals: LCP $<$2.5s, INP $<$200ms, CLS $<$0.1.
  \item Imagini: WebP/AVIF, \code{srcset}, \code{loading="lazy"} (cu exceptia imaginii hero), dimensiuni explicite.
\end{itemize}
\end{recap}

\partdivider

% =========================================================
\begin{joc}
\textbf{Platforma:} PageSpeed Insights (\url{https://pagespeed.web.dev}).

\medskip

\textbf{Activitatea.}
\begin{enumerate}[itemsep=1pt]
  \item \emph{Pasul 1 (3 minute).} Lipirea URL-ului paginii proprii deployate. Executarea analizei. Scorul tipic, dupa parcurgerea fazelor 1\textendash{}5, este peste 90.
  \item \emph{Pasul 2 (5 minute).} Identificarea uneia sau a doua probleme cu impact ridicat. Aplicarea fix-ului in cod local. \code{git push}. Reexecutarea analizei dupa 60 de secunde.
  \item \emph{Pasul 3 (2 minute).} Aplicarea aceluiasi audit pe varianta deployata a folder-ului \code{kit-starter/01-broken/}. Compararea scorurilor.
\end{enumerate}

\textbf{Continuare independenta:} folder-ul \code{01-broken} contine 10 probleme documentate in fisierul README. Tinta este atingerea scorului 95+ pe toate cele patru categorii.

\medskip

\textbf{Durata estimata:} 10 minute.
\end{joc}

% =========================================================
\section*{\faGraduationCap\ \ Resurse pentru aprofundare}

\begin{itemize}[itemsep=2pt]
  \item Google. \emph{PageSpeed Insights}. \url{https://pagespeed.web.dev}
  \item Google. \emph{Learn Performance}. \url{https://web.dev/learn/performance}
  \item Google. \emph{Web Vitals}. \url{https://web.dev/articles/vitals}
  \item Google. \emph{Lighthouse documentation}. \url{https://developer.chrome.com/docs/lighthouse/overview}
  \item WebAIM. \emph{Contrast Checker}. \url{https://webaim.org/resources/contrastchecker/}
  \item Google. \emph{Squoosh image optimizer}. \url{https://squoosh.app}
  \item Real Favicon Generator. \url{https://realfavicongenerator.net}
\end{itemize}

\bigskip

\bridge{Capitolul Resurse \& urmatorii pa\textcommabelow{s}i (cap. 7) propune un roadmap de 4 saptamani pentru consolidare \textcommabelow{s}i o lista curatata de tutoriale, comunita\textcommabelow{t}i \textcommabelow{s}i tool-uri pentru a continua singur.}

% Fragment -- inclus de main.tex prin \input
\clearpage
\setpartlabel{themeResurse}{\faGraduationCap}{Resurse \& urmatorii pasi}
\setcounter{section}{0}
\phantomsection
\addcontentsline{toc}{part}{\protect\textbf{Resurse \& urmatorii pasi}}
\handout{Resurse \& urmatorii pasi}{Aplicatii interactive. Tutoriale. Referinte. Comunitati.}

Acest capitol consolideaza resursele referentiate in capitolele anterioare si adauga sugestii pentru continuarea studiului dupa training. Materialele sunt grupate functional: aplicatii interactive (joculete), tutoriale structurate, referinte de cautare, comunitati, si tool-uri practice.

\bigskip

% =========================================================
\section{Roadmap pentru consolidare}

\begin{nota}[Saptamana 1]
Reconstruirea site-ului cu o tema personala (portofoliu, blog tematic, pagina pentru un eveniment). Folder-ul \code{kit-starter/02-reference/} serveste ca exemplu de implementare, nu ca sablon copiat direct.
\end{nota}

\begin{nota}[Saptamana 2]
Atingerea scorului 95+ pe toate cele patru categorii Lighthouse. Iteratia consta in audit \textendash{} identificare a problemelor \textendash{} remediere \textendash{} reaudit.
\end{nota}

\begin{nota}[Saptamana 3\textendash{}4]
Adaugarea unei functionalitati non-triviale: comutator pentru tema dark/light, formular cu trimitere efectiva (\url{https://formspree.io} ofera un endpoint gratuit), galerie de imagini cu lightbox, navigare cu mai multe pagini, animatii la scroll.
\end{nota}

\begin{nota}[Luna 2 si urmatoarele]
Inscrierea intr-un curriculum complet: \emph{The Odin Project} sau \emph{freeCodeCamp}. Ambele sunt gratuite si acopera traiectoria completa pana la dezvoltare full-stack, in 6\textendash{}12 luni de studiu sustinut.
\end{nota}

% =========================================================
\section{\faGamepad\ \ Aplicatii interactive (joculete)}

\begin{tcolorbox}[
  enhanced, colback=themeHtml!8, colframe=themeHtml, boxrule=0pt, leftrule=4pt,
  arc=4pt, sharp corners=downhill, rounded corners=northwest,
  left=12pt, right=12pt, top=8pt, bottom=8pt,
]
\textbf{\color{themeHtml}HTML / Accesibilitate}\\
\href{https://www.w3.org/WAI/demos/bad/}{W3C \emph{Before / After Demo}} \textendash{} comparatie intre o pagina inaccesibila si echivalentul sau accesibil.
\end{tcolorbox}

\begin{tcolorbox}[
  enhanced, colback=themeCss!8, colframe=themeCss, boxrule=0pt, leftrule=4pt,
  arc=4pt, sharp corners=downhill, rounded corners=northwest,
  left=12pt, right=12pt, top=8pt, bottom=8pt,
]
\textbf{\color{themeCss}CSS}\\
\textbf{Flexbox Froggy} \textendash{} \url{https://flexboxfroggy.com}\\
\textbf{Grid Garden} \textendash{} \url{https://cssgridgarden.com}\\
\textbf{CSS Diner} \textendash{} \url{https://flukeout.github.io} (selectori)\\
\textbf{CSS Battle} \textendash{} \url{https://cssbattle.dev} (reproducerea unei imagini cu CSS minimal)
\end{tcolorbox}

\begin{tcolorbox}[
  enhanced, colback=themeJs!8, colframe=themeJs, boxrule=0pt, leftrule=4pt,
  arc=4pt, sharp corners=downhill, rounded corners=northwest,
  left=12pt, right=12pt, top=8pt, bottom=8pt,
]
\textbf{\color{themeJs}JavaScript}\\
\textbf{JS Robot} \textendash{} \url{https://lifesphere.eu/jsrobot}\\
\textbf{Coding Fantasy} \textendash{} \url{https://codingfantasy.com} (lectiile JS sunt gratuite)\\
\textbf{Codewars} \textendash{} \url{https://codewars.com} (probleme de tip kata)
\end{tcolorbox}

\begin{tcolorbox}[
  enhanced, colback=themeDeploy!8, colframe=themeDeploy, boxrule=0pt, leftrule=4pt,
  arc=4pt, sharp corners=downhill, rounded corners=northwest,
  left=12pt, right=12pt, top=8pt, bottom=8pt,
]
\textbf{\color{themeDeploy}Git / Deploy}\\
\textbf{Learn Git Branching} \textendash{} \url{https://learngitbranching.js.org}\\
\textbf{Oh My Git!} \textendash{} \url{https://ohmygit.org} (joc desktop, gratuit, open source)
\end{tcolorbox}

% =========================================================
\section{\faBook\ \ Tutoriale structurate}

\begin{itemize}[itemsep=2pt]
  \item \textbf{The Odin Project.} \url{https://www.theodinproject.com} \\
  Curriculum complet, gratuit, parcurs estimat 6\textendash{}12 luni. Acoperire: HTML, CSS, JavaScript, Git, Node.js, Ruby on Rails sau React.
  \item \textbf{freeCodeCamp.} \url{https://www.freecodecamp.org} \\
  Certificari gratuite (Responsive Web Design, JavaScript Algorithms and Data Structures, Front End Development Libraries).
  \item \textbf{web.dev/learn.} \url{https://web.dev/learn} \\
  Cursuri scurte produse de Google, actualizate frecvent: HTML, CSS, JavaScript, Performance, Forms, Accessibility, PWA.
  \item \textbf{javascript.info.} \url{https://javascript.info} \\
  Manual extins pentru JavaScript, mai detaliat decat MDN, structurat didactic.
  \item \textbf{W3Schools.} \url{https://www.w3schools.com} \\
  Referin\textcommabelow{t}a sistematica cu \emph{,,Try it Yourself''} pe fiecare exemplu. Sec\textcommabelow{t}iuni separate: HTML, CSS, JS, TypeScript, React, SQL, Python.
  \item \textbf{Educative \textendash{} Web Development: Unraveling HTML, CSS, JS.} \url{https://www.educative.io} \\
  Curs comprehensiv (20h) cu pattern-ul ,,Hands-On / How it Works'' \textendash{} prezinta ac\textcommabelow{t}iunea, apoi explica.
  \item \textbf{Educative \textendash{} HTML5 Canvas: From Noob to Ninja.} \url{https://www.educative.io} \\
  Continuare pentru cei interesa\textcommabelow{t}i de animatii \textcommabelow{s}i grafica programatica (9h, peste 100 lec\textcommabelow{t}ii).
\end{itemize}

% =========================================================
\section{\faTrophy\ \ Pregatire pentru interviuri}

\begin{itemize}[itemsep=2pt]
  \item \textbf{GreatFrontend \textendash{} 50 must-know HTML/CSS/JS interview questions.} \\\url{https://www.greatfrontend.com/blog/50-must-know-html-css-and-javascript-interview-questions-by-ex-interviewers} \\
  Lista celor mai frecvente intrebari tehnice la interviurile front-end, formulata de fo\textcommabelow{s}ti intervievatori.
  \item \textbf{Frontend Interview Handbook.} \url{https://www.frontendinterviewhandbook.com} \\
  Resursa free pentru pregatirea interviurilor front-end (sistem design, coding, comportamentale).
  \item \textbf{LeetCode \textendash{} JavaScript track.} \url{https://leetcode.com} \\
  Probleme algoritmice si de structuri de date in JavaScript.
\end{itemize}

% =========================================================
\section{\faSearch\ \ Referinte de cautare}

\begin{itemize}[itemsep=2pt]
  \item \textbf{Mozilla Developer Network.} \url{https://developer.mozilla.org} \\
  Sursa de referinta pentru standardele web. Cautarea se efectueaza cu pattern-ul \emph{,,X mdn''} pe orice motor de cautare.
  \item \textbf{Can I Use.} \url{https://caniuse.com} \\
  Tabel detaliat pe browser si versiune pentru suportul fiecarui feature web.
  \item \textbf{CSS-Tricks.} \url{https://css-tricks.com} \\
  Articole de detaliu pentru concepte CSS specifice; ghidul de flexbox si cel de grid sunt referinte standard.
\end{itemize}

% =========================================================
\section{\faUsers\ \ Comunitati}

\begin{itemize}[itemsep=2pt]
  \item Subreddit-urile \emph{r/webdev}, \emph{r/learnprogramming}, \emph{r/Frontend} \textendash{} discutii zilnice si raspunsuri la intrebari concrete.
  \item \textbf{Stack Overflow.} \url{https://stackoverflow.com} \textendash{} platforma de intrebari/raspunsuri tehnice. Primul rezultat pentru majoritatea erorilor cautate pe Google.
  \item \textbf{Frontend Mentor.} \url{https://www.frontendmentor.io} \textendash{} provocari de tip design $\to$ implementare, cu peer review.
  \item Canalele Discord ale comunitatii BEST Iasi \textendash{} pentru intrebari specifice contextului local.
\end{itemize}

% =========================================================
\section{\faTools\ \ Tool-uri practice}

\begin{itemize}[itemsep=2pt]
  \item \textbf{Squoosh} \textendash{} \url{https://squoosh.app} (optimizare imagini).
  \item \textbf{TinyPNG} \textendash{} \url{https://tinypng.com} (compresie PNG / JPG).
  \item \textbf{Coolors} \textendash{} \url{https://coolors.co} (generare palete de culori).
  \item \textbf{Google Fonts} \textendash{} \url{https://fonts.google.com}.
  \item \textbf{Favicon Generator} \textendash{} \url{https://favicon.io}.
  \item \textbf{Excalidraw} \textendash{} \url{https://excalidraw.com} (diagrame, schite tehnice).
\end{itemize}

% =========================================================
\section{\faLightbulb\ \ Idei de proiecte (pentru cand te plictise\textcommabelow{s}ti)}

Cele mai bune skill-uri se invata facand ceva ce te intereseaza pe tine. Cateva idei, de la 30 minute la cateva weekenduri:

\subsection*{Personale (pentru CV / share pe social)}

\begin{itemize}[itemsep=2pt]
  \item Portofoliu personal cu poza, sectiune ,,Despre'', proiecte, contact.
  \item Blog despre orice te intereseaza (gaming, anime, gatit, drumetii, fotografie).
  \item ,,Card de vizita'' digital (CV web stylish, share-uibil prin URL).
  \item Pagina pentru un proiect BEST sau un eveniment universitar.
\end{itemize}

\subsection*{Distractive (1-2 ore fiecare)}

\begin{itemize}[itemsep=2pt]
  \item Tic-Tac-Toe sau Memory game in HTML/CSS/JS pur.
  \item Quiz interactiv cu intrebari preferate (\emph{,,Ce film e asta?''}, \emph{,,Ghice\textcommabelow{s}te tara dupa steag''}).
  \item Pomodoro timer (25 min lucru / 5 min pauza, cu sunet).
  \item Calculator simplu, dar cu un design mi\textcommabelow{s}to.
  \item Generator de palete de culori la apasare de buton.
  \item Joc snake clasic, in canvas.
\end{itemize}

\subsection*{Utile (le folose\textcommabelow{s}ti tu, prietenii tai)}

\begin{itemize}[itemsep=2pt]
  \item Habit tracker (verde la zilele cand ai facut treaba, ro\textcommabelow{s}u la cand n-ai).
  \item Lista de filme / serii / car\textcommabelow{t}i de vazut, cu rating.
  \item Tracker de cheltuieli (carbohidra\textcommabelow{t}i, bani, ore studiate).
  \item Pagina pentru proiectul tau de \textcommabelow{s}coala (mai stylish decat un PDF).
\end{itemize}

\subsection*{Fan / pasiune}

\begin{itemize}[itemsep=2pt]
  \item Pagina dedicata serialului / anime-ului / forma\textcommabelow{t}iei tale preferate.
  \item Galerie cu pozele tale de la concerte / drume\textcommabelow{t}ii / vacan\textcommabelow{t}e.
  \item Wiki personal pentru jocul tau favorit (Skyrim, Pokemon, etc.).
\end{itemize}

% =========================================================
\section{\faTrophy\ \ Provocari pentru cand vrei sa te depa\textcommabelow{s}e\textcommabelow{s}ti}

\begin{itemize}[itemsep=2pt]
  \item \textbf{,,Site in 1 ora''.} Cronometreaza-te. Pune un site decent online intr-o ora.
  \item \textbf{,,Lighthouse 100 / 100 / 100 / 100''.} Niciun fix nu este de prisos.
  \item \textbf{,,30 zile de cod''.} Un commit pe zi, orice marime, pe orice repo. Vezi cum se umple GitHub-ul de patrate verzi.
  \item \textbf{,,No CSS framework''.} Refa unul dintre proiectele tale fara Bootstrap/Tailwind \textendash{} doar CSS pur.
  \item \textbf{,,Cont GitHub cu 5 stele''.} Construie\textcommabelow{s}te ceva de care al\textcommabelow{t}ii vor sa puna stea. (Nu le cere stele \textendash{} fa-le de la sine.)
\end{itemize}

\partdivider

\begin{recap}[Observatie de incheiere]
Slide-urile prezentarii si pagina personala construita in training partajeaza acelasi set de tehnologii (HTML, CSS, JavaScript). Functionalitatea \emph{View Source} a oricarui browser ofera acces la codul sursa al oricarei pagini publice. Lectura acelui cod este, de acum inainte, accesibila.
\end{recap}


% Fragment -- inclus de main.tex prin \input
\clearpage
\setpartlabel{themePrereq}{\faBookOpen}{Anexe}
\setcounter{section}{0}
\phantomsection
\addcontentsline{toc}{part}{\protect\textbf{Anexe \textendash{} Glosar, Checklist, Erori comune}}
\handout{Anexe}{Glosar termeni. Checklist final. Erori comune si remediere.}

Sectiunea finala consolideaza referintele rapide: definitiile termenilor folositi pe parcursul manualului, lista de verificari inainte de publicarea unui site, si problemele intalnite frecvent impreuna cu solutiile lor.

% =========================================================
\section{Glosar de termeni}

\begin{description}[font=\sffamily\bfseries\color{partcolor}, leftmargin=0pt, labelindent=0pt, style=nextline, itemsep=4pt]
  \item[Accesibilitate (a11y)] Practica de a crea continut utilizabil de catre persoane cu dizabilitati. Standardul de referinta este WCAG (Web Content Accessibility Guidelines).
  \item[API] \emph{Application Programming Interface}. Set de functii prin care componentele software interactioneaza. In context web: \code{document.querySelector}, \code{fetch}, etc.
  \item[Arrow function] Sintaxa concisa pentru functii in JavaScript: \code{(a, b) => a + b}. Se comporta diferit fata de \code{function} privind \code{this}.
  \item[Cascade (CSS)] Algoritmul prin care browserul determina valoarea finala a unei proprietati cand mai multe reguli o vizeaza. Foloseste specificitate, ordine si importanta.
  \item[CDN] \emph{Content Delivery Network}. Reteaua de servere distribuite geografic pentru servirea resurselor cu latenta minima.
  \item[CLS] \emph{Cumulative Layout Shift}. Indicator Core Web Vitals pentru stabilitatea vizuala in timpul incarcarii.
  \item[Commit] Snapshot al starii fisierelor intr-un repository git, identificat printr-un hash unic.
  \item[CORS] \emph{Cross-Origin Resource Sharing}. Mecanism de securitate care controleaza accesul la resurse din alte domenii.
  \item[CSS Grid] Model de layout bidimensional, complementar lui Flexbox; util pentru aranjamente in retea.
  \item[DOM] \emph{Document Object Model}. Reprezentarea in memorie a documentului HTML, expusa programatic prin \code{document}.
  \item[Endpoint] URL specific care accepta cereri HTTP si returneaza un raspuns. Echivalent cu o ,,adresa'' a unui serviciu.
  \item[Flexbox] Model de layout uni-dimensional in CSS pentru distribuirea spatiului pe o axa.
  \item[Hosting] Servirea fisierelor unui site catre vizitatori. GitHub Pages, Netlify si Vercel sunt servicii de hosting pentru site-uri statice.
  \item[HTTP] \emph{HyperText Transfer Protocol}. Protocolul de baza al webului, prin care browserele cer si primesc resurse.
  \item[INP] \emph{Interaction to Next Paint}. Indicator Core Web Vitals pentru receptivitate; latenta intre interactiune si feedback vizual.
  \item[JSON-LD] \emph{JSON for Linked Data}. Format de date structurate folosit de motoarele de cautare pentru intelegerea continutului.
  \item[LCP] \emph{Largest Contentful Paint}. Indicator Core Web Vitals; momentul afisarii celui mai mare element vizibil.
  \item[Lighthouse] Tool de audit automatizat dezvoltat de Google; evalueaza Performance, Accessibility, Best Practices, SEO.
  \item[NodeList] Colectie de noduri DOM returnata de \code{querySelectorAll}; iterabila prin \code{forEach}.
  \item[Open Graph (OGP)] Standard initiat de Facebook pentru previzualizarea paginilor pe retele sociale (\code{og:title}, \code{og:image}, etc.).
  \item[PAT] \emph{Personal Access Token}. Token de autentificare pentru GitHub, folosit in locul parolei la operatiunile git.
  \item[Pseudo-clasa] Selector CSS care vizeaza o stare a unui element: \code{:hover}, \code{:focus}, \code{:checked}, \code{:first-child}.
  \item[Repository (repo)] Director versionat de git; contine codul, istoricul commit-urilor si configuratia.
  \item[Responsive design] Abordare prin care site-ul se adapteaza dimensiunii ecranului utilizatorului.
  \item[Schema.org] Vocabular standardizat pentru date structurate, suportat de motoarele de cautare majore.
  \item[Selector CSS] Pattern care identifica elementele HTML asupra carora se aplica o regula CSS.
  \item[Semantic HTML] Folosirea tag-urilor care descriu rolul continutului (\code{<header>}, \code{<nav>}) in locul \code{<div>} generic.
  \item[Specificitate] Mecanism prin care browserul rezolva conflictele intre reguli CSS care vizeaza acelasi element.
  \item[Static site] Site format din fisiere HTML/CSS/JS pre-generate, servite direct, fara procesare pe server.
  \item[Viewport] Zona vizibila a paginii in browser. Meta tag-ul \code{viewport} controleaza scalarea pe dispozitive mobile.
  \item[WCAG] \emph{Web Content Accessibility Guidelines}. Standardul international pentru accesibilitatea continutului web.
\end{description}

% =========================================================
\section{Checklist de publicare}

Lista de verificari aplicabila inainte de a considera un site finalizat. Se parcurge de sus in jos; orice item nebifat indica o lacuna remediabila.

\subsection*{HTML \textendash{} structura}

\begin{itemize}[leftmargin=*, label={$\square$}, itemsep=2pt]
  \item Documentul incepe cu \code{<!DOCTYPE html>}.
  \item \code{<html>} are atributul \code{lang}.
  \item \code{<head>} include \code{<meta charset>}, \code{<meta viewport>}, \code{<title>} specific paginii.
  \item Structura folosesc tag-uri semantice (\code{<header>}, \code{<main>}, \code{<section>}, \code{<footer>}, \code{<nav>}).
  \item Un singur \code{<h1>} per pagina; ierarhia titlurilor nu sare niveluri.
  \item Toate imaginile au \code{alt}; cele decorative au \code{alt=""}.
  \item Toate input-urile au \code{<label for>} corespunzator.
\end{itemize}

\subsection*{CSS \textendash{} aspect}

\begin{itemize}[leftmargin=*, label={$\square$}, itemsep=2pt]
  \item Reset minimal aplicat (\code{box-sizing: border-box}, \code{margin}/\code{padding} default 0).
  \item Variabilele CSS folosite pentru culori, spatiere, dimensiuni.
  \item Layout-ul rezista pe ecran 320px (test in DevTools \textendash{} mobile preview).
  \item Media queries pentru ecrane $\geq$768px.
  \item Contrast text/background $\geq$4.5:1 pentru text normal.
\end{itemize}

\subsection*{JavaScript \textendash{} comportament}

\begin{itemize}[leftmargin=*, label={$\square$}, itemsep=2pt]
  \item \code{<script>} are atributul \code{defer}.
  \item Niciun \code{var}; doar \code{const} si \code{let}.
  \item Comparatii doar cu \code{===}.
  \item Console-ul (F12) nu afiseaza erori la incarcarea paginii.
  \item Form-urile interceptate cu \code{e.preventDefault()} (sau au \code{action} valid).
\end{itemize}

\subsection*{SEO + Open Graph}

\begin{itemize}[leftmargin=*, label={$\square$}, itemsep=2pt]
  \item \code{<title>} specific, 50\textendash{}60 caractere.
  \item \code{<meta name="description">} 150\textendash{}160 caractere.
  \item Cinci tag-uri \code{og:*}: type, title, description, image, url.
  \item \code{og:image} este URL absolut, dimensiune $\geq$1200$\times$630.
  \item \code{<meta name="twitter:card" content="summary\_large\_image">} prezent.
  \item Favicon (\code{<link rel="icon">}) configurat.
\end{itemize}

\subsection*{Deploy + Lighthouse}

\begin{itemize}[leftmargin=*, label={$\square$}, itemsep=2pt]
  \item Site accesibil pe URL public (\code{username.github.io/repo}).
  \item \code{.gitignore} prezent in repository.
  \item Lighthouse Performance $\geq$90.
  \item Lighthouse Accessibility $\geq$95.
  \item Lighthouse Best Practices $\geq$95.
  \item Lighthouse SEO $\geq$95.
  \item Previzualizare validata pe opengraph.xyz.
\end{itemize}

% =========================================================
\section{Erori comune si remediere}

Lista neexhaustiva de probleme intalnite frecvent de incepatori, cu cauza si solutia.

\subsection*{HTML / CSS}

\begin{itemize}[itemsep=4pt]
  \item \textbf{Diacriticele apar ca \texttt{moz\u a}.} Cauza: lipsa \code{<meta charset="UTF-8">} sau fisier salvat in alta codificare. Solutie: adaugarea meta tag-ului si salvarea fisierului in UTF-8 din editor.
  \item \textbf{Stilurile CSS nu se aplica.} Cauze: (1) calea \code{href} incorecta in \code{<link>}; (2) numele clasei diferit fata de selector; (3) regula are specificitate mai mica decat alta. Verificare: DevTools $\to$ Inspect $\to$ tab \emph{Styles}.
  \item \textbf{Pe mobil pagina apare mica/dezordonat.} Cauza: lipsa \code{<meta name="viewport">}.
  \item \textbf{Flexbox-ul nu functioneaza.} Cauza: \code{display: flex} aplicat pe element gresit; trebuie pe \emph{containerul} parinte, nu pe copii.
\end{itemize}

\subsection*{JavaScript}

\begin{itemize}[itemsep=4pt]
  \item \textbf{\code{querySelector returns null}.} Cauze: (1) selectorul nu corespunde niciunui element; (2) script-ul ruleaza inainte de parsarea HTML-ului (lipseste \code{defer}); (3) typo in numele clasei/id-ului.
  \item \textbf{\code{addEventListener is not a function}.} Cauza: rezultatul \code{querySelector} este \code{null}. Verificare: \code{console.log(el)} inainte de \code{addEventListener}.
  \item \textbf{Form-ul reincarca pagina la submit.} Cauza: lipsa \code{e.preventDefault()} in handler.
  \item \textbf{\code{console.log not defined}.} Foarte rar; de obicei o tipografica (\code{Console.log} cu C mare).
\end{itemize}

\subsection*{Git / Deploy}

\begin{itemize}[itemsep=4pt]
  \item \textbf{\code{git push} cere parola si o respinge.} Cauza: GitHub a inlocuit autentificarea prin parola cu PAT. Solutie: generare PAT la \url{https://github.com/settings/tokens} cu scope \code{repo}; folosirea tokenului in locul parolei.
  \item \textbf{Pages returneaza 404 dupa push.} Cauze: (1) intervalul de propagare ($\sim$60s); (2) fisierul de pornire nu se numeste \code{index.html} (case-sensitive); (3) cache-ul browserului \textendash{} \code{Ctrl+F5} pentru hard refresh.
  \item \textbf{\code{fatal: not a git repository}.} Cauza: comanda rulata in afara unui director versionat. Solutie: \code{cd} in directorul corect sau \code{git init -b main}.
  \item \textbf{Modificarile nu apar pe Pages.} Cauza: nu s-a executat \code{git push} dupa \code{commit}; sau cache-ul Pages \textendash{} asteptare $\sim$60s.
\end{itemize}

\subsection*{SEO + Lighthouse}

\begin{itemize}[itemsep=4pt]
  \item \textbf{Card-ul WhatsApp nu afiseaza imaginea.} Cauza principala: \code{og:image} are URL relativ in loc de absolut. Verificare: opengraph.xyz raporteaza statusul.
  \item \textbf{Lighthouse Performance scor mic.} Cauze frecvente: (1) imagini nedimensionate, supradimensionate; (2) script-uri sincrone in \code{<head>}; (3) format imagini neoptimizat (PNG/JPEG cu fisiere mari).
  \item \textbf{Lighthouse Accessibility 50\textendash{}80.} Cauze frecvente: (1) \code{alt} lipsa; (2) contrast scazut; (3) lipsa \code{lang} pe \code{<html>}; (4) butoane fara nume accesibil.
\end{itemize}


\end{document}
